\documentclass{article}
\usepackage[utf8]{inputenc}
\usepackage[T1]{fontenc}
\usepackage{textcomp}
\usepackage{amssymb}
\usepackage[margin=1in, headheight=26pt]{geometry}
\usepackage{graphicx}
\usepackage[most]{tcolorbox}
\usepackage{mdframed}
\usepackage{xcolor}
\usepackage{fancyhdr}

\pagestyle{fancy}
\fancyhf{}
\newcommand{\currenttopic}{Automata Theory}
\fancyhead[L]{\textbf{\currenttopic}}
\fancyfoot[C]{\thepage}
\renewcommand{\headrulewidth}{0.4pt}

\graphicspath{ {./images/} }

\newcommand{\newtopic}[1]{
    \cleardoublepage
    \renewcommand{\currenttopic}{#1}
    \sectionmark{#1}
    \addcontentsline{toc}{section}{#1}
}

\begin{document}

\title{Automata Theory MCQ Collection}
\author{Sanfoundry Scraping Project}
\date{\today}
\maketitle
\tableofcontents
\newpage

\newtopic{Finite Automata – Basics}

\begin{mdframed}[linewidth=1pt, roundcorner=5pt, linecolor=gray!30, backgroundcolor=gray!5]
\textbf{Question 1} 

Assume the R is a relation on a set A, aRb is partially ordered such that a and b are \_\_\_\_\_\_\_\_\_\_\_\_\_

a) reflexive

b) transitive

c) symmetric

d) reflexive and transitive
\vspace{5pt} \par
\begin{tcolorbox}[colback=green!5, colframe=green!40!black, title=Solution]
Answer: d

Explanation: A partially ordered relation refers to one which is Reflexive, Transitive and Antisymmetric.
\end{tcolorbox}
\end{mdframed}
\newpage

\begin{mdframed}[linewidth=1pt, roundcorner=5pt, linecolor=gray!30, backgroundcolor=gray!5]
\textbf{Question 2} 

The non- Kleene Star operation accepts the following string of finite length over set A = \{0,1\} | where string s contains even number of 0 and 1

a) 01,0011,010101

b) 0011,11001100

c) $\epsilon$,0011,11001100

d) $\epsilon$,0011,11001100
\vspace{5pt} \par
\begin{tcolorbox}[colback=green!5, colframe=green!40!black, title=Solution]
Answer: b

Explanation: The Kleene star of A, denoted by A*, is the set of all strings obtained by concatenating zero or more strings from A.
\end{tcolorbox}
\end{mdframed}
\newpage

\begin{mdframed}[linewidth=1pt, roundcorner=5pt, linecolor=gray!30, backgroundcolor=gray!5]
\textbf{Question 3} 

A regular language over an alphabet $\Sigma$ is one that cannot be obtained from the basic languages using the operation

a) Union

b) Concatenation

c) Kleene*

d) All of the mentioned
\vspace{5pt} \par
\begin{tcolorbox}[colback=green!5, colframe=green!40!black, title=Solution]
Answer: d

Explanation: Union, Intersection, Concatenation, Kleene*, Reverse are all the closure properties of Regular Language.
\end{tcolorbox}
\end{mdframed}
\newpage

\begin{mdframed}[linewidth=1pt, roundcorner=5pt, linecolor=gray!30, backgroundcolor=gray!5]
\textbf{Question 4} 

Statement 1: A Finite automata can be represented graphically; Statement 2: The nodes can be its states;   Statement 3: The edges or arcs can be used for transitions

Hint: Nodes and Edges are for trees and forests too.

Which of the following make the correct combination?

a) Statement 1 is false but Statement 2 and 3 are correct

b) Statement 1 and 2 are correct while 3 is wrong

c) None of the mentioned statements are correct

d) All of the mentioned
\vspace{5pt} \par
\begin{tcolorbox}[colback=green!5, colframe=green!40!black, title=Solution]
Answer: d

Explanation: It is possible to represent a finite automaton graphically, with nodes for states, and arcs for transitions.
\end{tcolorbox}
\end{mdframed}
\newpage

\begin{mdframed}[linewidth=1pt, roundcorner=5pt, linecolor=gray!30, backgroundcolor=gray!5]
\textbf{Question 5} 

The minimum number of states required to recognize an octal number divisible by 3 are/is

a) 1

b) 3

c) 5

d) 7
\vspace{5pt} \par
\begin{tcolorbox}[colback=green!5, colframe=green!40!black, title=Solution]
Answer: b

Explanation: According to the question, minimum of 3 states are required to recognize an octal number divisible by 3.
\end{tcolorbox}
\end{mdframed}
\newpage

\begin{mdframed}[linewidth=1pt, roundcorner=5pt, linecolor=gray!30, backgroundcolor=gray!5]
\textbf{Question 6} 

Which of the following is not a part of 5-tuple finite automata?

a) Input alphabet

b) Transition function

c) Initial State

d) Output Alphabet
\vspace{5pt} \par
\begin{tcolorbox}[colback=green!5, colframe=green!40!black, title=Solution]
Answer: d

Explanation: A FA can be represented as FA = (Q, $\Sigma$, $\delta$, q0, F) where Q=Finite Set of States, $\Sigma$=Finite Input Alphabet, $\delta$=Transition Function, q0=Initial State, F=Final/Acceptance State).
\end{tcolorbox}
\end{mdframed}
\newpage

\begin{mdframed}[linewidth=1pt, roundcorner=5pt, linecolor=gray!30, backgroundcolor=gray!5]
\textbf{Question 7} 

If an Infinite language is passed to Machine M, the subsidiary which gives a finite solution to the infinite input tape is \_\_\_\_\_\_\_\_\_\_\_\_\_\_

a) Compiler

b) Interpreter

c) Loader and Linkers

d) None of the mentioned
\vspace{5pt} \par
\begin{tcolorbox}[colback=green!5, colframe=green!40!black, title=Solution]
Answer: a

Explanation: A Compiler is used to give a finite solution to an infinite phenomenon. An example of an infinite phenomenon is Language C, etc.
\end{tcolorbox}
\end{mdframed}
\newpage

\begin{mdframed}[linewidth=1pt, roundcorner=5pt, linecolor=gray!30, backgroundcolor=gray!5]
\textbf{Question 8} 

The number of elements in the set for the Language L=\{x$\epsilon$($\Sigma$r) *|length if x is at most 2\} and $\Sigma$=\{0,1\} is \_\_\_\_\_\_\_\_\_

a) 7

b) 6

c) 8

d) 5
\vspace{5pt} \par
\begin{tcolorbox}[colback=green!5, colframe=green!40!black, title=Solution]
Answer: a

Explanation: $\Sigma$r = \{1,0\} and a Kleene* operation would lead to the following set=COUNT\{$\epsilon$,0,1,00,11,01,10\} = 7.
\end{tcolorbox}
\end{mdframed}
\newpage

\begin{mdframed}[linewidth=1pt, roundcorner=5pt, linecolor=gray!30, backgroundcolor=gray!5]
\textbf{Question 9} 

For the following change of state in FA, which of the following codes is an incorrect option?

a) $\delta$ (m, 1) = n

b) $\delta$ (0, n) = m

c) $\delta$ (m,0) = $\epsilon$

d) 


s: accept = false; cin $>$$>$ char;

if char = "0" goto n;
\vspace{5pt} \par
\begin{tcolorbox}[colback=green!5, colframe=green!40!black, title=Solution]
Answer: b

Explanation: $\delta$(QX$\Sigma$) = Q1 is the correct representation of change of state. Here, $\delta$ is called the Transition function.
\end{tcolorbox}
\end{mdframed}
\newpage

\begin{mdframed}[linewidth=1pt, roundcorner=5pt, linecolor=gray!30, backgroundcolor=gray!5]
\textbf{Question 10} 

Given: $\Sigma$= \{a, b\}

L= \{x$\epsilon$$\Sigma$*|x is a string combination\}

$\Sigma$4 represents which among the following?

a) \{aa, ab, ba, bb\}

b) \{aaaa, abab, $\epsilon$, abaa, aabb\}

c) \{aaa, aab, aba, bbb\}

d) All of the mentioned
\vspace{5pt} \par
\begin{tcolorbox}[colback=green!5, colframe=green!40!black, title=Solution]
Answer: b

Explanation: $\Sigma$* represents any combination of the given set while $\Sigma$x represents the set of combinations with length x where x $\epsilon$ I.
\end{tcolorbox}
\end{mdframed}
\newpage

\newtopic{Deterministic Finite Automata-Introduction and Definition}

\begin{mdframed}[linewidth=1pt, roundcorner=5pt, linecolor=gray!30, backgroundcolor=gray!5]
\textbf{Question 1} 

Which of the following not an example Bounded Information?

a) fan switch outputs \{on, off\}

b) electricity meter reading

c) colour of the traffic light at the moment

d) none of the mentioned
\vspace{5pt} \par
\begin{tcolorbox}[colback=green!5, colframe=green!40!black, title=Solution]
Answer: b

Explanation: Bounded information refers to one whose output is limited and it cannot be said what were the recorded outputs previously until memorized.
\end{tcolorbox}
\end{mdframed}
\newpage

\begin{mdframed}[linewidth=1pt, roundcorner=5pt, linecolor=gray!30, backgroundcolor=gray!5]
\textbf{Question 2} 

A Language for which no DFA exist is a\_\_\_\_\_\_\_\_

a) Regular Language

b) Non-Regular Language

c) May be Regular

d) Cannot be said
\vspace{5pt} \par
\begin{tcolorbox}[colback=green!5, colframe=green!40!black, title=Solution]
Answer: b

Explanation: A language for which there is no existence of a deterministic finite automata is always Non Regular and methods like Pumping Lemma can be used to prove the same.
\end{tcolorbox}
\end{mdframed}
\newpage

\begin{mdframed}[linewidth=1pt, roundcorner=5pt, linecolor=gray!30, backgroundcolor=gray!5]
\textbf{Question 3} 

A DFA cannot be represented in the following format

a) Transition graph

b) Transition Table

c) C code

d) None of the mentioned
\vspace{5pt} \par
\begin{tcolorbox}[colback=green!5, colframe=green!40!black, title=Solution]
Answer: d

Explanation: A DFA can be represented in the following formats: Transition Graph, Transition Table, Transition tree/forest/Any programming Language.
\end{tcolorbox}
\end{mdframed}
\newpage

\begin{mdframed}[linewidth=1pt, roundcorner=5pt, linecolor=gray!30, backgroundcolor=gray!5]
\textbf{Question 4} 

What the following DFA accepts?


 
\begin{center}\includegraphics[width=0.5\textwidth]{automata-theory-interview-questions-answers-q4.png}\end{center}
 

a) x is a string such that it ends with ‘101'

b) x is a string such that it ends with ‘01'

c) x is a string such that it has odd 1's and even 0's

d) x is a strings such that it has starting and ending character as 1
\vspace{5pt} \par
\begin{tcolorbox}[colback=green!5, colframe=green!40!black, title=Solution]
Answer: a

Explanation: Strings such as \{1101,101,10101\} are being accepted while \{1001,11001\} are not. Thus, this conclusion leads to option a.
\end{tcolorbox}
\end{mdframed}
\newpage

\begin{mdframed}[linewidth=1pt, roundcorner=5pt, linecolor=gray!30, backgroundcolor=gray!5]
\textbf{Question 5} 

When are 2 finite states equivalent?

a) Same number of transitions

b) Same number of states

c) Same number of states as well as transitions

d) Both are final states
\vspace{5pt} \par
\begin{tcolorbox}[colback=green!5, colframe=green!40!black, title=Solution]
Answer: c

Explanation: Two states are said to be equivalent if and only if they have same number of states as well as transitions.





     Apply Now for 
Free C Certification
 - April 2026
\end{tcolorbox}
\end{mdframed}
\newpage

\begin{mdframed}[linewidth=1pt, roundcorner=5pt, linecolor=gray!30, backgroundcolor=gray!5]
\textbf{Question 6} 

What does the following figure most correctly represents?


 
\begin{center}\includegraphics[width=0.5\textwidth]{automata-theory-interview-questions-answers-q6.png}\end{center}
 

a) Final state with loop x

b) Transitional state with loop x

c) Initial state as well as final state with loop x

d) Insufficient Data
\vspace{5pt} \par
\begin{tcolorbox}[colback=green!5, colframe=green!40!black, title=Solution]
Answer: c

Explanation: The figure represents the initial as well as the final state with an iteration of x.
\end{tcolorbox}
\end{mdframed}
\newpage

\begin{mdframed}[linewidth=1pt, roundcorner=5pt, linecolor=gray!30, backgroundcolor=gray!5]
\textbf{Question 7} 

Which of the following will not be accepted by the following DFA?


 
\begin{center}\includegraphics[width=0.5\textwidth]{automata-theory-interview-questions-answers-q7.png}\end{center}
 

a) ababaabaa

b) abbbaa

c) abbbaabb

d) abbaabbaa
\vspace{5pt} \par
\begin{tcolorbox}[colback=green!5, colframe=green!40!black, title=Solution]
Answer: a

Explanation: All the Strings are getting accepted except ‘ababaabaa' as it is directed to dumping state. Dumping state also refers to the reject state of the automata.
\end{tcolorbox}
\end{mdframed}
\newpage

\begin{mdframed}[linewidth=1pt, roundcorner=5pt, linecolor=gray!30, backgroundcolor=gray!5]
\textbf{Question 8} 

Which of the following will the given DFA won't accept?


 
\begin{center}\includegraphics[width=0.5\textwidth]{automata-theory-interview-questions-answers-q8.png}\end{center}
 

a) $\epsilon$

b) 11010

c) 10001010

d) String of letter count 11
\vspace{5pt} \par
\begin{tcolorbox}[colback=green!5, colframe=green!40!black, title=Solution]
Answer: a

Explanation: As the initial state is not made an acceptance state, thus $\epsilon$ will not be accepted by the given DFA. For the automata to accept $\epsilon$ as an entity, one should make the initial state as also the final state.
\end{tcolorbox}
\end{mdframed}
\newpage

\begin{mdframed}[linewidth=1pt, roundcorner=5pt, linecolor=gray!30, backgroundcolor=gray!5]
\textbf{Question 9} 

Can a DFA recognize a palindrome number?

a) Yes

b) No

c) Yes, with input alphabet as $\Sigma$*

d) Can't be determined
\vspace{5pt} \par
\begin{tcolorbox}[colback=green!5, colframe=green!40!black, title=Solution]
Answer: b

Explanation: Language to accept a palindrome number or string will be non-regular and thus, its DFA cannot be obtained. Though, PDA is possible.
\end{tcolorbox}
\end{mdframed}
\newpage

\begin{mdframed}[linewidth=1pt, roundcorner=5pt, linecolor=gray!30, backgroundcolor=gray!5]
\textbf{Question 10} 

Which of the following is not an example of finite state machine system?

a) Control Mechanism of an elevator

b) Combinational Locks

c) Traffic Lights

d) Digital Watches
\vspace{5pt} \par
\begin{tcolorbox}[colback=green!5, colframe=green!40!black, title=Solution]
Answer: d

Explanation: Proper and sequential combination of events leads the machines to work in hand which includes The elevator, Combinational Locks, Traffic Lights, vending machine, etc. Other applications of Finite machine state system are Communication Protocol Design, Artificial Intelligence Research, A Turnstile, etc.
\end{tcolorbox}
\end{mdframed}
\newpage

\newtopic{DFA Processing Strings}

\begin{mdframed}[linewidth=1pt, roundcorner=5pt, linecolor=gray!30, backgroundcolor=gray!5]
\textbf{Question 1} 

The password to the admins account="administrator". The total number of states required to make a password-pass system using DFA would be \_\_\_\_\_\_\_\_\_\_

a) 14 states

b) 13 states

c) 12 states

d) A password pass system cannot be created using DFA
\vspace{5pt} \par
\begin{tcolorbox}[colback=green!5, colframe=green!40!black, title=Solution]
Answer: a

Explanation: For a string of n characters with no repetitive substrings, the number of states required to pass the string is n+1.
\end{tcolorbox}
\end{mdframed}
\newpage

\begin{mdframed}[linewidth=1pt, roundcorner=5pt, linecolor=gray!30, backgroundcolor=gray!5]
\textbf{Question 2} 

Which of the following is the corresponding Language to the given DFA?


 
\begin{center}\includegraphics[width=0.5\textwidth]{automata-theory-questions-answers-dfa-processing-strings-q2.png}\end{center}
 

a) L = \{x $\epsilon$ \{0, 1\} * | x ends in 1 and does not contain substring 01\}

b) L = \{x $\epsilon$ \{0,1\} * |x ends in 1 and does not contain substring 00\}

c) L = \{x $\epsilon$ \{0,1\} |x ends in 1 and does not contain substring 00\}

d) L = \{x $\epsilon$ \{0,1\} * |x ends in 1 and does not contain substring 11\}
\vspace{5pt} \par
\begin{tcolorbox}[colback=green!5, colframe=green!40!black, title=Solution]
Answer: b

Explanation: The Language can be anonymously checked and thus the answer can be predicted. The language needs to be accepted by the automata (acceptance state) in order to prove its regularity.
\end{tcolorbox}
\end{mdframed}
\newpage

\begin{mdframed}[linewidth=1pt, roundcorner=5pt, linecolor=gray!30, backgroundcolor=gray!5]
\textbf{Question 3} 

Let $\Sigma$= \{a, b, …. z\} and A = \{Hello, World\}, B= \{Input, Output\}, then (A*$\cap$B) U (B*$\cap$A) can be represented as:

a) \{Hello, World, Input, Output, $\epsilon$\}

b) \{Hello, World, $\epsilon$\}

c) \{Input, Output, $\epsilon$\}

d) \{\}
\vspace{5pt} \par
\begin{tcolorbox}[colback=green!5, colframe=green!40!black, title=Solution]
Answer: d

Explanation: Union operation creates the universal set by combining all the elements of first and second set while intersection operation creates a set of common elements of the first and the second state.
\end{tcolorbox}
\end{mdframed}
\newpage

\begin{mdframed}[linewidth=1pt, roundcorner=5pt, linecolor=gray!30, backgroundcolor=gray!5]
\textbf{Question 4} 

Let the given DFA consist of x states. Find x-y such that y is the number of states on minimization of DFA?


 
\begin{center}\includegraphics[width=0.5\textwidth]{automata-theory-questions-answers-dfa-processing-strings-q4.png}\end{center}
 

a) 3

b) 2

c) 1

d) 4
\vspace{5pt} \par
\begin{tcolorbox}[colback=green!5, colframe=green!40!black, title=Solution]
Answer: b

Explanation: Use the equivalence theorem or Myphill Nerode theorem to minimize the DFA.
\end{tcolorbox}
\end{mdframed}
\newpage

\begin{mdframed}[linewidth=1pt, roundcorner=5pt, linecolor=gray!30, backgroundcolor=gray!5]
\textbf{Question 5} 

For a machine to surpass all the letters of alphabet excluding vowels, how many number of states in DFA would be required?

a) 3

b) 2

c) 22

d) 27
\vspace{5pt} \par
\begin{tcolorbox}[colback=green!5, colframe=green!40!black, title=Solution]
Answer: a

Explanation:


 
\begin{center}\includegraphics[width=0.5\textwidth]{automata-theory-questions-answers-dfa-processing-strings-q5.png}\end{center}
 







     
Free Certifications
 on 300 subjects are live for April 2026. Register Now!
\end{tcolorbox}
\end{mdframed}
\newpage

\begin{mdframed}[linewidth=1pt, roundcorner=5pt, linecolor=gray!30, backgroundcolor=gray!5]
\textbf{Question 6} 

For the DFA given below compute the following:

Union of all possible combinations at state 7,8 and 9.


 
\begin{center}\includegraphics[width=0.5\textwidth]{automata-theory-questions-answers-dfa-processing-strings-q6.png}\end{center}
 

a) \{aba, ac, cc, ca, cb, bc, bab, ca\}

b) \{bab, bc, ac, aba, ca, aac, ccb\}

c) \{cc, ca, cb, aba, bab, ac\}

d) \{aba, ac, cc, ca, cb, bc, bab, caa\}
\vspace{5pt} \par
\begin{tcolorbox}[colback=green!5, colframe=green!40!black, title=Solution]
Answer: d

Explanation: The string a state receives is the combination of all input alphabets which lie across the path covered.
\end{tcolorbox}
\end{mdframed}
\newpage

\begin{mdframed}[linewidth=1pt, roundcorner=5pt, linecolor=gray!30, backgroundcolor=gray!5]
\textbf{Question 7} 

Given L= \{X$\epsilon$$\Sigma$*= \{a, b\} |x has equal number of a, s and b's\}.

 Which of the following property satisfy the regularity of the given language?

a) Regularity is dependent upon the length of the string

b) Regularity is not dependent upon the length of the string

c) Can't be said for a particular string of a language

d) It may depend on the length of the string
\vspace{5pt} \par
\begin{tcolorbox}[colback=green!5, colframe=green!40!black, title=Solution]
Answer: b

Explanation: DFA can be made for infinite language with an infinite length. Thus, dependency over length is unfruitful.
\end{tcolorbox}
\end{mdframed}
\newpage

\begin{mdframed}[linewidth=1pt, roundcorner=5pt, linecolor=gray!30, backgroundcolor=gray!5]
\textbf{Question 8} 

Given:

L= \{x$\epsilon$$\Sigma$= \{0,1\} |x=0n1n for n$>$=1\}; Can there be a DFA possible for the language?

a) Yes

b) No
\vspace{5pt} \par
\begin{tcolorbox}[colback=green!5, colframe=green!40!black, title=Solution]
Answer: b

Explanation: It is not possible to have a count of equal number of 0 and 1 at any instant in DFA. Thus, It is not possible to build a DFA for the given Language.
\end{tcolorbox}
\end{mdframed}
\newpage

\begin{mdframed}[linewidth=1pt, roundcorner=5pt, linecolor=gray!30, backgroundcolor=gray!5]
\textbf{Question 9} 

$\delta$(A,1) = B, $\delta$(A,0) =A

     $\Delta$ (B, (0,1)) =C

     $\delta$(C,0) = A (Initial state =A)

String="011001" is transit at which of the states?

a) A

b) C

c) B

d) Invalid String
\vspace{5pt} \par
\begin{tcolorbox}[colback=green!5, colframe=green!40!black, title=Solution]
Answer: a

Explanation: It is east and simple to create the table and then the corresponding transition graph in order to get the result, at which state the given string would be accepted.
\end{tcolorbox}
\end{mdframed}
\newpage

\newtopic{Simpler Notations}

\begin{mdframed}[linewidth=1pt, roundcorner=5pt, linecolor=gray!30, backgroundcolor=gray!5]
\textbf{Question 1} 

Given Language: L= \{x$\epsilon$$\Sigma$= \{a, b\} |x has a substring ‘aa' in the production\}. Which of the corresponding representation notate the same?

a)


 
\begin{center}\includegraphics[width=0.5\textwidth]{automata-theory-questions-answers-simpler-notations-q1.png}\end{center}
 












b)


 
\begin{center}\includegraphics[width=0.5\textwidth]{automata-theory-questions-answers-simpler-notations-q1a.png}\end{center}
 


c)


 
\begin{center}\includegraphics[width=0.5\textwidth]{automata-theory-questions-answers-simpler-notations-q1b.png}\end{center}
 















d)


 
\begin{center}\includegraphics[width=0.5\textwidth]{automata-theory-questions-answers-simpler-notations-q1c.png}\end{center}
\vspace{5pt} \par
\begin{tcolorbox}[colback=green!5, colframe=green!40!black, title=Solution]
Answer: a

Explanation: The states transited has been written corresponding to the transitions as per the row and column. The row represents the transitions made and the ultimate.
\end{tcolorbox}
\end{mdframed}
\newpage

\begin{mdframed}[linewidth=1pt, roundcorner=5pt, linecolor=gray!30, backgroundcolor=gray!5]
\textbf{Question 2} 

Let u='1101', v='0001', then uv=11010001 and vu= 00011101.Using the given information what is the identity element for the string?

a) u
-1

b) v
-1

c) u
-1
v
-1

d) $\epsilon$
\vspace{5pt} \par
\begin{tcolorbox}[colback=green!5, colframe=green!40!black, title=Solution]
Answer: d

Explanation:  Identity relation: $\epsilon$w = w$\epsilon$ = w, thus the one satisfying the given relation will be the identity element.





     Apply Now for 
Free C Certification
 - April 2026
\end{tcolorbox}
\end{mdframed}
\newpage

\begin{mdframed}[linewidth=1pt, roundcorner=5pt, linecolor=gray!30, backgroundcolor=gray!5]
\textbf{Question 3} 

Which of the following substring will the following notation result?


 
\begin{center}\includegraphics[width=0.5\textwidth]{automata-theory-questions-answers-simpler-notations-q3.png}\end{center}
 

a) 0101011

b) 0101010

c) 010100

d) 100001
\vspace{5pt} \par
\begin{tcolorbox}[colback=green!5, colframe=green!40!black, title=Solution]
Answer: c

Explanation: The given DFA notation accepts the string of even length and prefix ‘01'.
\end{tcolorbox}
\end{mdframed}
\newpage

\begin{mdframed}[linewidth=1pt, roundcorner=5pt, linecolor=gray!30, backgroundcolor=gray!5]
\textbf{Question 4} 

Predict the following step in the given bunch of steps which accepts a strings which is of even length and has a prefix='01'

$\delta$ (q0, $\epsilon$) =q0 $<$ $\delta$(q0,0) =$\delta$ ($\delta$ (q0, $\epsilon$),0) = $\delta$(q0,0) =q1 $<$ \_\_\_\_\_\_\_\_\_\_\_\_\_\_\_

a) $\delta$ (q0, 011) =$\delta$ ($\delta$ (q0,1), 1) =$\delta$ (q2, 1) =q3

b) $\delta$ (q0, 01) =$\delta$ ($\delta$ (q0, 0), 1) = $\delta$ (q1, 1) =q2

c) $\delta$ (q0, 011) =$\delta$ ($\delta$ (q01, 1), 1) =$\delta$ (q2, 0) =q3

d) $\delta$ (q0, 0111) =$\delta$ ($\delta$ (q0, 011), 0) = $\delta$ (q3, 1) =q2
\vspace{5pt} \par
\begin{tcolorbox}[colback=green!5, colframe=green!40!black, title=Solution]
Answer: b

Explanation: Here, $\delta$ refers to transition function and results into new state or function when an transition is performed over its state.
\end{tcolorbox}
\end{mdframed}
\newpage

\begin{mdframed}[linewidth=1pt, roundcorner=5pt, linecolor=gray!30, backgroundcolor=gray!5]
\textbf{Question 5} 

Fill the missing blank in the given Transition Table:

Language L= \{x$\epsilon$$\Sigma$= \{0,1\} |x accepts all the binary strings not divisible by 3\}


 
\begin{center}\includegraphics[width=0.5\textwidth]{automata-theory-questions-answers-simpler-notations-q5.png}\end{center}
 

a) Q0

b) Q1

c) Q2

d) No Transition
\vspace{5pt} \par
\begin{tcolorbox}[colback=green!5, colframe=green!40!black, title=Solution]
Answer: b

Explanation: The tabular representation of DFA is quite readable and can be used to some ore complex problems. Here, we need to form the transition graph and fill up the given blank.
\end{tcolorbox}
\end{mdframed}
\newpage

\begin{mdframed}[linewidth=1pt, roundcorner=5pt, linecolor=gray!30, backgroundcolor=gray!5]
\textbf{Question 6} 

Which among the following is the missing transition in the given DFA?

L= \{x$\epsilon$$\Sigma$= \{a, b\} | x starts with a and ends with b\}


 
\begin{center}\includegraphics[width=0.5\textwidth]{automata-theory-questions-answers-simpler-notations-q6.png}\end{center}
 

a) $\delta$ (q0, a) = q0

b) $\delta$ (F, a) = q1

c) $\delta$ (F, a) = D

d) $\delta$ (q1, a) = D
\vspace{5pt} \par
\begin{tcolorbox}[colback=green!5, colframe=green!40!black, title=Solution]
Answer: b

Explanation:  For the given Language, the transition missing is $\delta$ (F, a) = q1.
\end{tcolorbox}
\end{mdframed}
\newpage

\begin{mdframed}[linewidth=1pt, roundcorner=5pt, linecolor=gray!30, backgroundcolor=gray!5]
\textbf{Question 7} 

The complement of a language will only be defined when and only when the \_\_\_\_\_\_\_\_\_\_ over the language is defined.

a) String

b) Word

c) Alphabet

d) Grammar
\vspace{5pt} \par
\begin{tcolorbox}[colback=green!5, colframe=green!40!black, title=Solution]
Answer: c

Explanation: It is not possible to define the complement of a language without defining the input alphabets. Example: A language which does not consist of substring ‘ab' while the complement would be the language which does contain a substring ‘ab'.
\end{tcolorbox}
\end{mdframed}
\newpage

\begin{mdframed}[linewidth=1pt, roundcorner=5pt, linecolor=gray!30, backgroundcolor=gray!5]
\textbf{Question 8} 

Which among the following is not notated as infinite language?

a) Palindrome

b) Reverse

c) Factorial

d) L=\{ab\}*
\vspace{5pt} \par
\begin{tcolorbox}[colback=green!5, colframe=green!40!black, title=Solution]
Answer: c

Explanation: Factorial, here is the most appropriate non-infinite domain. Otherwise, palindrome and reverse have infinite domains.
\end{tcolorbox}
\end{mdframed}
\newpage

\begin{mdframed}[linewidth=1pt, roundcorner=5pt, linecolor=gray!30, backgroundcolor=gray!5]
\textbf{Question 9} 

Which among the following states would be notated as the final state/acceptance state?

L= \{x$\epsilon$$\Sigma$= \{a, b\} | length of x is 2\}


 
\begin{center}\includegraphics[width=0.5\textwidth]{automata-theory-questions-answers-simpler-notations-q9.png}\end{center}
 

a) q1

b) q2

c) q1, q2

d) q3
\vspace{5pt} \par
\begin{tcolorbox}[colback=green!5, colframe=green!40!black, title=Solution]
Answer: b

Explanation: According to the given language, q2 Is to become the final/acceptance state in order to satisfy.
\end{tcolorbox}
\end{mdframed}
\newpage

\begin{mdframed}[linewidth=1pt, roundcorner=5pt, linecolor=gray!30, backgroundcolor=gray!5]
\textbf{Question 10} 

Which of the following are the final states in the given DFA according to the Language given.?

L= \{x$\epsilon$$\Sigma$= \{a, b\} |length of x is at most 2\}


 
\begin{center}\includegraphics[width=0.5\textwidth]{automata-theory-questions-answers-simpler-notations-q9.png}\end{center}
 

a) q0, q1

b) q0, q2

c) q1, q2

d) q0, q1, q2
\vspace{5pt} \par
\begin{tcolorbox}[colback=green!5, colframe=green!40!black, title=Solution]
Answer: d

Explanation: According to the given language, the length is at most 2, thus the answer is found accordingly.
\end{tcolorbox}
\end{mdframed}
\newpage

\newtopic{The Language of DFA}

\begin{mdframed}[linewidth=1pt, roundcorner=5pt, linecolor=gray!30, backgroundcolor=gray!5]
\textbf{Question 1} 

How many languages are over the alphabet R?

a) countably infinite

b) countably finite

c) uncountable finite

d) uncountable infinite
\vspace{5pt} \par
\begin{tcolorbox}[colback=green!5, colframe=green!40!black, title=Solution]
Answer: d

Explanation: A language over an alphabet R is a set of strings over A which is uncountable and infinite.
\end{tcolorbox}
\end{mdframed}
\newpage

\begin{mdframed}[linewidth=1pt, roundcorner=5pt, linecolor=gray!30, backgroundcolor=gray!5]
\textbf{Question 2} 

According to the 5-tuple representation i.e. FA= \{Q, $\Sigma$, $\delta$, q, F\}

Statement 1: q $\epsilon$ Q'; Statement 2: F$\epsilon$Q

a) Statement 1 is true, Statement 2 is false

b) Statement 1 is false, Statement 2 is true

c) Statement 1 is false, Statement 2 may be true

d) Statement 1 may be true, Statement 2 is false
\vspace{5pt} \par
\begin{tcolorbox}[colback=green!5, colframe=green!40!black, title=Solution]
Answer: b

Explanation: Q is the Finite set of states, whose elements i.e. the states constitute the finite automata.
\end{tcolorbox}
\end{mdframed}
\newpage

\begin{mdframed}[linewidth=1pt, roundcorner=5pt, linecolor=gray!30, backgroundcolor=gray!5]
\textbf{Question 3} 

$\delta$ˆ tells us the best:

a) how the DFA S behaves on a word u

b) the state is the dumping state

c) the final state has been reached

d) Kleene operation is performed on the set
\vspace{5pt} \par
\begin{tcolorbox}[colback=green!5, colframe=green!40!black, title=Solution]
Answer: a

Explanation: $\delta$ or the Transition function describes the best, how a DFA behaves on a string where to transit next, which direction to take.
\end{tcolorbox}
\end{mdframed}
\newpage

\begin{mdframed}[linewidth=1pt, roundcorner=5pt, linecolor=gray!30, backgroundcolor=gray!5]
\textbf{Question 4} 

Which of the following option is correct?

A = \{\{abc, aaba\}. \{$\epsilon$, a, bb\}\}

a) abcbb \textcent  A

b) $\epsilon$\textcent A

c) $\epsilon$ may not belong to A

d) abca \textcent  A
\vspace{5pt} \par
\begin{tcolorbox}[colback=green!5, colframe=green!40!black, title=Solution]
Answer: b

Explanation: As the question has dot operation, $\epsilon$ will not be a part of the concatenated set. Had it been a union operation, $\epsilon$ would be a part of the operated set.
\end{tcolorbox}
\end{mdframed}
\newpage

\begin{mdframed}[linewidth=1pt, roundcorner=5pt, linecolor=gray!30, backgroundcolor=gray!5]
\textbf{Question 5} 

For a DFA accepting binary numbers whose decimal equivalent is divisible by 4, what are all the possible remainders?

a) 0

b) 0,2

c) 0,2,4

d) 0,1,2,3
\vspace{5pt} \par
\begin{tcolorbox}[colback=green!5, colframe=green!40!black, title=Solution]
Answer: d

Explanation: All the decimal numbers on division would lead to only 4 remainders i.e. 0,1,2,3 (Property of Decimal division).





     
Free Certifications
 on 300 subjects are live for April 2026. Register Now!
\end{tcolorbox}
\end{mdframed}
\newpage

\begin{mdframed}[linewidth=1pt, roundcorner=5pt, linecolor=gray!30, backgroundcolor=gray!5]
\textbf{Question 6} 

Which of the following x is accepted by the given DFA (x is a binary string $\Sigma$= \{0,1\})?


 
\begin{center}\includegraphics[width=0.5\textwidth]{automata-theory-questions-answers-the-language-dfa-q6.png}\end{center}
 

a) divisible by 3

b) divisible by 2

c) divisible by 2 and 3

d) divisible by 3 and 2
\vspace{5pt} \par
\begin{tcolorbox}[colback=green!5, colframe=green!40!black, title=Solution]
Answer: d

Explanation: The given DFA accepts all the binary strings such that they are divisible by 3 and 2.Thus, it can be said that it also accepts all the strings which is divisible by 6.
\end{tcolorbox}
\end{mdframed}
\newpage

\begin{mdframed}[linewidth=1pt, roundcorner=5pt, linecolor=gray!30, backgroundcolor=gray!5]
\textbf{Question 7} 

Given:

L1= \{x$\epsilon$ $\Sigma$*|x contains even no's of 0's\}

L2= \{x$\epsilon$ $\Sigma$*|x contains odd no's of 1's\}

No of final states in Language L1 U L2?

a) 1

b) 2

c) 3

d) 4
\vspace{5pt} \par
\begin{tcolorbox}[colback=green!5, colframe=green!40!black, title=Solution]
Answer: c

Explanation:


 
\begin{center}\includegraphics[width=0.5\textwidth]{automata-theory-questions-answers-the-language-dfa-q7.png}\end{center}
\end{tcolorbox}
\end{mdframed}
\newpage

\begin{mdframed}[linewidth=1pt, roundcorner=5pt, linecolor=gray!30, backgroundcolor=gray!5]
\textbf{Question 8} 

The maximum number of transition which can be performed over a state in a DFA?

$\Sigma$= \{a, b, c\}

a) 1

b) 2

c) 3

d) 4
\vspace{5pt} \par
\begin{tcolorbox}[colback=green!5, colframe=green!40!black, title=Solution]
Answer: c

Explanation: The maximum number of transitions which a DFA allows for a language is the number of elements the transitions constitute.
\end{tcolorbox}
\end{mdframed}
\newpage

\begin{mdframed}[linewidth=1pt, roundcorner=5pt, linecolor=gray!30, backgroundcolor=gray!5]
\textbf{Question 9} 

The maximum sum of in degree and out degree over a state in a DFA can be determined as:

$\Sigma$= \{a, b, c, d\}

a) 4+4

b) 4+16

c) 4+0

d) depends on the Language
\vspace{5pt} \par
\begin{tcolorbox}[colback=green!5, colframe=green!40!black, title=Solution]
Answer: d

Explanation: The out degree for a DFA I fixed while the in degree depends on the number of states in the DFA and that cannot be determined without the dependence over the Language.
\end{tcolorbox}
\end{mdframed}
\newpage

\begin{mdframed}[linewidth=1pt, roundcorner=5pt, linecolor=gray!30, backgroundcolor=gray!5]
\textbf{Question 10} 

The sum of minimum and maximum number of final states for a DFA n states is equal to:

a) n+1

b) n

c) n-1

d) n+2
\vspace{5pt} \par
\begin{tcolorbox}[colback=green!5, colframe=green!40!black, title=Solution]
Answer: a

Explanation: The maximum number of final states for a DFA can be total number of states itself and minimum would always be 1, as no DFA exits without a final state. Therefore, the solution is n+1.
\end{tcolorbox}
\end{mdframed}
\newpage

\newtopic{Finite Automata}

\begin{mdframed}[linewidth=1pt, roundcorner=5pt, linecolor=gray!30, backgroundcolor=gray!5]
\textbf{Question 1} 

There are \_\_\_\_\_\_\_\_ tuples in finite state machine.

a) 4

b) 5

c) 6

d) unlimited
\vspace{5pt} \par
\begin{tcolorbox}[colback=green!5, colframe=green!40!black, title=Solution]
Answer:b

Explanation: States, input symbols, initial state, accepting state and transition function.
\end{tcolorbox}
\end{mdframed}
\newpage

\begin{mdframed}[linewidth=1pt, roundcorner=5pt, linecolor=gray!30, backgroundcolor=gray!5]
\textbf{Question 2} 

Transition function maps.

a) $\Sigma$ * Q -$>$ $\Sigma$

b) Q * Q -$>$ $\Sigma$

c) $\Sigma$ * $\Sigma$ -$>$ Q

d) Q * $\Sigma$ -$>$ Q
\vspace{5pt} \par
\begin{tcolorbox}[colback=green!5, colframe=green!40!black, title=Solution]
Answer:d

Explanation: Inputs are state and input string output is states.
\end{tcolorbox}
\end{mdframed}
\newpage

\begin{mdframed}[linewidth=1pt, roundcorner=5pt, linecolor=gray!30, backgroundcolor=gray!5]
\textbf{Question 3} 

Number of states require to accept string ends with 10.

a) 3

b) 2

c) 1

d) can't be represented.
\vspace{5pt} \par
\begin{tcolorbox}[colback=green!5, colframe=green!40!black, title=Solution]
Answer:a

Explanation: This is minimal finite automata.
\end{tcolorbox}
\end{mdframed}
\newpage

\begin{mdframed}[linewidth=1pt, roundcorner=5pt, linecolor=gray!30, backgroundcolor=gray!5]
\textbf{Question 4} 

Extended transition function is .

a) Q * $\Sigma$* -$>$ Q

b) Q * $\Sigma$ -$>$ Q

c) Q* * $\Sigma$* -$>$ $\Sigma$

d) Q * $\Sigma$ -$>$ $\Sigma$
\vspace{5pt} \par
\begin{tcolorbox}[colback=green!5, colframe=green!40!black, title=Solution]
Answer:a

Explanation: This takes single state and string of input to produce a state.
\end{tcolorbox}
\end{mdframed}
\newpage

\begin{mdframed}[linewidth=1pt, roundcorner=5pt, linecolor=gray!30, backgroundcolor=gray!5]
\textbf{Question 5} 

$\delta$*(q,ya) is equivalent to .

a) $\delta$((q,y),a)

b) $\delta$($\delta$*(q,y),a)

c) $\delta$(q,ya)

d) independent from $\delta$ notation
\vspace{5pt} \par
\begin{tcolorbox}[colback=green!5, colframe=green!40!black, title=Solution]
Answer:b

Explanation: First it parse y string after that it parse a.
\end{tcolorbox}
\end{mdframed}
\newpage

\begin{mdframed}[linewidth=1pt, roundcorner=5pt, linecolor=gray!30, backgroundcolor=gray!5]
\textbf{Question 6} 

String X is accepted by finite automata if .

a) $\delta$*(q,x) E A

b) $\delta$(q,x) E A

c) $\delta$*(Q0,x) E A

d) $\delta$(Q0,x) E A
\vspace{5pt} \par
\begin{tcolorbox}[colback=green!5, colframe=green!40!black, title=Solution]
Answer:c

Explanation: If automata starts with starting state and after finite moves if reaches to final step then it called accepted.
\end{tcolorbox}
\end{mdframed}
\newpage

\begin{mdframed}[linewidth=1pt, roundcorner=5pt, linecolor=gray!30, backgroundcolor=gray!5]
\textbf{Question 7} 

Languages of a automata is

a) If it is accepted by automata

b) If it halts

c) If automata touch final state in its life time

d) All language are language of automata
\vspace{5pt} \par
\begin{tcolorbox}[colback=green!5, colframe=green!40!black, title=Solution]
Answer:a

Explanation: If a string accepted by automata it is called language of automata.
\end{tcolorbox}
\end{mdframed}
\newpage

\begin{mdframed}[linewidth=1pt, roundcorner=5pt, linecolor=gray!30, backgroundcolor=gray!5]
\textbf{Question 8} 

Language of finite automata is.

a) Type 0

b) Type 1

c) Type 2

d) Type 3
\vspace{5pt} \par
\begin{tcolorbox}[colback=green!5, colframe=green!40!black, title=Solution]
Answer:d

Explanation: According to Chomsky classification.
\end{tcolorbox}
\end{mdframed}
\newpage

\begin{mdframed}[linewidth=1pt, roundcorner=5pt, linecolor=gray!30, backgroundcolor=gray!5]
\textbf{Question 9} 

Finite automata requires minimum \_\_\_\_\_\_\_ number of stacks.

a) 1

b) 0

c) 2

d) None of the mentioned
\vspace{5pt} \par
\begin{tcolorbox}[colback=green!5, colframe=green!40!black, title=Solution]
Answer:b

Explanation: Finite automata doesn't require any stack operation.
\end{tcolorbox}
\end{mdframed}
\newpage

\begin{mdframed}[linewidth=1pt, roundcorner=5pt, linecolor=gray!30, backgroundcolor=gray!5]
\textbf{Question 10} 

Number of final state require to accept $\Phi$ in minimal finite automata.

a) 1

b) 2

c) 3

d) None of the mentioned
\vspace{5pt} \par
\begin{tcolorbox}[colback=green!5, colframe=green!40!black, title=Solution]
Answer:d

Explanation: No final state requires.
\end{tcolorbox}
\end{mdframed}
\newpage

\begin{mdframed}[linewidth=1pt, roundcorner=5pt, linecolor=gray!30, backgroundcolor=gray!5]
\textbf{Question 11} 

Regular expression for all strings starts with ab and ends with bba is.

a) aba*b*bba

b) ab(ab)*bba

c) ab(a+b)*bba

d) All of the mentioned
\vspace{5pt} \par
\begin{tcolorbox}[colback=green!5, colframe=green!40!black, title=Solution]
Answer:c

Explanation: Starts with ab then any number of a or b and ends with bba.
\end{tcolorbox}
\end{mdframed}
\newpage

\begin{mdframed}[linewidth=1pt, roundcorner=5pt, linecolor=gray!30, backgroundcolor=gray!5]
\textbf{Question 12} 

How many DFA's exits with two states over input alphabet \{0,1\} ?

a) 16

b) 26

c) 32

d) 64
\vspace{5pt} \par
\begin{tcolorbox}[colback=green!5, colframe=green!40!black, title=Solution]
Answer:d

Explanation: Number of DFA's = 2
n
 * n
(2*n)
.
\end{tcolorbox}
\end{mdframed}
\newpage

\begin{mdframed}[linewidth=1pt, roundcorner=5pt, linecolor=gray!30, backgroundcolor=gray!5]
\textbf{Question 13} 

The basic limitation of finite automata is that

a) It can't remember arbitrary large amount of information.

b) It sometimes recognize grammar that are not regular.

c) It sometimes fails to recognize regular grammar.

d) All of the mentioned
\vspace{5pt} \par
\begin{tcolorbox}[colback=green!5, colframe=green!40!black, title=Solution]
Answer:a

Explanation:Because there is no memory associated with automata.
\end{tcolorbox}
\end{mdframed}
\newpage

\begin{mdframed}[linewidth=1pt, roundcorner=5pt, linecolor=gray!30, backgroundcolor=gray!5]
\textbf{Question 14} 

Number of states require to simulate a computer with memory capable of storing ‘3' words each of length ‘8'.

a) 3 * 2
8

b) 2
(3*8)

c) 2
(3+8)

d) None of the mentioned
\vspace{5pt} \par
\begin{tcolorbox}[colback=green!5, colframe=green!40!black, title=Solution]
Answer:b

Explanation: 2
(m*n)
 states requires.
\end{tcolorbox}
\end{mdframed}
\newpage

\begin{mdframed}[linewidth=1pt, roundcorner=5pt, linecolor=gray!30, backgroundcolor=gray!5]
\textbf{Question 15} 

FSM with output capability can be used to add two given integer in binary representation. This is

a) True

b) False

c) May be true

d) None of the mentioned
\vspace{5pt} \par
\begin{tcolorbox}[colback=green!5, colframe=green!40!black, title=Solution]
Answer:a

Explanation: Use them as a flip flop output.
\end{tcolorbox}
\end{mdframed}
\newpage

\newtopic{Non Deterministic Finite Automata-Introduction}

\begin{mdframed}[linewidth=1pt, roundcorner=5pt, linecolor=gray!30, backgroundcolor=gray!5]
\textbf{Question 1} 

Which of the following options is correct?

Statement 1: Initial State of NFA is Initial State of DFA.

Statement 2: The final state of DFA will be every combination of final state of NFA.

a) Statement 1 is true and Statement 2 is true

b) Statement 1 is true and Statement 2 is false

c) Statement 1 can be true and Statement 2 is true

d) Statement 1 is false and Statement 2 is also false
\vspace{5pt} \par
\begin{tcolorbox}[colback=green!5, colframe=green!40!black, title=Solution]
Answer: a

Explanation: Statement 1 and 2 always true for a given Language.
\end{tcolorbox}
\end{mdframed}
\newpage

\begin{mdframed}[linewidth=1pt, roundcorner=5pt, linecolor=gray!30, backgroundcolor=gray!5]
\textbf{Question 2} 

Given Language: L= \{ab U aba\}*

If X is the minimum number of states for a DFA and Y is the number of states to construct the NFA,

|X-Y|=?

a) 2

b) 3

c) 4

d) 1
\vspace{5pt} \par
\begin{tcolorbox}[colback=green!5, colframe=green!40!black, title=Solution]
Answer: a

Explanation: Construct the DFA and NFA individually, and the attain the difference of states.
\end{tcolorbox}
\end{mdframed}
\newpage

\begin{mdframed}[linewidth=1pt, roundcorner=5pt, linecolor=gray!30, backgroundcolor=gray!5]
\textbf{Question 3} 

An automaton that presents output based on previous state or current input:

a) Acceptor

b) Classifier

c) Transducer

d) None of the mentioned.
\vspace{5pt} \par
\begin{tcolorbox}[colback=green!5, colframe=green!40!black, title=Solution]
Answer: c

Explanation: A transducer is an automaton that produces an output on the basis of what input has been given currently or previous state.
\end{tcolorbox}
\end{mdframed}
\newpage

\begin{mdframed}[linewidth=1pt, roundcorner=5pt, linecolor=gray!30, backgroundcolor=gray!5]
\textbf{Question 4} 

If NFA of 6 states excluding the initial state is converted into DFA, maximum possible number of states for the DFA is ?

a) 64

b) 32

c) 128

d) 127
\vspace{5pt} \par
\begin{tcolorbox}[colback=green!5, colframe=green!40!black, title=Solution]
Answer: c

Explanation: The maximum number of sets for DFA converted from NFA would be not greater than 2n.
\end{tcolorbox}
\end{mdframed}
\newpage

\begin{mdframed}[linewidth=1pt, roundcorner=5pt, linecolor=gray!30, backgroundcolor=gray!5]
\textbf{Question 5} 

NFA, in its name has 'non-deterministic' because of :

a) The result is undetermined

b) The choice of path is non-deterministic

c) The state to be transited next is non-deterministic

d) All of the mentioned
\vspace{5pt} \par
\begin{tcolorbox}[colback=green!5, colframe=green!40!black, title=Solution]
Answer: b

Explanation: Non deterministic or deterministic depends upon the definite path defined for the transition from one state to another or undefined(multiple paths).





     
Free Certifications
 on 300 subjects are live for April 2026. Register Now!
\end{tcolorbox}
\end{mdframed}
\newpage

\begin{mdframed}[linewidth=1pt, roundcorner=5pt, linecolor=gray!30, backgroundcolor=gray!5]
\textbf{Question 6} 

Which of the following is correct proposition?

Statement 1: Non determinism is a generalization of Determinism.

Statement 2: Every DFA is automatically an NFA

a) Statement 1 is correct because Statement 2 is correct

b) Statement 2 is correct because Statement 2 is correct

c) Statement 2 is false and Statement 1 is false

d) Statement 1 is false because Statement 2 is false
\vspace{5pt} \par
\begin{tcolorbox}[colback=green!5, colframe=green!40!black, title=Solution]
Answer: b

Explanation: DFA is a specific case of NFA.
\end{tcolorbox}
\end{mdframed}
\newpage

\begin{mdframed}[linewidth=1pt, roundcorner=5pt, linecolor=gray!30, backgroundcolor=gray!5]
\textbf{Question 7} 

Given Language L= \{x$\epsilon$ \{a, b\}*|x contains aba as its substring\}

Find the difference of transitions made in constructing a DFA and an equivalent NFA?

a) 2

b) 3

c) 4

d) Cannot be determined
\vspace{5pt} \par
\begin{tcolorbox}[colback=green!5, colframe=green!40!black, title=Solution]
Answer: a

Explanation: The individual Transition graphs can be made and the difference of transitions can be determined.
\end{tcolorbox}
\end{mdframed}
\newpage

\begin{mdframed}[linewidth=1pt, roundcorner=5pt, linecolor=gray!30, backgroundcolor=gray!5]
\textbf{Question 8} 

The construction time for DFA from an equivalent NFA (m number of node)is:

a) O(m
2
)

b) O(2
m
)

c) O(m)

d) O(log m)
\vspace{5pt} \par
\begin{tcolorbox}[colback=green!5, colframe=green!40!black, title=Solution]
Answer: b

Explanation: From the coded NFA-DFA conversion.
\end{tcolorbox}
\end{mdframed}
\newpage

\begin{mdframed}[linewidth=1pt, roundcorner=5pt, linecolor=gray!30, backgroundcolor=gray!5]
\textbf{Question 9} 

If n is the length of Input string and m is the number of nodes, the running time of DFA  is x that of NFA.Find x?

a) 1/m
2

b) 2
m

c) 1/m

d) log m
\vspace{5pt} \par
\begin{tcolorbox}[colback=green!5, colframe=green!40!black, title=Solution]
Answer: a

Explanation: Running time of DFA: O(n) and Running time of NFA =O(m
2
n).
\end{tcolorbox}
\end{mdframed}
\newpage

\begin{mdframed}[linewidth=1pt, roundcorner=5pt, linecolor=gray!30, backgroundcolor=gray!5]
\textbf{Question 10} 

Which of the following option is correct?

a) NFA is slower to process and its representation uses more memory than DFA

b) DFA is faster to process and its representation uses less memory than NFA

c) NFA is slower to process and its representation uses less memory than DFA

d) DFA is slower to process and its representation uses less memory than NFA
\vspace{5pt} \par
\begin{tcolorbox}[colback=green!5, colframe=green!40!black, title=Solution]
Answer: c

Explanation: NFA, while computing strings, take parallel paths, make different copies of input and goes along different paths in order to search for the result. This creates the difference in processing speed of DFA and NFA.
\end{tcolorbox}
\end{mdframed}
\newpage

\newtopic{Extended Transition Function}

\begin{mdframed}[linewidth=1pt, roundcorner=5pt, linecolor=gray!30, backgroundcolor=gray!5]
\textbf{Question 1} 

The number of tuples in an extended Non Deterministic Finite Automaton:

a) 5

b) 6

c) 7

d) 4
\vspace{5pt} \par
\begin{tcolorbox}[colback=green!5, colframe=green!40!black, title=Solution]
Answer: a

Explanation: For NFA or extended transition function on NFA, the tuple elements remains same i.e. 5.
\end{tcolorbox}
\end{mdframed}
\newpage

\begin{mdframed}[linewidth=1pt, roundcorner=5pt, linecolor=gray!30, backgroundcolor=gray!5]
\textbf{Question 2} 

Choose the correct option for the given statement:

Statement: The DFA shown represents all strings which has 1 at second last position.


 
\begin{center}\includegraphics[width=0.5\textwidth]{automata-theory-extended-transition-function-q2.png}\end{center}
 

a) Correct

b) Incorrect, Incomplete DFA

c) Wrong proposition

d) May be correct
\vspace{5pt} \par
\begin{tcolorbox}[colback=green!5, colframe=green!40!black, title=Solution]
Answer: c

Explanation: The given figure is an NFA. The statement contradicts itself.
\end{tcolorbox}
\end{mdframed}
\newpage

\begin{mdframed}[linewidth=1pt, roundcorner=5pt, linecolor=gray!30, backgroundcolor=gray!5]
\textbf{Question 3} 

What is wrong in the given definition?

Def: (\{q0, q1, q2\}, \{0,1\}, $\delta$, q3, \{q3\})

a) The definition does not satisfy 5 Tuple definition of NFA

b) There are no transition definition

c) Initial and Final states do not belong to the Graph

d) Initial and final states can't be same
\vspace{5pt} \par
\begin{tcolorbox}[colback=green!5, colframe=green!40!black, title=Solution]
Answer: c

Explanation: q3 does not belong to Q where Q= set of finite states.
\end{tcolorbox}
\end{mdframed}
\newpage

\begin{mdframed}[linewidth=1pt, roundcorner=5pt, linecolor=gray!30, backgroundcolor=gray!5]
\textbf{Question 4} 

If $\delta$ is the transition function for a given NFA, then we define the $\delta$' for the DFA accepting the same language would be:

Note: S is a subset of Q and a is a symbol.

a) $\delta$' (S, a) =U
p$\epsilon$s
 $\delta$ (p, a)

b) $\delta$' (S, a) =U
p$\neq$s
 $\delta$ (p, a)

c) $\delta$' (S, a) =U
p$\epsilon$s
 $\delta$(p)

d) $\delta$' (S) =U
p$\neq$s
 $\delta$(p)
\vspace{5pt} \par
\begin{tcolorbox}[colback=green!5, colframe=green!40!black, title=Solution]
Answer: a

Explanation: According to subset construction, equation 1 holds true.
\end{tcolorbox}
\end{mdframed}
\newpage

\begin{mdframed}[linewidth=1pt, roundcorner=5pt, linecolor=gray!30, backgroundcolor=gray!5]
\textbf{Question 5} 

What is the relation between DFA and NFA on the basis of computational power?

a) DFA $>$ NFA

b) NFA $>$ DFA

c) Equal

d) Can't be said
\vspace{5pt} \par
\begin{tcolorbox}[colback=green!5, colframe=green!40!black, title=Solution]
Answer: c

Explanation: DFA is said to be a specific case of NFA and for every NFA that exists for a given language, an equivalent DFA also exists.





     Apply Now for 
Free C Certification
 - April 2026
\end{tcolorbox}
\end{mdframed}
\newpage

\begin{mdframed}[linewidth=1pt, roundcorner=5pt, linecolor=gray!30, backgroundcolor=gray!5]
\textbf{Question 6} 

If a string S is accepted by a finite state automaton, S=s
1
s
2
s
3
……s
n
 where s
i
$\epsilon$$\Sigma$ and there exists a sequence of states r0, r1, r2…… rn such that $\delta$(r(i), s
i+1
) =r
i+1
 for each 0, 1, …n-1, then r(n) is:

a) initial state

b) transition symbol

c) accepting state

d) intermediate state
\vspace{5pt} \par
\begin{tcolorbox}[colback=green!5, colframe=green!40!black, title=Solution]
Answer: c

Explanation: r(n) is the final state and accepts the string S after the string being traversed through r(i) other states where I $\epsilon$ 01,2…(n-2).
\end{tcolorbox}
\end{mdframed}
\newpage

\begin{mdframed}[linewidth=1pt, roundcorner=5pt, linecolor=gray!30, backgroundcolor=gray!5]
\textbf{Question 7} 

According to the given table, compute the number of transitions with 1 as its symbol but not 0:


 
\begin{center}\includegraphics[width=0.5\textwidth]{automata-theory-questions-answers-extended-transition-function-q7.png}\end{center}
 

a) 4

b) 3

c) 2

d) 1
\vspace{5pt} \par
\begin{tcolorbox}[colback=green!5, colframe=green!40!black, title=Solution]
Answer: d

Explanation: The transition graph is made and thus the answer can be found.
\end{tcolorbox}
\end{mdframed}
\newpage

\begin{mdframed}[linewidth=1pt, roundcorner=5pt, linecolor=gray!30, backgroundcolor=gray!5]
\textbf{Question 8} 

From the given table, $\delta$*(q0, 011) =?


 
\begin{center}\includegraphics[width=0.5\textwidth]{automata-theory-questions-answers-extended-transition-function-q7.png}\end{center}
 

a) \{q0\}

b) \{q1\} U \{q0, q1, q2\}

c) \{q2, q1\}

d) \{q3, q1, q2, q0\}
\vspace{5pt} \par
\begin{tcolorbox}[colback=green!5, colframe=green!40!black, title=Solution]
Answer: b

Explanation: $\delta$*(q0,011) = U
r$\epsilon$
$\delta$*(q0,01) $\delta$ (r, 1) = \{q0, q1, q2\}.
\end{tcolorbox}
\end{mdframed}
\newpage

\begin{mdframed}[linewidth=1pt, roundcorner=5pt, linecolor=gray!30, backgroundcolor=gray!5]
\textbf{Question 9} 

Number of times the state q3 or q2 is being a part of extended 6 transition state is


 
\begin{center}\includegraphics[width=0.5\textwidth]{automata-theory-questions-answers-extended-transition-function-q9.png}\end{center}
 

a) 6

b) 5

c) 4

d) 7
\vspace{5pt} \par
\begin{tcolorbox}[colback=green!5, colframe=green!40!black, title=Solution]
Answer: a

Explanation: According to the question, presence of q2 or q1 would count so it does and the answer according to the diagram is 6.
\end{tcolorbox}
\end{mdframed}
\newpage

\begin{mdframed}[linewidth=1pt, roundcorner=5pt, linecolor=gray!30, backgroundcolor=gray!5]
\textbf{Question 10} 

Predict the missing procedure:


 
\begin{center}\includegraphics[width=0.5\textwidth]{automata-theory-questions-answers-extended-transition-function-q10.png}\end{center}
 

i.$\Delta$(Q0, $\epsilon$) =\{Q0\},

ii.$\Delta$(Q0, 01) = \{Q0, Q1\}

iii.$\delta$(Q0, 010) =?

a) \{Q0, Q1, Q2\}

b) \{Q0, Q1\}

c) \{Q0, Q2\}

d) \{Q1, Q2\}
\vspace{5pt} \par
\begin{tcolorbox}[colback=green!5, colframe=green!40!black, title=Solution]
Answer: c

Explanation: According to given table and extended transition state implementation, we can find the state at which it rests.
\end{tcolorbox}
\end{mdframed}
\newpage

\newtopic{The Language of NFA}

\begin{mdframed}[linewidth=1pt, roundcorner=5pt, linecolor=gray!30, backgroundcolor=gray!5]
\textbf{Question 1} 

Subset Construction method refers to:

a) Conversion of NFA to DFA

b) DFA minimization

c) Eliminating Null references

d) $\epsilon$-NFA to NFA
\vspace{5pt} \par
\begin{tcolorbox}[colback=green!5, colframe=green!40!black, title=Solution]
Answer: a

Explanation:  The conversion of a non-deterministic automata into a deterministic one is a process we call subset construction or power set construction.
\end{tcolorbox}
\end{mdframed}
\newpage

\begin{mdframed}[linewidth=1pt, roundcorner=5pt, linecolor=gray!30, backgroundcolor=gray!5]
\textbf{Question 2} 

Given Language:

L
n
= \{x$\epsilon$ \{0,1\} * | |x|$\geq$n, nth symbol from the right in x is 1\}

How many state are required to execute L
3
 using NFA?

a) 16

b) 15

c) 8

d) 7
\vspace{5pt} \par
\begin{tcolorbox}[colback=green!5, colframe=green!40!black, title=Solution]
Answer: b

Explanation: The finite automaton for the given language is made and thus, the answer can be obtained.
\end{tcolorbox}
\end{mdframed}
\newpage

\begin{mdframed}[linewidth=1pt, roundcorner=5pt, linecolor=gray!30, backgroundcolor=gray!5]
\textbf{Question 3} 

Which of the following does the given NFA represent?


 
\begin{center}\includegraphics[width=0.5\textwidth]{automata-theory-questions-answers-the-language-nfa-q3.png}\end{center}
 

a) \{11, 101\} * \{01\}

b) \{110, 01\} * \{11\}

c) \{11, 110\} * \{0\}

d) \{00, 110\} * \{1\}
\vspace{5pt} \par
\begin{tcolorbox}[colback=green!5, colframe=green!40!black, title=Solution]
Answer: c

Explanation: The given diagram can be analysed and thus the option can be seeked.
\end{tcolorbox}
\end{mdframed}
\newpage

\begin{mdframed}[linewidth=1pt, roundcorner=5pt, linecolor=gray!30, backgroundcolor=gray!5]
\textbf{Question 4} 

The number of transitions required to convert the following into equivalents DFA:


 
\begin{center}\includegraphics[width=0.5\textwidth]{automata-theory-questions-answers-the-language-nfa-q4.png}\end{center}
 

a) 2

b) 3

c) 1

d) 0
\vspace{5pt} \par
\begin{tcolorbox}[colback=green!5, colframe=green!40!black, title=Solution]
Answer: a

Explanation:


 
\begin{center}\includegraphics[width=0.5\textwidth]{automata-theory-questions-answers-the-language-nfa-q4a.png}\end{center}
\end{tcolorbox}
\end{mdframed}
\newpage

\begin{mdframed}[linewidth=1pt, roundcorner=5pt, linecolor=gray!30, backgroundcolor=gray!5]
\textbf{Question 5} 

If L is a regular language, L
c
 and L
r
 both will be:

a) Accepted by NFA

b) Rejected by NFA

c) One of them will be accepted

d) Cannot be said
\vspace{5pt} \par
\begin{tcolorbox}[colback=green!5, colframe=green!40!black, title=Solution]
Answer: a

Explanation: If L is a regular Language, L
c
 and L
r
  both are regular even.
\end{tcolorbox}
\end{mdframed}
\newpage

\begin{mdframed}[linewidth=1pt, roundcorner=5pt, linecolor=gray!30, backgroundcolor=gray!5]
\textbf{Question 6} 

In NFA, this very state is like dead-end non final state:

a) ACCEPT

b) REJECT

c) DISTINCT

d) START
\vspace{5pt} \par
\begin{tcolorbox}[colback=green!5, colframe=green!40!black, title=Solution]
Answer: b

Explanation: REJECT state will be like a halting state which rejects a particular invalid input.
\end{tcolorbox}
\end{mdframed}
\newpage

\begin{mdframed}[linewidth=1pt, roundcorner=5pt, linecolor=gray!30, backgroundcolor=gray!5]
\textbf{Question 7} 

We can represent one language in more one FSMs, true or false?

a) TRUE

b) FALSE

c) May be true

d) Cannot be said
\vspace{5pt} \par
\begin{tcolorbox}[colback=green!5, colframe=green!40!black, title=Solution]
Answer: a

Explanation: We can represent one language in more one FSMs, example for a same language we have a DFA and an equivalent NFA.
\end{tcolorbox}
\end{mdframed}
\newpage

\begin{mdframed}[linewidth=1pt, roundcorner=5pt, linecolor=gray!30, backgroundcolor=gray!5]
\textbf{Question 8} 

The production of form non-terminal -$>$ $\epsilon$ is called:

a) Sigma Production

b) Null Production

c) Epsilon Production

d) All of the mentioned
\vspace{5pt} \par
\begin{tcolorbox}[colback=green!5, colframe=green!40!black, title=Solution]
Answer: b

Explanation: The production of form non-terminal -$>$$\epsilon$ is call null production.
\end{tcolorbox}
\end{mdframed}
\newpage

\begin{mdframed}[linewidth=1pt, roundcorner=5pt, linecolor=gray!30, backgroundcolor=gray!5]
\textbf{Question 9} 

Which of the following is a regular language?

a) String whose length is a sequence of prime numbers

b) String with substring ww
r
 in between

c) Palindrome string

d) String with even number of Zero's
\vspace{5pt} \par
\begin{tcolorbox}[colback=green!5, colframe=green!40!black, title=Solution]
Answer: d

Explanation: DFSM's for the first three option is not possible; hence they aren't regular.
\end{tcolorbox}
\end{mdframed}
\newpage

\begin{mdframed}[linewidth=1pt, roundcorner=5pt, linecolor=gray!30, backgroundcolor=gray!5]
\textbf{Question 10} 

Which of the following recognizes the same formal language as of DFA and NFA?

a) Power set Construction

b) Subset Construction

c) Robin-Scott Construction

d) All of the mentioned
\vspace{5pt} \par
\begin{tcolorbox}[colback=green!5, colframe=green!40!black, title=Solution]
Answer: d

Explanation: All the three option refers to same technique if distinguishing similar constructions for different type of automata.
\end{tcolorbox}
\end{mdframed}
\newpage

\newtopic{Equivalence of NFA and DFA}

\begin{mdframed}[linewidth=1pt, roundcorner=5pt, linecolor=gray!30, backgroundcolor=gray!5]
\textbf{Question 1} 

Under which of the following operation, NFA is not closed?

a) Negation

b) Kleene

c) Concatenation

d) None of the mentioned
\vspace{5pt} \par
\begin{tcolorbox}[colback=green!5, colframe=green!40!black, title=Solution]
Answer: d

Explanation: NFA is said to be closed under the following operations:

a) Union

b) Intersection

c) Concatenation

d) Kleene

e) Negation
\end{tcolorbox}
\end{mdframed}
\newpage

\begin{mdframed}[linewidth=1pt, roundcorner=5pt, linecolor=gray!30, backgroundcolor=gray!5]
\textbf{Question 2} 

It is less complex to prove the closure properties over regular languages using

a) NFA

b) DFA

c) PDA

d) Can't be said
\vspace{5pt} \par
\begin{tcolorbox}[colback=green!5, colframe=green!40!black, title=Solution]
Answer: a

Explanation: We use the construction method to prove the validity of closure properties of regular languages. Thus, it can be observe, how tedious and complex is the construction of a DFA as compared to an NFA with respect to space.
\end{tcolorbox}
\end{mdframed}
\newpage

\begin{mdframed}[linewidth=1pt, roundcorner=5pt, linecolor=gray!30, backgroundcolor=gray!5]
\textbf{Question 3} 

Which of the following is an application of Finite Automaton?

a) Compiler Design

b) Grammar Parsers

c) Text Search

d) All of the mentioned
\vspace{5pt} \par
\begin{tcolorbox}[colback=green!5, colframe=green!40!black, title=Solution]
Answer: d

Explanation: There are many applications of finite automata, mainly in the field of Compiler Design and Parsers and Search Engines.
\end{tcolorbox}
\end{mdframed}
\newpage

\begin{mdframed}[linewidth=1pt, roundcorner=5pt, linecolor=gray!30, backgroundcolor=gray!5]
\textbf{Question 4} 

John is asked to make an automaton which accepts a given string for all the occurrence of ‘1001' in it. How many number of transitions would John use such that, the string processing application works?

a) 9

b) 11

c) 12

d) 15
\vspace{5pt} \par
\begin{tcolorbox}[colback=green!5, colframe=green!40!black, title=Solution]
Answer: a

Explanation:


 
\begin{center}\includegraphics[width=0.5\textwidth]{automata-theory-questions-answers-equivalence-nfa-dfa-q4.png}\end{center}
\end{tcolorbox}
\end{mdframed}
\newpage

\begin{mdframed}[linewidth=1pt, roundcorner=5pt, linecolor=gray!30, backgroundcolor=gray!5]
\textbf{Question 5} 

Which of the following do we use to form an NFA from a regular expression?

a) Subset Construction Method

b) Power Set Construction Method

c) Thompson Construction Method

d) Scott Construction Method
\vspace{5pt} \par
\begin{tcolorbox}[colback=green!5, colframe=green!40!black, title=Solution]
Answer: c

Explanation: Thompson Construction method is used to turn a regular expression in an NFA by fragmenting the given regular expression through the operations performed on the input alphabets.





     Join Sanfoundry classes at 
Telegram
 or 
Youtube
\end{tcolorbox}
\end{mdframed}
\newpage

\begin{mdframed}[linewidth=1pt, roundcorner=5pt, linecolor=gray!30, backgroundcolor=gray!5]
\textbf{Question 6} 

Which among the following can be an example of application of finite state machine(FSM)?

a) Communication Link

b) Adder

c) Stack

d) None of the mentioned
\vspace{5pt} \par
\begin{tcolorbox}[colback=green!5, colframe=green!40!black, title=Solution]
Answer: a

Explanation: Idle is the state when data in form of packets is send and returns if NAK is received else waits for the NAK to be received.
\end{tcolorbox}
\end{mdframed}
\newpage

\begin{mdframed}[linewidth=1pt, roundcorner=5pt, linecolor=gray!30, backgroundcolor=gray!5]
\textbf{Question 7} 

Which among the following is not an application of FSM?

a) Lexical Analyser

b) BOT

c) State charts

d) None of the mentioned
\vspace{5pt} \par
\begin{tcolorbox}[colback=green!5, colframe=green!40!black, title=Solution]
Answer: d

Explanation: Finite state automation is used in Lexical Analyser, Computer BOT (used in games), State charts, etc.
\end{tcolorbox}
\end{mdframed}
\newpage

\begin{mdframed}[linewidth=1pt, roundcorner=5pt, linecolor=gray!30, backgroundcolor=gray!5]
\textbf{Question 8} 

L1= \{w | w does not contain the string tr \}

L2 = \{w | w does contain the string tr\}

Given $\Sigma$ = \{t, r\}, The difference of the minimum number of states required to form L1 and L2?

a) 0

b) 1

c) 2

d) Cannot be said
\vspace{5pt} \par
\begin{tcolorbox}[colback=green!5, colframe=green!40!black, title=Solution]
Answer: a

Explanation:


 
\begin{center}\includegraphics[width=0.5\textwidth]{automata-theory-questions-answers-equivalence-nfa-dfa-q8.png}\end{center}
\end{tcolorbox}
\end{mdframed}
\newpage

\begin{mdframed}[linewidth=1pt, roundcorner=5pt, linecolor=gray!30, backgroundcolor=gray!5]
\textbf{Question 9} 

Predict the number of transitions required to automate the following language using only 3 states:

L = \{w | w ends with 00\}

a) 3

b) 2

c) 4

d) Cannot be said
\vspace{5pt} \par
\begin{tcolorbox}[colback=green!5, colframe=green!40!black, title=Solution]
Answer: a

Explanation:


 
\begin{center}\includegraphics[width=0.5\textwidth]{automata-theory-questions-answers-equivalence-nfa-dfa-q9.png}\end{center}
\end{tcolorbox}
\end{mdframed}
\newpage

\begin{mdframed}[linewidth=1pt, roundcorner=5pt, linecolor=gray!30, backgroundcolor=gray!5]
\textbf{Question 10} 

The total number of states to build the given language using DFA:

L = \{w | w has exactly 2 a's and at least 2 b's\}

a) 10

b) 11

c) 12

d) 13
\vspace{5pt} \par
\begin{tcolorbox}[colback=green!5, colframe=green!40!black, title=Solution]
Answer: a

Explanation: We need to make the number of a as fixed i.e. 2 and b can be 2 or more. Thus, using this condition a finite automata can be created using 1 states.
\end{tcolorbox}
\end{mdframed}
\newpage

\newtopic{Applications of DFA}

\begin{mdframed}[linewidth=1pt, roundcorner=5pt, linecolor=gray!30, backgroundcolor=gray!5]
\textbf{Question 1} 

Given Language: \{x | it is divisible by 3\}

The total number of final states to be assumed in order to pass the number constituting \{0, 1\} is

a) 0

b) 1

c) 2

d) 3
\vspace{5pt} \par
\begin{tcolorbox}[colback=green!5, colframe=green!40!black, title=Solution]
Answer: c

Explanation: The DFA for the given language can be constructed as follows:


 
\begin{center}\includegraphics[width=0.5\textwidth]{automata-theory-questions-answers-aplications-dfa-q1.png}\end{center}
\end{tcolorbox}
\end{mdframed}
\newpage

\begin{mdframed}[linewidth=1pt, roundcorner=5pt, linecolor=gray!30, backgroundcolor=gray!5]
\textbf{Question 2} 

A binary string is divisible by 4 if and only if it ends with:

a) 100

b) 1000

c) 1100

d) 0011
\vspace{5pt} \par
\begin{tcolorbox}[colback=green!5, colframe=green!40!black, title=Solution]
Answer: a

Explanation: If the string is divisible by four, it surely ends with the substring ‘100' while a binary string divisible by 2 would surely end with the substring ‘10'.
\end{tcolorbox}
\end{mdframed}
\newpage

\begin{mdframed}[linewidth=1pt, roundcorner=5pt, linecolor=gray!30, backgroundcolor=gray!5]
\textbf{Question 3} 

Let L be a language whose FA consist of 5 acceptance states and 11 non final states. It further consists of a dumping state. Predict the number of acceptance states in L
c
.

a) 16

b) 11

c) 5

d) 6
\vspace{5pt} \par
\begin{tcolorbox}[colback=green!5, colframe=green!40!black, title=Solution]
Answer: a

Explanation: If L leads to FA1, then for L
c
, the FA can be obtained by exchanging the final and non-final states.
\end{tcolorbox}
\end{mdframed}
\newpage

\begin{mdframed}[linewidth=1pt, roundcorner=5pt, linecolor=gray!30, backgroundcolor=gray!5]
\textbf{Question 4} 

If L1 and L2 are regular languages, which among the following is an exception?

a) L1 U L2

b) L1 – L2

c) L1 $\cap$ L2

d) All of the mentioned
\vspace{5pt} \par
\begin{tcolorbox}[colback=green!5, colframe=green!40!black, title=Solution]
Answer: d

Explanation: It the closure property of Regular language which lays down the following statement:

If L1, L2 are 2- regular languages, then L1 U L2, L1 $\cap$ L2, L1
C
, L1 – L2 are regular language.
\end{tcolorbox}
\end{mdframed}
\newpage

\begin{mdframed}[linewidth=1pt, roundcorner=5pt, linecolor=gray!30, backgroundcolor=gray!5]
\textbf{Question 5} 

Predict the analogous operation for the given language:

A: \{[p, q] | p $\epsilon$ A1, q does not belong to A2\}

a) A1-A2

b) A2-A1

c) A1.A2

d) A1+A2
\vspace{5pt} \par
\begin{tcolorbox}[colback=green!5, colframe=green!40!black, title=Solution]
Answer: a

Explanation: When set operation ‘-‘ is performed between two sets, it points to those values of prior set which belongs to it but not to the latter set analogous to basic subtraction operation.
\end{tcolorbox}
\end{mdframed}
\newpage

\begin{mdframed}[linewidth=1pt, roundcorner=5pt, linecolor=gray!30, backgroundcolor=gray!5]
\textbf{Question 6} 

Which among the following NFA's is correct corresponding to the given Language?

L= \{x$\epsilon$ \{0, 1\} | 3rd bit from right is 0\}

a) 
 
\begin{center}\includegraphics[width=0.5\textwidth]{automata-theory-questions-answers-aplications-dfa-q6a.png}\end{center}
 

b) 
 
\begin{center}\includegraphics[width=0.5\textwidth]{automata-theory-questions-answers-aplications-dfa-q6b.png}\end{center}
 

c) 
 
\begin{center}\includegraphics[width=0.5\textwidth]{automata-theory-questions-answers-aplications-dfa-q6c.png}\end{center}
 

d) None of the mentioned
\vspace{5pt} \par
\begin{tcolorbox}[colback=green!5, colframe=green!40!black, title=Solution]
Answer: a

Explanation: The NFA accepts all binary strings such that the third bit from right end is 1 and if not, is send to Dumping state. Note: It is assumed that the input is given from the right end bit by bit.
\end{tcolorbox}
\end{mdframed}
\newpage

\begin{mdframed}[linewidth=1pt, roundcorner=5pt, linecolor=gray!30, backgroundcolor=gray!5]
\textbf{Question 7} 

Statement 1: NFA computes the string along parallel paths.

Statement 2: An input can be accepted at more than one place in an NFA.

Which among the following options are most appropriate?

a) Statement 1 is true while 2 is not

b) Statement 1 is false while is not

c) Statement 1 and 2, both are true

d) Statement 1 and 2, both are false
\vspace{5pt} \par
\begin{tcolorbox}[colback=green!5, colframe=green!40!black, title=Solution]
Answer: c

Explanation: While the machine runs on some input string, if it has the choice to split, it goes in all possible way and each one is different copy of the machine. The machine takes subsequent choice to split further giving rise to more copies of the machine getting each copy run parallel. If any one copy of the machine accepts the strings, then NFA accepts, otherwise it rejects.
\end{tcolorbox}
\end{mdframed}
\newpage

\begin{mdframed}[linewidth=1pt, roundcorner=5pt, linecolor=gray!30, backgroundcolor=gray!5]
\textbf{Question 8} 

Which of the following options is correct for the given statement?

Statement: If K is the number of states in NFA, the DFA simulating the same language would have states less than 2
k
.

a) True

b) False
\vspace{5pt} \par
\begin{tcolorbox}[colback=green!5, colframe=green!40!black, title=Solution]
Answer: a

Explanation: If K is the number of states in NFA, the DFA simulating the same language would have states equal to or less than 2
k
.
\end{tcolorbox}
\end{mdframed}
\newpage

\begin{mdframed}[linewidth=1pt, roundcorner=5pt, linecolor=gray!30, backgroundcolor=gray!5]
\textbf{Question 9} 

Let N (Q, $\Sigma$, $\delta$, q0, A) be the NFA recognizing a language L. Then for a DFA (Q', $\Sigma$, $\delta$', q0', A'), which among the following is true?

a) Q' = P(Q)

b) $\Delta$' = $\delta$' (R, a) = \{q $\epsilon$ Q | q $\epsilon$ $\delta$ (r, a), for some r $\epsilon$ R\}

c) Q' = \{q0\}

d) All of the mentioned
\vspace{5pt} \par
\begin{tcolorbox}[colback=green!5, colframe=green!40!black, title=Solution]
Answer: d

Explanation: All the optioned mentioned are the instruction formats of how to convert a NFA to a DFA.
\end{tcolorbox}
\end{mdframed}
\newpage

\begin{mdframed}[linewidth=1pt, roundcorner=5pt, linecolor=gray!30, backgroundcolor=gray!5]
\textbf{Question 10} 

There exists an initial state, 17 transition states, 7 final states and one dumping state, Predict the maximum number of states in its equivalent DFA?

a) 226

b) 224

c) 225

d) 223
\vspace{5pt} \par
\begin{tcolorbox}[colback=green!5, colframe=green!40!black, title=Solution]
Answer: a

Explanation: The maximum number of states an equivalent DFA can comprise for its respective NFA with k states will be 2
k
.
\end{tcolorbox}
\end{mdframed}
\newpage

\newtopic{Applications of NFA}

\begin{mdframed}[linewidth=1pt, roundcorner=5pt, linecolor=gray!30, backgroundcolor=gray!5]
\textbf{Question 1} 

Under which of the following operation, NFA is not closed?

a) Negation

b) Kleene

c) Concatenation

d) None of the mentioned
\vspace{5pt} \par
\begin{tcolorbox}[colback=green!5, colframe=green!40!black, title=Solution]
Answer: d

Explanation: NFA is said to be closed under the following operations:

a) Union

b) Intersection

c) Concatenation

d) Kleene

e) Negation.
\end{tcolorbox}
\end{mdframed}
\newpage

\begin{mdframed}[linewidth=1pt, roundcorner=5pt, linecolor=gray!30, backgroundcolor=gray!5]
\textbf{Question 2} 

It is less complex to prove the closure properties over regular languages using:

a) NFA

b) DFA

c) PDA

d) Can't be said
\vspace{5pt} \par
\begin{tcolorbox}[colback=green!5, colframe=green!40!black, title=Solution]
Answer: a

Explanation: None.
\end{tcolorbox}
\end{mdframed}
\newpage

\begin{mdframed}[linewidth=1pt, roundcorner=5pt, linecolor=gray!30, backgroundcolor=gray!5]
\textbf{Question 3} 

Which of the following is an application of Finite Automaton?

a) Compiler Design

b) Grammar Parsers

c) Text Search

d) All of the mentioned
\vspace{5pt} \par
\begin{tcolorbox}[colback=green!5, colframe=green!40!black, title=Solution]
Answer: d

Explanation: There are many applications of finite automata, mainly in the field of Compiler Design and Parsers and Search Engines.
\end{tcolorbox}
\end{mdframed}
\newpage

\begin{mdframed}[linewidth=1pt, roundcorner=5pt, linecolor=gray!30, backgroundcolor=gray!5]
\textbf{Question 4} 

John is asked to make an automaton which accepts a given string for all the occurrence of ‘1001' in it. How many number of transitions would John use such that, the string processing application works?

a) 9

b) 11

c) 12

d) 15
\vspace{5pt} \par
\begin{tcolorbox}[colback=green!5, colframe=green!40!black, title=Solution]
Answer: a

Explanation: None.
\end{tcolorbox}
\end{mdframed}
\newpage

\begin{mdframed}[linewidth=1pt, roundcorner=5pt, linecolor=gray!30, backgroundcolor=gray!5]
\textbf{Question 5} 

Which of the following do we use to form an NFA from a regular expression?

a) Subset Construction Method

b) Power Set Construction Method

c) Thompson Construction Method

d) Scott Construction Method
\vspace{5pt} \par
\begin{tcolorbox}[colback=green!5, colframe=green!40!black, title=Solution]
Answer: c

Explanation: Thompson Construction method is used to turn a regular expression in an NFA by fragmenting the given regular expression through the operations performed on the input alphabets.





     
Free Certifications
 on 300 subjects are live for April 2026. Register Now!
\end{tcolorbox}
\end{mdframed}
\newpage

\begin{mdframed}[linewidth=1pt, roundcorner=5pt, linecolor=gray!30, backgroundcolor=gray!5]
\textbf{Question 6} 

Which among the following can be an example of application of finite state machine(FSM)?

a) Communication Link

b) Adder

c) Stack

d) None of the mentioned
\vspace{5pt} \par
\begin{tcolorbox}[colback=green!5, colframe=green!40!black, title=Solution]
Answer: a

Explanation: Idle is the state when data in form of packets is send and returns if NAK is received else waits for the NAK to be received.
\end{tcolorbox}
\end{mdframed}
\newpage

\begin{mdframed}[linewidth=1pt, roundcorner=5pt, linecolor=gray!30, backgroundcolor=gray!5]
\textbf{Question 7} 

Which among the following is not an application of FSM?

a) Lexical Analyser

b) BOT

c) State charts

d) None of the mentioned
\vspace{5pt} \par
\begin{tcolorbox}[colback=green!5, colframe=green!40!black, title=Solution]
Answer: d

Explanation: Finite state automation is used in Lexical Analyser, Computer BOT (used in games), State charts, etc.
\end{tcolorbox}
\end{mdframed}
\newpage

\begin{mdframed}[linewidth=1pt, roundcorner=5pt, linecolor=gray!30, backgroundcolor=gray!5]
\textbf{Question 8} 

L1= \{w | w does not contain the string tr \}

L2= \{w | w does contain the string tr\}

Given $\Sigma$= \{t, r\}, The difference of the minimum number of states required to form L1 and L2?

a) 0

b) 1

c) 2

d) Cannot be said
\vspace{5pt} \par
\begin{tcolorbox}[colback=green!5, colframe=green!40!black, title=Solution]
Answer: a

Explanation: None.
\end{tcolorbox}
\end{mdframed}
\newpage

\begin{mdframed}[linewidth=1pt, roundcorner=5pt, linecolor=gray!30, backgroundcolor=gray!5]
\textbf{Question 9} 

Predict the number of transitions required to automate the following language using only 3 states:

L= \{w | w ends with 00\}

a) 3

b) 2

c) 4

d) Cannot be said
\vspace{5pt} \par
\begin{tcolorbox}[colback=green!5, colframe=green!40!black, title=Solution]
Answer: a

Explanation: None.
\end{tcolorbox}
\end{mdframed}
\newpage

\begin{mdframed}[linewidth=1pt, roundcorner=5pt, linecolor=gray!30, backgroundcolor=gray!5]
\textbf{Question 10} 

The total number of states to build the given language using DFA:

L= \{w | w has exactly 2 a's and at least 2 b's\}

a) 10

b) 11

c) 12

d) 13
\vspace{5pt} \par
\begin{tcolorbox}[colback=green!5, colframe=green!40!black, title=Solution]
Answer: a

Explanation: We need to make the number of a as fixed i.e. 2 and b can be 2 or more. Thus, using this condition a finite automata can be created using 1 states.
\end{tcolorbox}
\end{mdframed}
\newpage

\newtopic{Finite Automata with Epsilon Transition}

\begin{mdframed}[linewidth=1pt, roundcorner=5pt, linecolor=gray!30, backgroundcolor=gray!5]
\textbf{Question 1} 

According to the given transitions, which among the following are the epsilon closures of q1 for the given NFA?

$\Delta$ (q1, $\epsilon$) = \{q2, q3, q4\}

$\Delta$ (q4, 1) = q1

$\Delta$ (q1, $\epsilon$) = q1

a) q4

b) q2

c) q1

d) q1, q2, q3, q4
\vspace{5pt} \par
\begin{tcolorbox}[colback=green!5, colframe=green!40!black, title=Solution]
Answer: d

Explanation: The set of states which can be reached from q using $\epsilon$-transitions, is called the $\epsilon$-closure over state q.
\end{tcolorbox}
\end{mdframed}
\newpage

\begin{mdframed}[linewidth=1pt, roundcorner=5pt, linecolor=gray!30, backgroundcolor=gray!5]
\textbf{Question 2} 

State true or false?

Statement: An NFA can be modified to allow transition without input alphabets, along with one or more transitions on input symbols.

a) True

b) False
\vspace{5pt} \par
\begin{tcolorbox}[colback=green!5, colframe=green!40!black, title=Solution]
Answer: a

Explanation: It is possible to construct an NFA with $\epsilon$-transitions, presence of no input symbols, and that is called NFA with $\epsilon$-moves.
\end{tcolorbox}
\end{mdframed}
\newpage

\begin{mdframed}[linewidth=1pt, roundcorner=5pt, linecolor=gray!30, backgroundcolor=gray!5]
\textbf{Question 3} 

State true or false?

Statement: $\epsilon$ (Input) does not appears on Input tape.

a) True

b) False
\vspace{5pt} \par
\begin{tcolorbox}[colback=green!5, colframe=green!40!black, title=Solution]
Answer: a

Explanation: $\epsilon$ does not appears on Input tape, $\epsilon$ transition means a transition without scanning a symbol i.e. without moving the read head.
\end{tcolorbox}
\end{mdframed}
\newpage

\begin{mdframed}[linewidth=1pt, roundcorner=5pt, linecolor=gray!30, backgroundcolor=gray!5]
\textbf{Question 4} 

Statement 1: $\epsilon$- transition can be called as hidden non-determinism.

Statement 2: $\delta$ (q, $\epsilon$) = p means from q it can jump to p with a shift in read head.

Which among the following options is correct?

a) Statement 1 and 2, both are correct

b) Statement 1 and 2, both are wrong

c) Statement 1 is correct while Statement 2 is wrong

d) Statement 1 is wrong while Statement 2 is correct
\vspace{5pt} \par
\begin{tcolorbox}[colback=green!5, colframe=green!40!black, title=Solution]
Answer: c

Explanation: The transition with $\epsilon$ leads to a jump but without any shift in read head. Further, the method can be called one to introduce hidden non-determinism.
\end{tcolorbox}
\end{mdframed}
\newpage

\begin{mdframed}[linewidth=1pt, roundcorner=5pt, linecolor=gray!30, backgroundcolor=gray!5]
\textbf{Question 5} 

$\epsilon$- closure of q1 in the given transition graph:

a) \{q1\}

b) \{q0, q2\}

c) \{q1, q2\}

d) \{q0, q1, q2\}
\vspace{5pt} \par
\begin{tcolorbox}[colback=green!5, colframe=green!40!black, title=Solution]
Answer: c

Explanation: $\epsilon$-closure is defined as the set of states being reached through $\epsilon$-transitions from a starting state.
\end{tcolorbox}
\end{mdframed}
\newpage

\begin{mdframed}[linewidth=1pt, roundcorner=5pt, linecolor=gray!30, backgroundcolor=gray!5]
\textbf{Question 6} 

Predict the total number of final states after removing the $\epsilon$-moves from the given NFA?

a) 1

b) 2

c) 3

d) 0
\vspace{5pt} \par
\begin{tcolorbox}[colback=green!5, colframe=green!40!black, title=Solution]
Answer: c

Explanation: The NFA which would result after eliminating $\epsilon$-moves can be shown diagramatically.
\end{tcolorbox}
\end{mdframed}
\newpage

\begin{mdframed}[linewidth=1pt, roundcorner=5pt, linecolor=gray!30, backgroundcolor=gray!5]
\textbf{Question 7} 

For NFA with $\epsilon$-moves, which among the following is correct?

a) $\Delta$: Q X ($\Sigma$ U \{$\epsilon$\}) -$>$ P(Q)

b) $\Delta$: Q X ($\Sigma$) -$>$ P(Q)

c) $\Delta$: Q X ($\Sigma$*) -$>$ P(Q)

d) All of the mentioned
\vspace{5pt} \par
\begin{tcolorbox}[colback=green!5, colframe=green!40!black, title=Solution]
Answer: a

Explanation: Due to the presence of $\epsilon$ symbol, or rather an epsilon-move, the input alphabets unites with it to form a set including $\epsilon$.
\end{tcolorbox}
\end{mdframed}
\newpage

\begin{mdframed}[linewidth=1pt, roundcorner=5pt, linecolor=gray!30, backgroundcolor=gray!5]
\textbf{Question 8} 

Which among the following is false?

$\epsilon$-closure of a subset S of Q is:

a) Every element of S $\epsilon$ Q

b) For any q $\epsilon$ $\epsilon$(S), every element of $\delta$ (q, $\epsilon$) is in $\epsilon$(S)

c) No other element is in $\epsilon$(S)

d) None of the mentioned
\vspace{5pt} \par
\begin{tcolorbox}[colback=green!5, colframe=green!40!black, title=Solution]
Answer: d

Explanation: All the mentioned are the closure properties of $\epsilon$ and encircles all the elements if it satisfies the following options:

a) Every element of S $\epsilon$ Q

b) For any q $\epsilon$ $\epsilon$(S), every element of $\delta$ (q, $\epsilon$) is in $\epsilon$(S)

c) No other element is in $\epsilon$(S)
\end{tcolorbox}
\end{mdframed}
\newpage

\begin{mdframed}[linewidth=1pt, roundcorner=5pt, linecolor=gray!30, backgroundcolor=gray!5]
\textbf{Question 9} 

The automaton which allows transformation to a new state without consuming any input symbols:

a) NFA

b) DFA

c) NFA-l

d) All of the mentioned
\vspace{5pt} \par
\begin{tcolorbox}[colback=green!5, colframe=green!40!black, title=Solution]
Answer: c

Explanation: NFA-l or e-NFA is an extension of Non deterministic Finite Automata which are usually called NFA with epsilon moves or lambda transitions.
\end{tcolorbox}
\end{mdframed}
\newpage

\begin{mdframed}[linewidth=1pt, roundcorner=5pt, linecolor=gray!30, backgroundcolor=gray!5]
\textbf{Question 10} 

e-transitions are

a) conditional

b) unconditional

c) input dependent

d) none of the mentioned
\vspace{5pt} \par
\begin{tcolorbox}[colback=green!5, colframe=green!40!black, title=Solution]
Answer: b

Explanation: An epsilon move is a transition from one state to another that doesnt require any specific condition.
\end{tcolorbox}
\end{mdframed}
\newpage

\begin{mdframed}[linewidth=1pt, roundcorner=5pt, linecolor=gray!30, backgroundcolor=gray!5]
\textbf{Question 11} 

The \_\_\_\_\_\_\_\_\_\_ of a set of states, P, of an NFA is defined as the set of states reachable from any state in P following e-transitions.

a) e-closure

b) e-pack

c) Q in the tuple

d) None of the mentioned
\vspace{5pt} \par
\begin{tcolorbox}[colback=green!5, colframe=green!40!black, title=Solution]
Answer: a

Explanation: The e-closure of a set of states, P, of an NFA is defined as the set of states reachable from any state in P following e-transitions.
\end{tcolorbox}
\end{mdframed}
\newpage

\begin{mdframed}[linewidth=1pt, roundcorner=5pt, linecolor=gray!30, backgroundcolor=gray!5]
\textbf{Question 12} 

The e-NFA  recognizable languages are not closed under :

a) Union

b) Negation

c) Kleene Closure

d) None of the mentioned
\vspace{5pt} \par
\begin{tcolorbox}[colback=green!5, colframe=green!40!black, title=Solution]
Answer: d

Explanation: The languages which are recognized by an epsilon Non deterministic automata are closed under the following operations:

a) Union

b) Intersection

c) Concatenation

d) Negation

e) Star

f) Kleene closure
\end{tcolorbox}
\end{mdframed}
\newpage

\newtopic{Uses of Epsilon-Transitions}

\begin{mdframed}[linewidth=1pt, roundcorner=5pt, linecolor=gray!30, backgroundcolor=gray!5]
\textbf{Question 1} 

The automaton which allows transformation to a new state without consuming any input symbols:

a) NFA

b) DFA

c) NFA-l

d) All of the mentioned
\vspace{5pt} \par
\begin{tcolorbox}[colback=green!5, colframe=green!40!black, title=Solution]
Answer: c

Explanation: NFA-l or e-NFA is an extension of Non deterministic Finite Automata which are usually called NFA with epsilon moves or lambda transitions.
\end{tcolorbox}
\end{mdframed}
\newpage

\begin{mdframed}[linewidth=1pt, roundcorner=5pt, linecolor=gray!30, backgroundcolor=gray!5]
\textbf{Question 2} 

e-transitions are

a) conditional

b) unconditional

c) input dependent

d) none of the mentioned
\vspace{5pt} \par
\begin{tcolorbox}[colback=green!5, colframe=green!40!black, title=Solution]
Answer: b

Explanation: An epsilon move is a transition from one state to another that doesn't require any specific condition.
\end{tcolorbox}
\end{mdframed}
\newpage

\begin{mdframed}[linewidth=1pt, roundcorner=5pt, linecolor=gray!30, backgroundcolor=gray!5]
\textbf{Question 3} 

The \_\_\_\_\_\_\_\_\_\_ of a set of states, P, of an NFA is defined as the set of states reachable from any state in P following e-transitions.

a) e-closure

b) e-pack

c) Q in the tuple

d) None of the mentioned
\vspace{5pt} \par
\begin{tcolorbox}[colback=green!5, colframe=green!40!black, title=Solution]
Answer: a

Explanation: The e-closure of a set of states, P, of an NFAis defined as the set of states reachable from any state in P following e-transitions.
\end{tcolorbox}
\end{mdframed}
\newpage

\begin{mdframed}[linewidth=1pt, roundcorner=5pt, linecolor=gray!30, backgroundcolor=gray!5]
\textbf{Question 4} 

The e-NFA  recognizable languages are not closed under \_\_\_\_\_\_\_\_\_\_\_

a) Union

b) Negation

c) Kleene Closure

d) None of the mentioned
\vspace{5pt} \par
\begin{tcolorbox}[colback=green!5, colframe=green!40!black, title=Solution]
Answer: d

Explanation: The languages which are recognized by an epsilon Non deterministic automata are closed under the following operations:

i) Union

ii) Intersection

iii) Concatenation

iv) Negation

v) Star

vi) Kleene closure
\end{tcolorbox}
\end{mdframed}
\newpage

\begin{mdframed}[linewidth=1pt, roundcorner=5pt, linecolor=gray!30, backgroundcolor=gray!5]
\textbf{Question 5} 

Is the language preserved in all the steps while eliminating epsilon transitions from a NFA?

a) yes

b) no
\vspace{5pt} \par
\begin{tcolorbox}[colback=green!5, colframe=green!40!black, title=Solution]
Answer: a

Explanation: Yes, the language is preserved during the dteps of construction: L(N)=L(N1)=L(N2)=L(3).





     
Free Certifications
 on 300 subjects are live for April 2026. Register Now!
\end{tcolorbox}
\end{mdframed}
\newpage

\begin{mdframed}[linewidth=1pt, roundcorner=5pt, linecolor=gray!30, backgroundcolor=gray!5]
\textbf{Question 6} 

An e-NFA is \_\_\_\_\_\_\_\_\_\_\_ in representation.

a) Quadruple

b) Quintuple

c) Triple

d) None of the mentioned
\vspace{5pt} \par
\begin{tcolorbox}[colback=green!5, colframe=green!40!black, title=Solution]
Answer: b

Explanation: An e-NFA consist of 5 tuples: A=(Q, S, d, q0, F)

Note: e is never a member of S.
\end{tcolorbox}
\end{mdframed}
\newpage

\begin{mdframed}[linewidth=1pt, roundcorner=5pt, linecolor=gray!30, backgroundcolor=gray!5]
\textbf{Question 7} 

State true or false:

Statement: Both NFA and e-NFA recognize exactly the same languages.

a) true

b) false
\vspace{5pt} \par
\begin{tcolorbox}[colback=green!5, colframe=green!40!black, title=Solution]
Answer: a

Explanation: e-NFA do come up with a convenient feature but nothing new.They do not extend the class of languages that can be represented.
\end{tcolorbox}
\end{mdframed}
\newpage

\newtopic{Epsilon Closures}

\begin{mdframed}[linewidth=1pt, roundcorner=5pt, linecolor=gray!30, backgroundcolor=gray!5]
\textbf{Question 1} 

Which of the following does not belong to input alphabet if S=\{a, b\}* for any language?

a) a

b) b

c) e

d) none of the mentioned
\vspace{5pt} \par
\begin{tcolorbox}[colback=green!5, colframe=green!40!black, title=Solution]
Answer: c

Explanation: The automaton may be allowed to change its state without reading the input symbol using epsilon but this does not mean that epsilon has become an input symbol. On the contrary, one assumes that the symbol epsilon does not belong to any alphabet.
\end{tcolorbox}
\end{mdframed}
\newpage

\begin{mdframed}[linewidth=1pt, roundcorner=5pt, linecolor=gray!30, backgroundcolor=gray!5]
\textbf{Question 2} 

The number of final states we need as per the given language?

Language L: \{an| n is even or divisible by 3\}

a) 1

b) 2

c) 3

d) 4
\vspace{5pt} \par
\begin{tcolorbox}[colback=green!5, colframe=green!40!black, title=Solution]
Answer: b

Explanation:


 
\begin{center}\includegraphics[width=0.5\textwidth]{automata-theory-questions-answers-epsilon-closures-q2.png}\end{center}
\end{tcolorbox}
\end{mdframed}
\newpage

\begin{mdframed}[linewidth=1pt, roundcorner=5pt, linecolor=gray!30, backgroundcolor=gray!5]
\textbf{Question 3} 

State true or false:

Statement: Both NFA and e-NFA recognize exactly the same languages.

a) true

b) false
\vspace{5pt} \par
\begin{tcolorbox}[colback=green!5, colframe=green!40!black, title=Solution]
Answer: a

Explanation: e-NFA do come up with a convenient feature but nothing new.They do not extend the class of languages that can be represented.
\end{tcolorbox}
\end{mdframed}
\newpage

\begin{mdframed}[linewidth=1pt, roundcorner=5pt, linecolor=gray!30, backgroundcolor=gray!5]
\textbf{Question 4} 

Design a NFA for the language:

L: \{an| n is even or divisible by 3\}

Which of the following methods can be used to simulate the same.

a) e-NFA

b) Power Construction Method

c) e-NFA and Power Construction Method

d) None of the mentioned
\vspace{5pt} \par
\begin{tcolorbox}[colback=green!5, colframe=green!40!black, title=Solution]
Answer: c

Explanation: It is more convenient to simulate a machine using e-NFA else the method of Power Construction is used from the union-closure of DFA's.
\end{tcolorbox}
\end{mdframed}
\newpage

\begin{mdframed}[linewidth=1pt, roundcorner=5pt, linecolor=gray!30, backgroundcolor=gray!5]
\textbf{Question 5} 

Which of the following belongs to the epsilon closure set of a?


 
\begin{center}\includegraphics[width=0.5\textwidth]{automata-theory-questions-answers-epsilon-closures-q5.png}\end{center}
 

a) \{f1, f2, f3\}

b) \{a, f1, f2, f3\}

c) \{f1, f2\}

d) none of the mentioned
\vspace{5pt} \par
\begin{tcolorbox}[colback=green!5, colframe=green!40!black, title=Solution]
Answer: b

Explanation: The epsilon closure of the set q is the set that contains q, together with all the states which can be reached starting at q by following only epsilon transitions.





     Apply Now for 
Free C Certification
 - April 2026
\end{tcolorbox}
\end{mdframed}
\newpage

\begin{mdframed}[linewidth=1pt, roundcorner=5pt, linecolor=gray!30, backgroundcolor=gray!5]
\textbf{Question 6} 

The number of elements present in the e-closure(f2) in the given diagram:


 
\begin{center}\includegraphics[width=0.5\textwidth]{automata-theory-questions-answers-epsilon-closures-q5.png}\end{center}
 

a) 0

b) 1

c) 2

d) 3
\vspace{5pt} \par
\begin{tcolorbox}[colback=green!5, colframe=green!40!black, title=Solution]
Answer: c

Explanation: The epsilon closure set of f2 consist of the elements:\{f2, f3\}. Thus the count of the element in the closure set is 2.
\end{tcolorbox}
\end{mdframed}
\newpage

\begin{mdframed}[linewidth=1pt, roundcorner=5pt, linecolor=gray!30, backgroundcolor=gray!5]
\textbf{Question 7} 

Which of the steps are non useful while eliminating the e-transitions for the given diagram?


 
\begin{center}\includegraphics[width=0.5\textwidth]{automata-theory-questions-answers-epsilon-closures-q5.png}\end{center}
 

a) Make a as accepting state of N' if ECLOSE(p) contains an accepting state of N

b) Add an arc a to f1 labelled a if there is an arc labelled a in N from some state in ECLOSE(a) to f1

c) Delete all arcs labelled as e

d) None of the mentioned
\vspace{5pt} \par
\begin{tcolorbox}[colback=green!5, colframe=green!40!black, title=Solution]
Answer: d

Explanation: The given are the steps followed while eliminating epsilon transitions from a NFA or converting an e-NFA to just NFA.
\end{tcolorbox}
\end{mdframed}
\newpage

\begin{mdframed}[linewidth=1pt, roundcorner=5pt, linecolor=gray!30, backgroundcolor=gray!5]
\textbf{Question 8} 

Remove all the epsilon transitions in the given diagram and compute the number of a-transitions in the result?


 
\begin{center}\includegraphics[width=0.5\textwidth]{automata-theory-questions-answers-epsilon-closures-q2.png}\end{center}
 

a) 5

b) 7

c) 9

d) 6
\vspace{5pt} \par
\begin{tcolorbox}[colback=green!5, colframe=green!40!black, title=Solution]
Answer: b

Explanation:


 
\begin{center}\includegraphics[width=0.5\textwidth]{automata-theory-questions-answers-epsilon-closures-q8.png}\end{center}
\end{tcolorbox}
\end{mdframed}
\newpage

\newtopic{Union, Intersection \& Complement}

\begin{mdframed}[linewidth=1pt, roundcorner=5pt, linecolor=gray!30, backgroundcolor=gray!5]
\textbf{Question 1} 

Regular sets are closed under union,concatenation and kleene closure.

a) True

b) False

c) Depends on regular set

d) Can't say
\vspace{5pt} \par
\begin{tcolorbox}[colback=green!5, colframe=green!40!black, title=Solution]
Answer:a

Explanation: Regular sets are closed under these three operation.
\end{tcolorbox}
\end{mdframed}
\newpage

\begin{mdframed}[linewidth=1pt, roundcorner=5pt, linecolor=gray!30, backgroundcolor=gray!5]
\textbf{Question 2} 

Complement of a DFA can be obtained by

a) making starting state as final state

b) no trival method

c) making final states non-final and non-final to final

d) make final as a starting state
\vspace{5pt} \par
\begin{tcolorbox}[colback=green!5, colframe=green!40!black, title=Solution]
Answer:c

Explanation: String accepted in previous DFA will not be accepted and non accepting string will be accepted .
\end{tcolorbox}
\end{mdframed}
\newpage

\begin{mdframed}[linewidth=1pt, roundcorner=5pt, linecolor=gray!30, backgroundcolor=gray!5]
\textbf{Question 3} 

Complement of regular sets are \_\_\_\_\_\_\_\_\_

a) Regular

b) CFG

c) CSG

d) RE
\vspace{5pt} \par
\begin{tcolorbox}[colback=green!5, colframe=green!40!black, title=Solution]
Answer:a

Explanation: Regular sets are closed under complement operation.
\end{tcolorbox}
\end{mdframed}
\newpage

\begin{mdframed}[linewidth=1pt, roundcorner=5pt, linecolor=gray!30, backgroundcolor=gray!5]
\textbf{Question 4} 

If L1 and L2 are regular sets then intersection of these two will be

a) Regular

b) Non Regular

c) Recursive

d) Non Recursive
\vspace{5pt} \par
\begin{tcolorbox}[colback=green!5, colframe=green!40!black, title=Solution]
Answer:a

Explanation: Regular expression are also colsed under intersection.
\end{tcolorbox}
\end{mdframed}
\newpage

\begin{mdframed}[linewidth=1pt, roundcorner=5pt, linecolor=gray!30, backgroundcolor=gray!5]
\textbf{Question 5} 

If L1 is regular L2 is unknown but L1-L2 is regular ,then L2 must be

a) Empty set

b) CFG

c) Decidable

d) Regular
\vspace{5pt} \par
\begin{tcolorbox}[colback=green!5, colframe=green!40!black, title=Solution]
Answer:d

Explanation: Regular is closed under difference.







     
Free Certifications
 on 300 subjects are live for April 2026. Register Now!
\end{tcolorbox}
\end{mdframed}
\newpage

\begin{mdframed}[linewidth=1pt, roundcorner=5pt, linecolor=gray!30, backgroundcolor=gray!5]
\textbf{Question 6} 

Reverse of a DFA can be formed by

a) using PDA

b) making final state as non-final

c) making final as starting state and starting state as final state

d) None of the mentioned
\vspace{5pt} \par
\begin{tcolorbox}[colback=green!5, colframe=green!40!black, title=Solution]
Answer:c

Explanation: By making final state as starting state string starting from end will be accepted.
\end{tcolorbox}
\end{mdframed}
\newpage

\begin{mdframed}[linewidth=1pt, roundcorner=5pt, linecolor=gray!30, backgroundcolor=gray!5]
\textbf{Question 7} 

Reverse of (0+1)* will be

a) Phi

b) Null

c) (0+1)*

d) (0+1)
\vspace{5pt} \par
\begin{tcolorbox}[colback=green!5, colframe=green!40!black, title=Solution]
Answer:c

Explanation: There is only one state which is start and final state of DFA so interchanging starting start and final state doesn't change DFA.
\end{tcolorbox}
\end{mdframed}
\newpage

\begin{mdframed}[linewidth=1pt, roundcorner=5pt, linecolor=gray!30, backgroundcolor=gray!5]
\textbf{Question 8} 

A \_\_\_\_\_\_\_\_\_\_\_ is a substitution such that h(a) contains a string for each a.

a) Closure

b) Interchange

c) Homomorphism

d) Inverse Homomorphism
\vspace{5pt} \par
\begin{tcolorbox}[colback=green!5, colframe=green!40!black, title=Solution]
Answer:c

Explanation: This operation replace using a function .
\end{tcolorbox}
\end{mdframed}
\newpage

\begin{mdframed}[linewidth=1pt, roundcorner=5pt, linecolor=gray!30, backgroundcolor=gray!5]
\textbf{Question 9} 

Homomorphism of a regular set is \_\_\_\_\_\_\_

a) Universal set

b) Null set

c) Regular set

d) Non regular set
\vspace{5pt} \par
\begin{tcolorbox}[colback=green!5, colframe=green!40!black, title=Solution]
Answer:c

Explanation: Regular set are closed under homomorphism.
\end{tcolorbox}
\end{mdframed}
\newpage

\begin{mdframed}[linewidth=1pt, roundcorner=5pt, linecolor=gray!30, backgroundcolor=gray!5]
\textbf{Question 10} 

(a \textasciicircum{} 5b \textasciicircum{} 5)* is example of \_\_\_\_\_\_\_\_

a) Type 0 language

b) Type 1 language

c) Type 2 language

d) Type 3 language
\vspace{5pt} \par
\begin{tcolorbox}[colback=green!5, colframe=green!40!black, title=Solution]
Answer:d

Explanation: It is a regular expression.
\end{tcolorbox}
\end{mdframed}
\newpage

\begin{mdframed}[linewidth=1pt, roundcorner=5pt, linecolor=gray!30, backgroundcolor=gray!5]
\textbf{Question 11} 

Which of the following is type 3 language ?

a) Strings of 0's whose length is perfect square

b) Palindromes string

c) Strings of 0's having length prime number

d) String of odd number of 0's
\vspace{5pt} \par
\begin{tcolorbox}[colback=green!5, colframe=green!40!black, title=Solution]
Answer:d

Explanation: Only d is regular language.
\end{tcolorbox}
\end{mdframed}
\newpage

\begin{mdframed}[linewidth=1pt, roundcorner=5pt, linecolor=gray!30, backgroundcolor=gray!5]
\textbf{Question 12} 

a \textasciicircum{} nb \textasciicircum{} n where (n+m) is even .

a) Type 0

b) Type 1

c) Type 2

d) Type 3
\vspace{5pt} \par
\begin{tcolorbox}[colback=green!5, colframe=green!40!black, title=Solution]
Answer:d

Explanation: It is a regular expression.
\end{tcolorbox}
\end{mdframed}
\newpage

\begin{mdframed}[linewidth=1pt, roundcorner=5pt, linecolor=gray!30, backgroundcolor=gray!5]
\textbf{Question 13} 

Complement of a \textasciicircum{} nb \textasciicircum{} m where n $>$= 4 and m $<$= 3 is example of

a) Type 0

b) Type 1

c) Type 2

d) Type 3
\vspace{5pt} \par
\begin{tcolorbox}[colback=green!5, colframe=green!40!black, title=Solution]
Answer:d

Explanation: It is a regular expression.
\end{tcolorbox}
\end{mdframed}
\newpage

\begin{mdframed}[linewidth=1pt, roundcorner=5pt, linecolor=gray!30, backgroundcolor=gray!5]
\textbf{Question 14} 

a \textasciicircum{} nb \textasciicircum{} m where n $>$= 1, m $>$= 1, nm $>$= 3 is example of

a) Type 0

b) Type 1

c) Type 2

d) Type 3
\vspace{5pt} \par
\begin{tcolorbox}[colback=green!5, colframe=green!40!black, title=Solution]
Answer:d

Explanation: It is a regular expression.
\end{tcolorbox}
\end{mdframed}
\newpage

\begin{mdframed}[linewidth=1pt, roundcorner=5pt, linecolor=gray!30, backgroundcolor=gray!5]
\textbf{Question 15} 

Complement of (a + b)* will be

a) phi

b) null

c) a

d) b
\vspace{5pt} \par
\begin{tcolorbox}[colback=green!5, colframe=green!40!black, title=Solution]
Answer:a

Explanation: Given expression accept all string so complement will accept nothing.
\end{tcolorbox}
\end{mdframed}
\newpage

\newtopic{Regular Expression – Introduction}

\begin{mdframed}[linewidth=1pt, roundcorner=5pt, linecolor=gray!30, backgroundcolor=gray!5]
\textbf{Question 1} 

L is a regular Language if and only If the set of \_\_\_\_\_\_\_\_\_\_ classes of IL is finite.

a) Equivalence

b) Reflexive

c) Myhill

d) Nerode
\vspace{5pt} \par
\begin{tcolorbox}[colback=green!5, colframe=green!40!black, title=Solution]
Answer: a

Explanation: According to Myhill Nerode theorem, the corollary proves the given statement correct for equivalence classes.
\end{tcolorbox}
\end{mdframed}
\newpage

\begin{mdframed}[linewidth=1pt, roundcorner=5pt, linecolor=gray!30, backgroundcolor=gray!5]
\textbf{Question 2} 

A language can be generated from simple primitive language in a simple way if and only if

a) It is recognized by a device of infinite states

b) It takes no auxiliary memory

c) All of the mentioned

d) None of the mentioned
\vspace{5pt} \par
\begin{tcolorbox}[colback=green!5, colframe=green!40!black, title=Solution]
Answer: b

Explanation: A language is regular if and only if it can be accepted by a finite automaton. Secondly, It supports no concept of auxiliary memory as it loses the data as soon as the device is shut down.
\end{tcolorbox}
\end{mdframed}
\newpage

\begin{mdframed}[linewidth=1pt, roundcorner=5pt, linecolor=gray!30, backgroundcolor=gray!5]
\textbf{Question 3} 

Which of the following does not represents the given language?

Language: \{0,01\}

a) 0+01

b) \{0\} U \{01\}

c) \{0\} U \{0\}\{1\}

d) \{0\} \textasciicircum{} \{01\}
\vspace{5pt} \par
\begin{tcolorbox}[colback=green!5, colframe=green!40!black, title=Solution]
Answer: d

Explanation: The given option represents \{0, 01\} in different forms using set operations and Regular Expressions. The operator like \textasciicircum{}, v, etc. are logical operation and they form invalid regular expressions when used.
\end{tcolorbox}
\end{mdframed}
\newpage

\begin{mdframed}[linewidth=1pt, roundcorner=5pt, linecolor=gray!30, backgroundcolor=gray!5]
\textbf{Question 4} 

According to the given language, which among the following expressions does it corresponds to?

Language L=\{x$\epsilon$\{0,1\}|x is of length 4 or less\}

a) (0+1+0+1+0+1+0+1)
4

b) (0+1)
4

c) (01)
4

d) (0+1+$\epsilon$)
4
\vspace{5pt} \par
\begin{tcolorbox}[colback=green!5, colframe=green!40!black, title=Solution]
Answer: d

Explanation: The extended notation would be (0+1)
4
 but however, we may allow some or all the factors to be $\epsilon$. Thus $\epsilon$ needs to be included in the given regular expression.
\end{tcolorbox}
\end{mdframed}
\newpage

\begin{mdframed}[linewidth=1pt, roundcorner=5pt, linecolor=gray!30, backgroundcolor=gray!5]
\textbf{Question 5} 

Which among the following looks similar to the given expression?

       ((0+1). (0+1)) *

a) \{x$\epsilon$ \{0,1\} *|x is all binary number with even length\}

b) \{x$\epsilon$ \{0,1\} |x is all binary number with even length\}

c) \{x$\epsilon$ \{0,1\} *|x is all binary number with odd length\}

d) \{x$\epsilon$ \{0,1\} |x is all binary number with odd length\}
\vspace{5pt} \par
\begin{tcolorbox}[colback=green!5, colframe=green!40!black, title=Solution]
Answer: a

Explanation: The given regular expression corresponds to a language of binary strings which is of even length including a length of 0.
\end{tcolorbox}
\end{mdframed}
\newpage

\begin{mdframed}[linewidth=1pt, roundcorner=5pt, linecolor=gray!30, backgroundcolor=gray!5]
\textbf{Question 6} 

If R represents a regular language, which of the following represents the Venn-diagram most correctly?


 
\begin{center}\includegraphics[width=0.5\textwidth]{automata-theory-questions-answers-regular-expression-introduction-q6.png}\end{center}
 

a) An Irregular Set

b) R*

c) R complement

d) R reverse
\vspace{5pt} \par
\begin{tcolorbox}[colback=green!5, colframe=green!40!black, title=Solution]
Answer: b

Explanation: The given diagram represents the Kleene operation over the Regular Language R in which the final states become the initial and the initial state becomes final.
\end{tcolorbox}
\end{mdframed}
\newpage

\begin{mdframed}[linewidth=1pt, roundcorner=5pt, linecolor=gray!30, backgroundcolor=gray!5]
\textbf{Question 7} 

The given NFA corresponds to which of the following Regular expressions?


 
\begin{center}\includegraphics[width=0.5\textwidth]{automata-theory-questions-answers-regular-expression-introduction-q7.png}\end{center}
 

a) (0+1) *(00+11) (0+1) *

b) (0+1) *(00+11) *(0+1) *

c) (0+1) *(00+11) (0+1)

d) (0+1) (00+11) (0+1) *
\vspace{5pt} \par
\begin{tcolorbox}[colback=green!5, colframe=green!40!black, title=Solution]
Answer: a

Explanation: The transition states shown are the result of breaking down the given regular expression in fragments. For dot operation, we change a state, for union (plus) operation, we diverge into two transitions and for Kleene Operation, we apply a loop.
\end{tcolorbox}
\end{mdframed}
\newpage

\begin{mdframed}[linewidth=1pt, roundcorner=5pt, linecolor=gray!30, backgroundcolor=gray!5]
\textbf{Question 8} 

Concatenation Operation refers to which of the following set operations:

a) Union

b) Dot

c) Kleene

d) Two of the options are correct
\vspace{5pt} \par
\begin{tcolorbox}[colback=green!5, colframe=green!40!black, title=Solution]
Answer: b

Explanation: Two operands are said to be performing Concatenation operation AB = A•B = \{xy: x $\in$ A \& y $\in$ B\}.
\end{tcolorbox}
\end{mdframed}
\newpage

\begin{mdframed}[linewidth=1pt, roundcorner=5pt, linecolor=gray!30, backgroundcolor=gray!5]
\textbf{Question 9} 

Concatenation of R with $\Phi$ outputs:

a) R

b) $\Phi$

c) R.$\Phi$

d) None of the mentioned
\vspace{5pt} \par
\begin{tcolorbox}[colback=green!5, colframe=green!40!black, title=Solution]
Answer: b

Explanation: By distributive property (Regular expression identities), we can prove the given identity to be $\Phi$.
\end{tcolorbox}
\end{mdframed}
\newpage

\begin{mdframed}[linewidth=1pt, roundcorner=5pt, linecolor=gray!30, backgroundcolor=gray!5]
\textbf{Question 10} 

RR* can be expressed in which of the forms:

a) R+

b) R-

c) R+ U R-

d) R
\vspace{5pt} \par
\begin{tcolorbox}[colback=green!5, colframe=green!40!black, title=Solution]
Answer: a

Explanation: RR*=R+ as R+ means the occurrence to be at least once.
\end{tcolorbox}
\end{mdframed}
\newpage

\newtopic{Operators of Regular Expression}

\begin{mdframed}[linewidth=1pt, roundcorner=5pt, linecolor=gray!30, backgroundcolor=gray!5]
\textbf{Question 1} 

A finite automaton accepts which type of language:

a) Type 0

b) Type 1

c) Type 2

d) Type 3
\vspace{5pt} \par
\begin{tcolorbox}[colback=green!5, colframe=green!40!black, title=Solution]
Answer: d

Explanation: Type 3 refers to Regular Languages which is accepted by a finite automaton.
\end{tcolorbox}
\end{mdframed}
\newpage

\begin{mdframed}[linewidth=1pt, roundcorner=5pt, linecolor=gray!30, backgroundcolor=gray!5]
\textbf{Question 2} 

Which among the following are incorrect regular identities?

a) $\epsilon$R=R

b) $\epsilon$*=$\epsilon$

c) $\Phi$*=$\epsilon$

d) R$\Phi$=R
\vspace{5pt} \par
\begin{tcolorbox}[colback=green!5, colframe=green!40!black, title=Solution]
Answer: d

Explanation: There are few identities over Regular Expressions which include: R$\Phi$=$\Phi$R=$\Phi$$\neq$R
\end{tcolorbox}
\end{mdframed}
\newpage

\begin{mdframed}[linewidth=1pt, roundcorner=5pt, linecolor=gray!30, backgroundcolor=gray!5]
\textbf{Question 3} 

Simplify the following regular expression:

$\epsilon$+1*(011) *(1*(011) *) *

a) (1+011) *

b) (1*(011) *)

c) (1+(011) *) *

d) (1011) *
\vspace{5pt} \par
\begin{tcolorbox}[colback=green!5, colframe=green!40!black, title=Solution]
Answer: a

Explanation: $\epsilon$+1*(011) *(1*(011) *) *

$\epsilon$ + RR*= $\epsilon$ + R*R= $\epsilon$ + R+= R*
\end{tcolorbox}
\end{mdframed}
\newpage

\begin{mdframed}[linewidth=1pt, roundcorner=5pt, linecolor=gray!30, backgroundcolor=gray!5]
\textbf{Question 4} 

P, O, R be regular expression over $\Sigma$, P is not $\epsilon$, then

R=Q + RP has a unique solution:

a) Q*P

b) QP*

c) Q*P*

d) (P*O*) *
\vspace{5pt} \par
\begin{tcolorbox}[colback=green!5, colframe=green!40!black, title=Solution]
Answer: b

Explanation: The given statement is the Arden's Theorem and it tends to have a unique solution as QP*.

Let P and Q be regular expressions,

R=Q+RP

R=Q+(Q+RP) P

R=Q+((Q+RP) +RP) +P=Q+QP+RPP+RPP=Q+QP+(Q+RP) PP+(Q+RP) PP=Q+QP+QPP+RPPP+QPP+RPPP,

If we do this recursively, we get:

R= QP*
\end{tcolorbox}
\end{mdframed}
\newpage

\begin{mdframed}[linewidth=1pt, roundcorner=5pt, linecolor=gray!30, backgroundcolor=gray!5]
\textbf{Question 5} 

Arden's theorem is true for:

a) More than one initial states

b) Null transitions

c) Non-null transitions

d) None of the mentioned
\vspace{5pt} \par
\begin{tcolorbox}[colback=green!5, colframe=green!40!black, title=Solution]
Answer: c

Explanation: Arden's theorem strictly assumes the following;

a) No null transitions in the transition diagrams

b) True for only single initial state
\end{tcolorbox}
\end{mdframed}
\newpage

\begin{mdframed}[linewidth=1pt, roundcorner=5pt, linecolor=gray!30, backgroundcolor=gray!5]
\textbf{Question 6} 

The difference between number of states with regular expression (a + b) and (a + b) * is:

a) 1

b) 2

c) 3

d) 0
\vspace{5pt} \par
\begin{tcolorbox}[colback=green!5, colframe=green!40!black, title=Solution]
Answer: a

Explanation:
\end{tcolorbox}
\end{mdframed}
\newpage

\begin{mdframed}[linewidth=1pt, roundcorner=5pt, linecolor=gray!30, backgroundcolor=gray!5]
\textbf{Question 7} 

In order to represent a regular expression, the first step to create the transition diagram is:

a) Create the NFA using Null moves

b)  Null moves are not acceptable, thus should not be used

c) Predict the number of states to be used in order to construct the Regular expression

d) None of the mentioned
\vspace{5pt} \par
\begin{tcolorbox}[colback=green!5, colframe=green!40!black, title=Solution]
Answer: a

Explanation: Two steps are to be followed while converting a regular expression into a transition diagram:

a) Construct the NFA using null moves.

b) Remove the null transitions and convert it into its equivalent DFA.
\end{tcolorbox}
\end{mdframed}
\newpage

\begin{mdframed}[linewidth=1pt, roundcorner=5pt, linecolor=gray!30, backgroundcolor=gray!5]
\textbf{Question 8} 

(0+$\epsilon$) (1+$\epsilon$) represents

a) \{0, 1, 01, $\epsilon$\}

b) \{0, 1, $\epsilon$\}

c) \{0, 1, 01 ,11, 00, 10, $\epsilon$\}

d) \{0, 1\}
\vspace{5pt} \par
\begin{tcolorbox}[colback=green!5, colframe=green!40!black, title=Solution]
Answer: a

Explanation: The regular expression is fragmented and the set of the strings eligible is formed. ‘+' represents union while ‘.' Represents concatenation.
\end{tcolorbox}
\end{mdframed}
\newpage

\begin{mdframed}[linewidth=1pt, roundcorner=5pt, linecolor=gray!30, backgroundcolor=gray!5]
\textbf{Question 9} 

The minimum number of states required to automate the following Regular Expression:

(1) *(01+10) (1) *

a) 4

b) 3

c) 2

d) 5
\vspace{5pt} \par
\begin{tcolorbox}[colback=green!5, colframe=green!40!black, title=Solution]
Answer: a

Explanation: None.
\end{tcolorbox}
\end{mdframed}
\newpage

\begin{mdframed}[linewidth=1pt, roundcorner=5pt, linecolor=gray!30, backgroundcolor=gray!5]
\textbf{Question 10} 

Regular Expression denote precisely the \_\_\_\_\_\_\_\_ of Regular Language.

a) Class

b) Power Set

c) Super Set

d) None of the mentioned
\vspace{5pt} \par
\begin{tcolorbox}[colback=green!5, colframe=green!40!black, title=Solution]
Answer: a

Explanation: Regular Expression denote precisely the class of regular language. Given any regular expression, L(R) is a regular language. Given any regular language L, there is a regular expression R, such that L(R)=L.
\end{tcolorbox}
\end{mdframed}
\newpage

\newtopic{Building Regular Expressions}

\begin{mdframed}[linewidth=1pt, roundcorner=5pt, linecolor=gray!30, backgroundcolor=gray!5]
\textbf{Question 1} 

Which of the following is correct?

Statement 1: $\epsilon$ represents a single string in the set.

Statement 2: $\Phi$ represents the language that consist of no string.

a) Statement 1 and 2 both are correct

b) Statement 1 is false but 2 is correct

c) Statement 1 and 2 both are false

d) There is no difference between both the statements, $\epsilon$ and $\Phi$ are different notation for same reason
\vspace{5pt} \par
\begin{tcolorbox}[colback=green!5, colframe=green!40!black, title=Solution]
Answer: a

Explanation: $\epsilon$ represents a single string in the set namely, the empty string while Statement 2 is also correct.
\end{tcolorbox}
\end{mdframed}
\newpage

\begin{mdframed}[linewidth=1pt, roundcorner=5pt, linecolor=gray!30, backgroundcolor=gray!5]
\textbf{Question 2} 

The appropriate precedence order of operations over a Regular Language is

a) Kleene, Union, Concatenate

b) Kleene, Star, Union

c) Kleene, Dot, Union

d) Star, Union, Dot
\vspace{5pt} \par
\begin{tcolorbox}[colback=green!5, colframe=green!40!black, title=Solution]
Answer: c

Explanation: If a regular language expression is given, the appropriate order of precedence if the parenthesis is ignored is: Star or Kleene, Dot or Concatenation, Union or Plus.
\end{tcolorbox}
\end{mdframed}
\newpage

\begin{mdframed}[linewidth=1pt, roundcorner=5pt, linecolor=gray!30, backgroundcolor=gray!5]
\textbf{Question 3} 

Regular Expression R and the language it describes can be represented as:

a) R, R(L)

b) L(R), R(L)

c) R, L(R)

d) All of the mentioned
\vspace{5pt} \par
\begin{tcolorbox}[colback=green!5, colframe=green!40!black, title=Solution]
Answer: c

Explanation: When we wish to distinguish between a regular expression R and the language it represents; we write L(R) to be the language of R.
\end{tcolorbox}
\end{mdframed}
\newpage

\begin{mdframed}[linewidth=1pt, roundcorner=5pt, linecolor=gray!30, backgroundcolor=gray!5]
\textbf{Question 4} 

Let for $\Sigma$= \{0,1\} R= ($\Sigma$$\Sigma$$\Sigma$) *, the language of R would be

a) \{w | w is a string of odd length\}

b) \{w | w is a string of length multiple of 3\}

c) \{w | w is a string of length 3\}

d) All of the mentioned
\vspace{5pt} \par
\begin{tcolorbox}[colback=green!5, colframe=green!40!black, title=Solution]
Answer: b

Explanation: This regular expression can be used to eliminate the answers and get the result. The length can be even and as well more than 3 when R= ($\Sigma$$\Sigma$$\Sigma$) ($\Sigma$$\Sigma$$\Sigma$) (particular case).
\end{tcolorbox}
\end{mdframed}
\newpage

\begin{mdframed}[linewidth=1pt, roundcorner=5pt, linecolor=gray!30, backgroundcolor=gray!5]
\textbf{Question 5} 

If $\Sigma$= \{0,1\}, then $\Phi$* will result to:

a) $\epsilon$

b) $\Phi$

c) $\Sigma$

d) None of the mentioned
\vspace{5pt} \par
\begin{tcolorbox}[colback=green!5, colframe=green!40!black, title=Solution]
Answer: a

Explanation: The star operation brings together any number of strings from the language to get a string in the result. If the language is empty, the star operation can put together 0 strings, resulting only the empty string.
\end{tcolorbox}
\end{mdframed}
\newpage

\begin{mdframed}[linewidth=1pt, roundcorner=5pt, linecolor=gray!30, backgroundcolor=gray!5]
\textbf{Question 6} 

The given NFA represents which of the following NFA

a) (ab U a) *

b) (a*b* U a*)

c) (ab U a*)

d) (ab)* U a*
\vspace{5pt} \par
\begin{tcolorbox}[colback=green!5, colframe=green!40!black, title=Solution]
Answer: a

Explanation: The Regular expression (ab U a) * is converted to NFA in a sequence of stages as it can be clearly seen in the diagram. This NFA consist of 8 stated while its minimized form only contains 2 states.
\end{tcolorbox}
\end{mdframed}
\newpage

\begin{mdframed}[linewidth=1pt, roundcorner=5pt, linecolor=gray!30, backgroundcolor=gray!5]
\textbf{Question 7} 

Which of the following represents a language which has no pair of consecutive 1's if $\Sigma$= \{0,1\}?

a) (0+10)*(1+$\epsilon$)

b) (0+10)*(1+$\epsilon$)*

c) (0+101)*(0+$\epsilon$)

d) (1+010)*(1+$\epsilon$)
\vspace{5pt} \par
\begin{tcolorbox}[colback=green!5, colframe=green!40!black, title=Solution]
Answer: a

Explanation: All the options except ‘a' accept those strings which comprises minimum one pair of 1's together.
\end{tcolorbox}
\end{mdframed}
\newpage

\begin{mdframed}[linewidth=1pt, roundcorner=5pt, linecolor=gray!30, backgroundcolor=gray!5]
\textbf{Question 8} 

The finite automata accept the following languages:

a) Context Free Languages

b) Context Sensitive Languages

c) Regular Languages

d) All the mentioned
\vspace{5pt} \par
\begin{tcolorbox}[colback=green!5, colframe=green!40!black, title=Solution]
Answer: c

Explanation: A finite automaton accepts the languages which are regular and for which a DFA can be constructed.
\end{tcolorbox}
\end{mdframed}
\newpage

\begin{mdframed}[linewidth=1pt, roundcorner=5pt, linecolor=gray!30, backgroundcolor=gray!5]
\textbf{Question 9} 

(a + b*c) most correctly represents:

a) (a +b) *c

b) (a)+((b)*.c)

c) (a + (b*)).c

d) a+ ((b*).c)
\vspace{5pt} \par
\begin{tcolorbox}[colback=green!5, colframe=green!40!black, title=Solution]
Answer: d

Explanation: Following the rules of precedence, Kleene or star operation would be done first, then concatenation and finally union or plus operation.
\end{tcolorbox}
\end{mdframed}
\newpage

\begin{mdframed}[linewidth=1pt, roundcorner=5pt, linecolor=gray!30, backgroundcolor=gray!5]
\textbf{Question 10} 

Which of the following regular expressions represents the set of strings which do not contain a substring ‘rt' if $\Sigma$= \{r, t\}

a) (rt)*

b) (tr)*

c) (r*t*)

d) (t*r*)
\vspace{5pt} \par
\begin{tcolorbox}[colback=green!5, colframe=green!40!black, title=Solution]
Answer: d

Explanation: As Kleene operation is not on the whole of the substring, it will not repeat and maintain the order of t, r.
\end{tcolorbox}
\end{mdframed}
\newpage

\begin{mdframed}[linewidth=1pt, roundcorner=5pt, linecolor=gray!30, backgroundcolor=gray!5]
\textbf{Question 11} 

According to the precedence rules, x-y-z is equivalent to which of the following?

a) (x-y)-z

b) x-(y-z)

c) Both (x-y)-z and x-(y-z)

d) None of the mentioned
\vspace{5pt} \par
\begin{tcolorbox}[colback=green!5, colframe=green!40!black, title=Solution]
Answer: a

Explanation: In arithmetic, we group two of the same operators from the left, hence x-y-z is equivalent to (x-y)-z and not x-(y—z).
\end{tcolorbox}
\end{mdframed}
\newpage

\begin{mdframed}[linewidth=1pt, roundcorner=5pt, linecolor=gray!30, backgroundcolor=gray!5]
\textbf{Question 12} 

Dot operator in regular expression resembles which of the following?

a) Expressions are juxtaposed

b) Expressions are multiplied

c) Cross operation

d) None of the mentioned
\vspace{5pt} \par
\begin{tcolorbox}[colback=green!5, colframe=green!40!black, title=Solution]
Answer: a

Explanation: Dot operation or concatenation operation means that the two expressions are juxtaposed i.e. there are no intervening operators in between. In fact, UNIX regular expressions use the dot for an entirely different purpose: representing any ASCII character.
\end{tcolorbox}
\end{mdframed}
\newpage

\begin{mdframed}[linewidth=1pt, roundcorner=5pt, linecolor=gray!30, backgroundcolor=gray!5]
\textbf{Question 13} 

Which among the following is not an associative operation?

a) Union

b) Concatenation

c) Dot

d) None of the mentioned
\vspace{5pt} \par
\begin{tcolorbox}[colback=green!5, colframe=green!40!black, title=Solution]
Answer: d

Explanation: It does not matter in which order we group the expression with the operators as they are associative. If one gets a chance to group the expression, one should group them from left for convenience. For instance, 012 is grouped as (01)2.
\end{tcolorbox}
\end{mdframed}
\newpage

\begin{mdframed}[linewidth=1pt, roundcorner=5pt, linecolor=gray!30, backgroundcolor=gray!5]
\textbf{Question 14} 

Which among the following is equivalent to the given regular expression?

01*+1

a) (01)*+1

b) 0((1)*+1)

c) (0(1)*)+1

d) ((0*1)1*)*
\vspace{5pt} \par
\begin{tcolorbox}[colback=green!5, colframe=green!40!black, title=Solution]
Answer: c

Explanation: Using the rules of precedence on the give expression, c is the appropriate choice with the order of: Bracket$>$Kleene$>$Dot$>$Union
\end{tcolorbox}
\end{mdframed}
\newpage

\newtopic{DFA to Regular Expressions}

\begin{mdframed}[linewidth=1pt, roundcorner=5pt, linecolor=gray!30, backgroundcolor=gray!5]
\textbf{Question 1} 

Which of the following is same as the given DFA?


 
\begin{center}\includegraphics[width=0.5\textwidth]{automata-theory-questions-answers-dfa-regular-expressions-q1.png}\end{center}
 

a) (0+1)*001(0+1)*

b) 1*001(0+1)*

c) (01)*(0+0+1)(01)*

d) None of the mentioned
\vspace{5pt} \par
\begin{tcolorbox}[colback=green!5, colframe=green!40!black, title=Solution]
Answer: a

Explanation: There needs to be 001 together in the string as an essential substring. Thus, the other components can be anything, 0 or 1 or e.
\end{tcolorbox}
\end{mdframed}
\newpage

\begin{mdframed}[linewidth=1pt, roundcorner=5pt, linecolor=gray!30, backgroundcolor=gray!5]
\textbf{Question 2} 

Which of the following statements is not true?

a) Every language defined by any of the automata is also defined by a regular expression

b) Every language defined by a regular expression can be represented using a DFA

c) Every language defined by a regular expression can be represented using NFA with e moves

d) Regular expression is just another representation for any automata definition
\vspace{5pt} \par
\begin{tcolorbox}[colback=green!5, colframe=green!40!black, title=Solution]
Answer: b

Explanation: Using NFA with e moves, we can represent all the regular expressions as an automata. As regular expressions include e, we need to use  e moves.
\end{tcolorbox}
\end{mdframed}
\newpage

\begin{mdframed}[linewidth=1pt, roundcorner=5pt, linecolor=gray!30, backgroundcolor=gray!5]
\textbf{Question 3} 

The total number of states required to automate the given regular expression

(00)*(11)*

a) 3

b) 4

c) 5

d) 6
\vspace{5pt} \par
\begin{tcolorbox}[colback=green!5, colframe=green!40!black, title=Solution]
Answer: c

Explanation:  
 
\begin{center}\includegraphics[width=0.5\textwidth]{automata-theory-questions-answers-dfa-regular-expressions-q3.png}\end{center}
\end{tcolorbox}
\end{mdframed}
\newpage

\begin{mdframed}[linewidth=1pt, roundcorner=5pt, linecolor=gray!30, backgroundcolor=gray!5]
\textbf{Question 4} 

Which of the given regular expressions correspond to the automata shown?


 
\begin{center}\includegraphics[width=0.5\textwidth]{automata-theory-questions-answers-dfa-regular-expressions-q4.png}\end{center}
 

a) (110+1)*0

b) (11+110)*1

c) (110+11)*0

d) (1+110)*1
\vspace{5pt} \par
\begin{tcolorbox}[colback=green!5, colframe=green!40!black, title=Solution]
Answer: c

Explanation: There is no state change for union operation, but has two different paths while for concatenation or dot operation, we have a state change for every element of the string.
\end{tcolorbox}
\end{mdframed}
\newpage

\begin{mdframed}[linewidth=1pt, roundcorner=5pt, linecolor=gray!30, backgroundcolor=gray!5]
\textbf{Question 5} 

Generate a regular expression for the following problem statement:

Password Validation: String should be 8-15 characters long. String must contain a number, an Uppercase letter and a Lower case letter.

a) \textasciicircum{}(?=.*[a-z])(?=.*[A-Z])(?=.*\textbackslash d).\{8,15\}\$

b) \textasciicircum{}(?=.*[a-z])(?=.*[A-Z])(?=.*\textbackslash d).\{9,16\}\$

c) \textasciicircum{}(?=.[a-z])(?=.[A-Z])(?=.\textbackslash d).\{8,15\}\$

d) None of the mentioned
\vspace{5pt} \par
\begin{tcolorbox}[colback=green!5, colframe=green!40!black, title=Solution]
Answer: a

Explanation: Passwords like abc123, 123XYZ, should not be accepted . If one also wants to include special characters as one of the constraint, one can use the following regular expression:

               \textasciicircum{}(?=.*[a-z])(?=.*[A-Z])(?=.*\textbackslash d)(?=.*[\textasciicircum{}\textbackslash da-za-Z]).\{8,15\}\$
\end{tcolorbox}
\end{mdframed}
\newpage

\begin{mdframed}[linewidth=1pt, roundcorner=5pt, linecolor=gray!30, backgroundcolor=gray!5]
\textbf{Question 6} 

Generate a regular expression for the following problem statement:

P(x): String of length 6 or less for å=\{0,1\}*

a) (1+0+e)6

b) (10)6

c) (1+0)(1+0)(1+0)(1+0)(1+0)(1+0)

d) More than one of the mentioned is correct
\vspace{5pt} \par
\begin{tcolorbox}[colback=green!5, colframe=green!40!black, title=Solution]
Answer: a

Explanation: As the input variables are under Kleene Operation, we need to include e,thus option c is not correct,thereby option (a) is the right answer.
\end{tcolorbox}
\end{mdframed}
\newpage

\begin{mdframed}[linewidth=1pt, roundcorner=5pt, linecolor=gray!30, backgroundcolor=gray!5]
\textbf{Question 7} 

The minimum number of states required in a DFA (along with a dumping state)  to check whether the 3rd bit is 1 or not for |n|$>$=3

a) 3

b) 4

c) 5

d) 1
\vspace{5pt} \par
\begin{tcolorbox}[colback=green!5, colframe=green!40!black, title=Solution]
Answer: c

Explanation:
 
\begin{center}\includegraphics[width=0.5\textwidth]{automata-theory-questions-answers-dfa-regular-expressions-q7.png}\end{center}
\end{tcolorbox}
\end{mdframed}
\newpage

\begin{mdframed}[linewidth=1pt, roundcorner=5pt, linecolor=gray!30, backgroundcolor=gray!5]
\textbf{Question 8} 

Which of the regular expressions corresponds to the given problem statement:

P(x): Express the identifiers in C Programming language

l=letters

d=digits

a) (l+\_)(d+\_)*

b) (l+d+\_)*

c) (l+\_)(l+d+\_)*

d) (\_+d)(l+d+\_)*
\vspace{5pt} \par
\begin{tcolorbox}[colback=green!5, colframe=green!40!black, title=Solution]
Answer: c

Explanation: Identifiers in C Programming Language follows the following identifiers rule:

a) The name of the identifier should not begin with a digit.

b) It can only begin with a letter or a underscore.

c) It can be of length 1 or more.
\end{tcolorbox}
\end{mdframed}
\newpage

\begin{mdframed}[linewidth=1pt, roundcorner=5pt, linecolor=gray!30, backgroundcolor=gray!5]
\textbf{Question 9} 

Generate a regular expression for the given language:l

L(x): \{xÎ\{0,1\}*| x ends with 1 nd does not contain a substring 01\}

a) (0+01)*

b) (0+01)*1

c) (0+01)*(1+01)

d) All of the mentioned
\vspace{5pt} \par
\begin{tcolorbox}[colback=green!5, colframe=green!40!black, title=Solution]
Answer: c

Explanation: (a) and (b) are the general cases where we restrict the acceptance of a string witrh substring 00 but we ignore the case where the string needs to end with 1 which therby, does not allows the acceptance of e.
\end{tcolorbox}
\end{mdframed}
\newpage

\begin{mdframed}[linewidth=1pt, roundcorner=5pt, linecolor=gray!30, backgroundcolor=gray!5]
\textbf{Question 10} 

The minimum number of transitions to pass to reach the final state as per the following regular expression is:

 \{a,b\}*\{baaa\}

a) 4

b) 5

c) 6

d) 3
\vspace{5pt} \par
\begin{tcolorbox}[colback=green!5, colframe=green!40!black, title=Solution]
Answer: a

Explanation: 
 
\begin{center}\includegraphics[width=0.5\textwidth]{automata-theory-questions-answers-dfa-regular-expressions-q10.png}\end{center}
\end{tcolorbox}
\end{mdframed}
\newpage

\newtopic{Conversion by Eliminating states}

\begin{mdframed}[linewidth=1pt, roundcorner=5pt, linecolor=gray!30, backgroundcolor=gray!5]
\textbf{Question 1} 

Which of the following is an utility of state elimination phenomenon?

a) DFA to NFA

b) NFA to DFA

c) DFA to Regular Expression

d) All of the mentioned
\vspace{5pt} \par
\begin{tcolorbox}[colback=green!5, colframe=green!40!black, title=Solution]
Answer: c

Explanation: We use this algorithm to simplify a finite automaton to regular expression or vice versa. We eliminate states while converting given finite automata to its corresponding regular expression.
\end{tcolorbox}
\end{mdframed}
\newpage

\begin{mdframed}[linewidth=1pt, roundcorner=5pt, linecolor=gray!30, backgroundcolor=gray!5]
\textbf{Question 2} 

If we have more than one accepting states or an accepting state with an outdegree, which of the following actions will be taken?

a) addition of new state

b) removal of a state

c) make the newly added state as final

d) more than one option is correct
\vspace{5pt} \par
\begin{tcolorbox}[colback=green!5, colframe=green!40!black, title=Solution]
Answer: d

Explanation: If there is more than one accepting state or if the single accepting state as an out degree, add a new accepting state, make all other states non accepting, and hold an e-transitions from each former accepting state to the new accepting state.
\end{tcolorbox}
\end{mdframed}
\newpage

\begin{mdframed}[linewidth=1pt, roundcorner=5pt, linecolor=gray!30, backgroundcolor=gray!5]
\textbf{Question 3} 

Which of the following is not a step in elimination of states procedure?

a) Unifying all the final states into one using e-transitions

b) Unify single transitions to multi transitions that contains union of input

c) Remove states until there is only starting and accepting states

d) Get the resulting regular expression by direct calculation
\vspace{5pt} \par
\begin{tcolorbox}[colback=green!5, colframe=green!40!black, title=Solution]
Answer: b

Explanation: While eliminating the states, we unify multiple transitions to one transition that contains union of input and not the vice versa.
\end{tcolorbox}
\end{mdframed}
\newpage

\begin{mdframed}[linewidth=1pt, roundcorner=5pt, linecolor=gray!30, backgroundcolor=gray!5]
\textbf{Question 4} 

Can the given state diagram be reduced?


 
\begin{center}\includegraphics[width=0.5\textwidth]{automata-theory-questions-answers-mcqs-q4.png}\end{center}
 

a) Yes

b) No
\vspace{5pt} \par
\begin{tcolorbox}[colback=green!5, colframe=green!40!black, title=Solution]
Answer: a

Explanation: The state q2 can be eliminated with ease and the reduced state diagram can be represented as:


 
\begin{center}\includegraphics[width=0.5\textwidth]{automata-theory-questions-answers-mcqs-q4a.png}\end{center}
\end{tcolorbox}
\end{mdframed}
\newpage

\begin{mdframed}[linewidth=1pt, roundcorner=5pt, linecolor=gray!30, backgroundcolor=gray!5]
\textbf{Question 5} 

Which of the following methods is suitable for conversion of DFA to RE?

a) Brzozowski method

b) Arden's method

c) Walter's method

d) All of the mentioned
\vspace{5pt} \par
\begin{tcolorbox}[colback=green!5, colframe=green!40!black, title=Solution]
Answer: a

Explanation: Brzozowski method takes a unique approach to generating regular expressions. We create a system of regular expressions with one regular expression unknown for each state in M, and then we solve the system for R$\Lambda$ where R$\Lambda$ is the regular expression associated with starting state q$\Lambda$.
\end{tcolorbox}
\end{mdframed}
\newpage

\begin{mdframed}[linewidth=1pt, roundcorner=5pt, linecolor=gray!30, backgroundcolor=gray!5]
\textbf{Question 6} 

State true or false:

Statement: The state removal approach identifies patterns within the graph and removes state, building up regular expressions along each transition.

a) true

b) false
\vspace{5pt} \par
\begin{tcolorbox}[colback=green!5, colframe=green!40!black, title=Solution]
Answer: a

Explanation: This method has the advantage over the transitive closure technique as it can easily be visualized.
\end{tcolorbox}
\end{mdframed}
\newpage

\begin{mdframed}[linewidth=1pt, roundcorner=5pt, linecolor=gray!30, backgroundcolor=gray!5]
\textbf{Question 7} 

The behaviour of NFA can be simulated using DFA.

a) always

b) never

c) sometimes

d) none of the mentioned
\vspace{5pt} \par
\begin{tcolorbox}[colback=green!5, colframe=green!40!black, title=Solution]
Answer: a

Explanation: For every NFA, there exists an equivalent DFA and vice versa.
\end{tcolorbox}
\end{mdframed}
\newpage

\begin{mdframed}[linewidth=1pt, roundcorner=5pt, linecolor=gray!30, backgroundcolor=gray!5]
\textbf{Question 8} 

It is suitable to use \_\_\_\_\_\_\_\_\_\_\_\_ method/methods to convert a DFA to regular expression.

a) Transitive Closure properties

b) Brzozowski method

c) State elimination method

d) All of the mentioned
\vspace{5pt} \par
\begin{tcolorbox}[colback=green!5, colframe=green!40!black, title=Solution]
Answer: d

Explanation: For converting RE to DFA , first we convert RE to NFA (Thompson Construction), and then NFA is converted into DFA(Subset Construction).
\end{tcolorbox}
\end{mdframed}
\newpage

\begin{mdframed}[linewidth=1pt, roundcorner=5pt, linecolor=gray!30, backgroundcolor=gray!5]
\textbf{Question 9} 

State true or false:

Statement: For every removed state, there is a regular expression produced.

a) true

b) false
\vspace{5pt} \par
\begin{tcolorbox}[colback=green!5, colframe=green!40!black, title=Solution]
Answer: a

Explanation: For every state which is eliminated, a new regular expression is produced. The newly generated regular expression act as an input for a state which is next to removed state.
\end{tcolorbox}
\end{mdframed}
\newpage

\begin{mdframed}[linewidth=1pt, roundcorner=5pt, linecolor=gray!30, backgroundcolor=gray!5]
\textbf{Question 10} 

Is it possible to obtain more than one regular expression from a given DFA using the state elimination method?

a) Yes

b) No
\vspace{5pt} \par
\begin{tcolorbox}[colback=green!5, colframe=green!40!black, title=Solution]
Answer: a

Explanation: Using different sequence of removal of state, we can have different possible solution of regular expressions. For n-state deterministic finite automata excluding starting and final states, n! Removal sequences are there. It is very tough to try all the possible removal sequences for smaller expressions.
\end{tcolorbox}
\end{mdframed}
\newpage

\newtopic{Regular Language \& Expression – 1}

\begin{mdframed}[linewidth=1pt, roundcorner=5pt, linecolor=gray!30, backgroundcolor=gray!5]
\textbf{Question 1} 

A regular language over an alphabet a is one that can be obtained from

a) union

b) concatenation

c) kleene

d) All of the mentioned
\vspace{5pt} \par
\begin{tcolorbox}[colback=green!5, colframe=green!40!black, title=Solution]
Answer: d

Explanation: None.
\end{tcolorbox}
\end{mdframed}
\newpage

\begin{mdframed}[linewidth=1pt, roundcorner=5pt, linecolor=gray!30, backgroundcolor=gray!5]
\textbf{Question 2} 

Regular expression \{0,1\} is equivalent to

a) 0 U 1

b) 0 / 1

c) 0 + 1

d) All of the mentioned
\vspace{5pt} \par
\begin{tcolorbox}[colback=green!5, colframe=green!40!black, title=Solution]
Answer: d

Explanation: All are equivalent to union operation.
\end{tcolorbox}
\end{mdframed}
\newpage

\begin{mdframed}[linewidth=1pt, roundcorner=5pt, linecolor=gray!30, backgroundcolor=gray!5]
\textbf{Question 3} 

Precedence of regular expression in decreasing order is

a) * , . , +

b) . , * , +

c) . , + , *

d) + , a , *
\vspace{5pt} \par
\begin{tcolorbox}[colback=green!5, colframe=green!40!black, title=Solution]
Answer: a

Explanation: None.
\end{tcolorbox}
\end{mdframed}
\newpage

\begin{mdframed}[linewidth=1pt, roundcorner=5pt, linecolor=gray!30, backgroundcolor=gray!5]
\textbf{Question 4} 

Regular expression $\Phi$* is equivalent to

a) $\epsilon$

b) $\Phi$

c) 0

d) 1
\vspace{5pt} \par
\begin{tcolorbox}[colback=green!5, colframe=green!40!black, title=Solution]
Answer: a

Explanation: None.
\end{tcolorbox}
\end{mdframed}
\newpage

\begin{mdframed}[linewidth=1pt, roundcorner=5pt, linecolor=gray!30, backgroundcolor=gray!5]
\textbf{Question 5} 

a? is equivalent to

a) a

b) a+$\Phi$

c) a+$\epsilon$

d) wrong expression
\vspace{5pt} \par
\begin{tcolorbox}[colback=green!5, colframe=green!40!black, title=Solution]
Answer: c

Explanation: Zero or one time repetition of previous character .
\end{tcolorbox}
\end{mdframed}
\newpage

\begin{mdframed}[linewidth=1pt, roundcorner=5pt, linecolor=gray!30, backgroundcolor=gray!5]
\textbf{Question 6} 

$\epsilon$L is equivalent to

a) $\epsilon$

b) $\Phi$

c) L

d) $\Phi$$\epsilon$
\vspace{5pt} \par
\begin{tcolorbox}[colback=green!5, colframe=green!40!black, title=Solution]
Answer: c

Explanation: None.
\end{tcolorbox}
\end{mdframed}
\newpage

\begin{mdframed}[linewidth=1pt, roundcorner=5pt, linecolor=gray!30, backgroundcolor=gray!5]
\textbf{Question 7} 

(a+b)* is equivalent to

a) b*a*

b) (a*b*)*

c) a*b*

d) none of the mentioned
\vspace{5pt} \par
\begin{tcolorbox}[colback=green!5, colframe=green!40!black, title=Solution]
Answer: b

Explanation: None.
\end{tcolorbox}
\end{mdframed}
\newpage

\begin{mdframed}[linewidth=1pt, roundcorner=5pt, linecolor=gray!30, backgroundcolor=gray!5]
\textbf{Question 8} 

$\Phi$L is equivalent to

a) L$\Phi$ \& $\Phi$

b) $\Phi$ \& L

c) L \& L

d) $\epsilon$ \& L
\vspace{5pt} \par
\begin{tcolorbox}[colback=green!5, colframe=green!40!black, title=Solution]
Answer: a

Explanation: None.
\end{tcolorbox}
\end{mdframed}
\newpage

\begin{mdframed}[linewidth=1pt, roundcorner=5pt, linecolor=gray!30, backgroundcolor=gray!5]
\textbf{Question 9} 

Which of the following pair of regular expression are not equivalent?

a) 1(01)* and (10)*1

b) x(xx)* and (xx)*x

c) (ab)* and a*b*

d) x+ and x*x+
\vspace{5pt} \par
\begin{tcolorbox}[colback=green!5, colframe=green!40!black, title=Solution]
Answer: c

Explanation: (ab)*=(a*b*)*.
\end{tcolorbox}
\end{mdframed}
\newpage

\begin{mdframed}[linewidth=1pt, roundcorner=5pt, linecolor=gray!30, backgroundcolor=gray!5]
\textbf{Question 10} 

Consider following regular expression

   i) (a/b)*    ii) (a*/b*)*   iii) (($\epsilon$/a)b*)*

Which of the following statements is correct

a) i,ii are equal and ii,iii are not

b) i,ii are equal and i,iii are not

c) ii,iii are equal and i,ii are not

d) all are equal
\vspace{5pt} \par
\begin{tcolorbox}[colback=green!5, colframe=green!40!black, title=Solution]
Answer: d

Explanation: All are equivalent to (a+b)*.
\end{tcolorbox}
\end{mdframed}
\newpage

\newtopic{Regular Language \& Expression – 2}

\begin{mdframed}[linewidth=1pt, roundcorner=5pt, linecolor=gray!30, backgroundcolor=gray!5]
\textbf{Question 1} 

How many strings of length less than 4 contains the language described by the regular expression (x+y)*y(a+ab)*?

a) 7

b) 10

c) 12

d) 11
\vspace{5pt} \par
\begin{tcolorbox}[colback=green!5, colframe=green!40!black, title=Solution]
Answer: c

Explanation: string of length 0 = Not possible (because y is always present).

string of length 1 = 1 (y)

string of length 2 = 3 (xy,yy,ya)

string of length 3 = 8 (xxy,xyy,yxy,yyy,yaa,yab,xya,yya)
\end{tcolorbox}
\end{mdframed}
\newpage

\begin{mdframed}[linewidth=1pt, roundcorner=5pt, linecolor=gray!30, backgroundcolor=gray!5]
\textbf{Question 2} 

Which of the following is true?

a) (01)*0 = 0(10)*

b) (0+1)*0(0+1)*1(0+1) = (0+1)*01(0+1)*

c) (0+1)*01(0+1)*+1*0* = (0+1)*

d) All of the mentioned
\vspace{5pt} \par
\begin{tcolorbox}[colback=green!5, colframe=green!40!black, title=Solution]
Answer: d

Explanation: None.
\end{tcolorbox}
\end{mdframed}
\newpage

\begin{mdframed}[linewidth=1pt, roundcorner=5pt, linecolor=gray!30, backgroundcolor=gray!5]
\textbf{Question 3} 

A language is regular if and only if

a) accepted by DFA

b) accepted by PDA

c) accepted by LBA

d) accepted by Turing machine
\vspace{5pt} \par
\begin{tcolorbox}[colback=green!5, colframe=green!40!black, title=Solution]
Answer: a

Explanation: All of above machine can accept regular language but all string accepted by machine is regular only for DFA.
\end{tcolorbox}
\end{mdframed}
\newpage

\begin{mdframed}[linewidth=1pt, roundcorner=5pt, linecolor=gray!30, backgroundcolor=gray!5]
\textbf{Question 4} 

Regular grammar is

a) context free grammar

b) non context free grammar

c) english grammar

d) none of the mentioned
\vspace{5pt} \par
\begin{tcolorbox}[colback=green!5, colframe=green!40!black, title=Solution]
Answer: a

Explanation: Regular grammar is subset of context free grammar.
\end{tcolorbox}
\end{mdframed}
\newpage

\begin{mdframed}[linewidth=1pt, roundcorner=5pt, linecolor=gray!30, backgroundcolor=gray!5]
\textbf{Question 5} 

Let the class of language accepted by finite state machine be L1 and the class of languages represented by regular expressions be L2 then

a) L1$<$L2

b) L1$>$=L2

c) L1 U L2 = .*

d) L1=L2
\vspace{5pt} \par
\begin{tcolorbox}[colback=green!5, colframe=green!40!black, title=Solution]
Answer: d

Explanation: Finite state machine and regular expression have same power to express a language.
\end{tcolorbox}
\end{mdframed}
\newpage

\begin{mdframed}[linewidth=1pt, roundcorner=5pt, linecolor=gray!30, backgroundcolor=gray!5]
\textbf{Question 6} 

Which of the following is not a regular expression?

a) [(a+b)*-(aa+bb)]*

b) [(0+1)-(0b+a1)*(a+b)]*

c) (01+11+10)*

d) (1+2+0)*(1+2)*
\vspace{5pt} \par
\begin{tcolorbox}[colback=green!5, colframe=green!40!black, title=Solution]
Answer: b

Explanation: Except b all are regular expression*.
\end{tcolorbox}
\end{mdframed}
\newpage

\begin{mdframed}[linewidth=1pt, roundcorner=5pt, linecolor=gray!30, backgroundcolor=gray!5]
\textbf{Question 7} 

Regular expression are

a) Type 0 language

b) Type 1 language

c) Type 2 language

d) Type 3 language
\vspace{5pt} \par
\begin{tcolorbox}[colback=green!5, colframe=green!40!black, title=Solution]
Answer: d

Explanation: According to Chomsky hierarchy,

Type 0 – Unrestricted Grammar.

Type 1 – Context Sensitive Grammar.

Type 2 – Context Free Grammar.

Type 3 – Regular Grammar.
\end{tcolorbox}
\end{mdframed}
\newpage

\begin{mdframed}[linewidth=1pt, roundcorner=5pt, linecolor=gray!30, backgroundcolor=gray!5]
\textbf{Question 8} 

Which of the following is true?

a) Every subset of a regular set is regular

b) Every finite subset of non-regular set is regular

c) The union of two non regular set is not regular

d) Infinite union of finite set is regular
\vspace{5pt} \par
\begin{tcolorbox}[colback=green!5, colframe=green!40!black, title=Solution]
Answer: b

Explanation: None.
\end{tcolorbox}
\end{mdframed}
\newpage

\begin{mdframed}[linewidth=1pt, roundcorner=5pt, linecolor=gray!30, backgroundcolor=gray!5]
\textbf{Question 9} 

L and ~L are recursive enumerable then L is

a) Regular

b) Context free

c) Context sensitive

d) Recursive
\vspace{5pt} \par
\begin{tcolorbox}[colback=green!5, colframe=green!40!black, title=Solution]
Answer: d

Explanation:If L is recursive enumerable and its complement too if and only if L is recursive.
\end{tcolorbox}
\end{mdframed}
\newpage

\begin{mdframed}[linewidth=1pt, roundcorner=5pt, linecolor=gray!30, backgroundcolor=gray!5]
\textbf{Question 10} 

Regular expressions are closed under

a) Union

b) Intersection

c) Kleen star

d) All of the mentioned
\vspace{5pt} \par
\begin{tcolorbox}[colback=green!5, colframe=green!40!black, title=Solution]
Answer: d

Explanation: According to definition of regular expression.
\end{tcolorbox}
\end{mdframed}
\newpage

\newtopic{Converting Regular Expressions to Automata}

\begin{mdframed}[linewidth=1pt, roundcorner=5pt, linecolor=gray!30, backgroundcolor=gray!5]
\textbf{Question 1} 

What kind of expressions do we used for pattern matching?

a) Regular Expression

b) Rational Expression

c) Regular \& Rational Expression

d) None of the mentioned
\vspace{5pt} \par
\begin{tcolorbox}[colback=green!5, colframe=green!40!black, title=Solution]
Answer: c

Explanation: In automata theory, Regular Expression(sometimes also called the Rational Expression  ) is a sequence or set of characters that define a search pattern, mainly for the use in pattern matching with strings or string matching.
\end{tcolorbox}
\end{mdframed}
\newpage

\begin{mdframed}[linewidth=1pt, roundcorner=5pt, linecolor=gray!30, backgroundcolor=gray!5]
\textbf{Question 2} 

Which of the following do Regexps do not find their use in?

a) search engines

b) word processors

c) sed

d) none of the mentioned
\vspace{5pt} \par
\begin{tcolorbox}[colback=green!5, colframe=green!40!black, title=Solution]
Answer: d

Explanation: Regexp processors are found in several search engines, seach and replace mechanisms, and text processing utilities.
\end{tcolorbox}
\end{mdframed}
\newpage

\begin{mdframed}[linewidth=1pt, roundcorner=5pt, linecolor=gray!30, backgroundcolor=gray!5]
\textbf{Question 3} 

Which of the following languages have built in regexps support?

a) Perl

b) Java

c) Python

d) C++
\vspace{5pt} \par
\begin{tcolorbox}[colback=green!5, colframe=green!40!black, title=Solution]
Answer: a

Explanation: Many languages come with built in support of regexps like Perl, Javascript, Ruby etc. While some provide support using standard libraries like .NET, Java, Python, C++, C and POSIX.
\end{tcolorbox}
\end{mdframed}
\newpage

\begin{mdframed}[linewidth=1pt, roundcorner=5pt, linecolor=gray!30, backgroundcolor=gray!5]
\textbf{Question 4} 

The following is/are an approach to process a regexp:

a) Construction of NFA and subsequently, a DFA

b) Thompson's Contruction Algorithm

c) Thompson's Contruction Algorithm \& Construction of NFA and subsequently, a DFA

d) None of the mentioned
\vspace{5pt} \par
\begin{tcolorbox}[colback=green!5, colframe=green!40!black, title=Solution]
Answer: c

Explanation: A regexp processor translates the syntax into internal representation which can be executed and matched with a string and that internal representation can have several approaches like the ones mentioned.
\end{tcolorbox}
\end{mdframed}
\newpage

\begin{mdframed}[linewidth=1pt, roundcorner=5pt, linecolor=gray!30, backgroundcolor=gray!5]
\textbf{Question 5} 

Are the given two patterns equivalent?

(1) gray|grey

(2) gr(a|e)y

a) yes

b) no
\vspace{5pt} \par
\begin{tcolorbox}[colback=green!5, colframe=green!40!black, title=Solution]
Answer: a

Explanation: Paranthesis can be used to define the scope and precedence of operators. Thus, both the expression represents the same pattern.
\end{tcolorbox}
\end{mdframed}
\newpage

\begin{mdframed}[linewidth=1pt, roundcorner=5pt, linecolor=gray!30, backgroundcolor=gray!5]
\textbf{Question 6} 

Which of the following are not quantifiers?

a) Kleene plus +

b) Kleene star *

c) Question mark ?

d) None of the mentioned
\vspace{5pt} \par
\begin{tcolorbox}[colback=green!5, colframe=green!40!black, title=Solution]
Answer: d

Explanation: A quantifier after a  token specifies how often the preceding element is allowed to occur. ?, *, +, \{n\}, \{min, \}, \{min, max\} are few quantifiers we use in regexps implementations.
\end{tcolorbox}
\end{mdframed}
\newpage

\begin{mdframed}[linewidth=1pt, roundcorner=5pt, linecolor=gray!30, backgroundcolor=gray!5]
\textbf{Question 7} 

Which of the following cannot be used to decide whether and how a given regexp matches a string:

a) NFA to DFA

b) Lazy DFA algorithm

c) Backtracking

d) None of the mentioned
\vspace{5pt} \par
\begin{tcolorbox}[colback=green!5, colframe=green!40!black, title=Solution]
Answer: d

Explanation: There are at least three algorithms which decides for us, whether and how a regexp matches a string which included the transformation of Non deterministic automaton to deterministic finite automaton, The lazy DFA algorithm where one simulates the NFA directly, building each DFA on demand and then discarding it at the next step and the process of backtracking whose running time is exponential.
\end{tcolorbox}
\end{mdframed}
\newpage

\begin{mdframed}[linewidth=1pt, roundcorner=5pt, linecolor=gray!30, backgroundcolor=gray!5]
\textbf{Question 8} 

What does the following segment of code output?


\$string1 
=
 
"Hello World
\textbackslash n
"
;


if
 
(
\$string1 
=
~ m
/
(
H..
)
.
(
l..
)
/
)
 
\{

  print 
"We matched '\$1' and '\$2'.
\textbackslash n
"
;


\}


a) We matched ‘Hel' and ‘ld'

b) We matched ‘Hel' and ‘lld'

c) We matched ‘Hel' and ‘lo ‘

d) None of the mentioned
\vspace{5pt} \par
\begin{tcolorbox}[colback=green!5, colframe=green!40!black, title=Solution]
Answer: c

Explanation: () groups a series of pattern element to a single element.

When we use pattern in parenthesis, we can use any of ‘\$1', ‘\$2' later to refer to the previously matched pattern.
\end{tcolorbox}
\end{mdframed}
\newpage

\begin{mdframed}[linewidth=1pt, roundcorner=5pt, linecolor=gray!30, backgroundcolor=gray!5]
\textbf{Question 9} 

Given segment of code:


\$string1 
=
 
"Hello
\textbackslash n
World
\textbackslash n
"
;


if
 
(
\$string1 
=
~ m
/
d\textbackslash n\textbackslash z
/
)
 
\{

  print 
"\$string1 is a string "
;

  print 
"that ends with 'd
\textbackslash \textbackslash 
n'.
\textbackslash n
"
;


\}


What does the symbol /z does?

a) changes line

b) matches the beginning of a string

c) matches the end of a string

d) none of the mentioned
\vspace{5pt} \par
\begin{tcolorbox}[colback=green!5, colframe=green!40!black, title=Solution]
Answer: c

Explanation: It matches the end of a string and not an internal line.The given segment of code outputs:

Hello

World

 is a string that ends with ‘d\textbackslash n'
\end{tcolorbox}
\end{mdframed}
\newpage

\begin{mdframed}[linewidth=1pt, roundcorner=5pt, linecolor=gray!30, backgroundcolor=gray!5]
\textbf{Question 10} 

Conversion of a regular expression into its corresponding NFA :

a) Thompson's Construction Algorithm

b) Powerset Construction

c) Kleene's algorithm

d) None of the mentioned
\vspace{5pt} \par
\begin{tcolorbox}[colback=green!5, colframe=green!40!black, title=Solution]
Answer: a

Explanation: Thompson construction algorithm is an algorithm in automata theory used to convert a given regular expression into NFA. Similarly, Kleene algorithm is used to convert a finite automaton to a regular expression.
\end{tcolorbox}
\end{mdframed}
\newpage

\newtopic{Finding Patterns in Text, Algebric Laws and Derivatives}

\begin{mdframed}[linewidth=1pt, roundcorner=5pt, linecolor=gray!30, backgroundcolor=gray!5]
\textbf{Question 1} 

The minimum length of a string \{0,1\}* not in the language corresponding to the given regular expression:

(0*+1*)(0*+1*)(0*+1*)

a) 3

b) 4

c) 5

d) 6
\vspace{5pt} \par
\begin{tcolorbox}[colback=green!5, colframe=green!40!black, title=Solution]
Answer: b

Explanation: 0101 or 1010 the strings with minimum length on \{0,1\}* which does not belong to the language of the given regular expression.Other strings like 111, 000, 1101, etc are accepted by the language .
\end{tcolorbox}
\end{mdframed}
\newpage

\begin{mdframed}[linewidth=1pt, roundcorner=5pt, linecolor=gray!30, backgroundcolor=gray!5]
\textbf{Question 2} 

Which of the following regular expression is equivalent to R(1,0)?

R(1,0)=\{111*\}*

a) (11+111)*

b) (111+1111)*

c) (111+11*)*

d) All of the mentioned
\vspace{5pt} \par
\begin{tcolorbox}[colback=green!5, colframe=green!40!black, title=Solution]
Answer: a

Explanation: What we observe from the question is that, it includes e and 11 and any number of 1's then. Therefore, its simplifies when we write the same reg. Expression as (11+111)*.
\end{tcolorbox}
\end{mdframed}
\newpage

\begin{mdframed}[linewidth=1pt, roundcorner=5pt, linecolor=gray!30, backgroundcolor=gray!5]
\textbf{Question 3} 

The minimum number of 1's to be used in a regular expression of the given language:

R(x): The language of all strings containing exactly 2 zeroes.

a) 2

b) 3

c) 0

d) 1
\vspace{5pt} \par
\begin{tcolorbox}[colback=green!5, colframe=green!40!black, title=Solution]
Answer: b

Explanation: It is not required to automate the question if asked theoretically.The number of zeroes fixed is 2. Therefore, we can represent the regular expression as 1*01*01*.
\end{tcolorbox}
\end{mdframed}
\newpage

\begin{mdframed}[linewidth=1pt, roundcorner=5pt, linecolor=gray!30, backgroundcolor=gray!5]
\textbf{Question 4} 

The given regular language corresponds to which of the given regular language

e+1+(1+0)*0+(0+1)*11

a) The language of all strings that end with 11 or 00

b) The language of all strings that end with 0 or 1

c) The language of all strings which does not end with 01

d) None of the mentioned
\vspace{5pt} \par
\begin{tcolorbox}[colback=green!5, colframe=green!40!black, title=Solution]
Answer: c

Explanation: According to the given regular expression, e is accepted by its language and it does not end with 00 or 11 or 0 or 1. Thus option a and b are eliminated. Further, the regular expression is valid for the third option.
\end{tcolorbox}
\end{mdframed}
\newpage

\begin{mdframed}[linewidth=1pt, roundcorner=5pt, linecolor=gray!30, backgroundcolor=gray!5]
\textbf{Question 5} 

Statement: If we take the union of two identical expression, we can replace them by one copy of the expression.

Which of the following is a correct option for the given statement?

a) Absorption Law

b) Idempotent Law

c) Closure Law

d) Commutative Law
\vspace{5pt} \par
\begin{tcolorbox}[colback=green!5, colframe=green!40!black, title=Solution]
Answer: b

Explanation: Idempotent Law states that if we take the union of two like expression, we can use a copy of the expression instead i.e. L+L=L. The common arithmetic operators are not idempotent.
\end{tcolorbox}
\end{mdframed}
\newpage

\begin{mdframed}[linewidth=1pt, roundcorner=5pt, linecolor=gray!30, backgroundcolor=gray!5]
\textbf{Question 6} 

Which among the following can be an annihilator for multiplication operation?

a) 0

b) 1

c) 100

d) 22/7
\vspace{5pt} \par
\begin{tcolorbox}[colback=green!5, colframe=green!40!black, title=Solution]
Answer: a

Explanation: An annihilator for an operator is a value such that when the operator is applied to the annihilator and some other value, the result is the annihilator.
\end{tcolorbox}
\end{mdframed}
\newpage

\begin{mdframed}[linewidth=1pt, roundcorner=5pt, linecolor=gray!30, backgroundcolor=gray!5]
\textbf{Question 7} 

Statement: A digit, when used in the CFG notation, will always be used as a terminal.

State true or false?

a) True

b) False
\vspace{5pt} \par
\begin{tcolorbox}[colback=green!5, colframe=green!40!black, title=Solution]
Answer: a

Explanation: Lowercase letters near the beginning of an alphabet, a, b and so on are terminal symbols. We shall also assume that digits and other characters such as + or parenthesis are terminals.
\end{tcolorbox}
\end{mdframed}
\newpage

\begin{mdframed}[linewidth=1pt, roundcorner=5pt, linecolor=gray!30, backgroundcolor=gray!5]
\textbf{Question 8} 

Choose the incorrect process to check whether the string belongs to the language of certain variable or not?

a) recursive inference

b) derivations

c) head to body method

d) All of the mentioned
\vspace{5pt} \par
\begin{tcolorbox}[colback=green!5, colframe=green!40!black, title=Solution]
Answer: d

Explanation: There are two approaches to infer that certain string are in the language of a certain variable. The most conventional way is to use the rules from body to head, recursive inference. The second approach is expanding the starting variable using one of its productions whose head is tart symbol and derive a string consisting entirely of terminals(head to body or derivations).
\end{tcolorbox}
\end{mdframed}
\newpage

\begin{mdframed}[linewidth=1pt, roundcorner=5pt, linecolor=gray!30, backgroundcolor=gray!5]
\textbf{Question 9} 

Statement: Left most derivations are lengthy as compared to Right most derivations.

Choose the correct option:

a) correct statement

b) incorrect statement

c) may or may not be correct

d) depends on the language of the grammar
\vspace{5pt} \par
\begin{tcolorbox}[colback=green!5, colframe=green!40!black, title=Solution]
Answer: c

Explanation: It completely depends on the person who develops the grammar of any language, how to make use of the tools i.e. leftmost and rightmost derivations.
\end{tcolorbox}
\end{mdframed}
\newpage

\begin{mdframed}[linewidth=1pt, roundcorner=5pt, linecolor=gray!30, backgroundcolor=gray!5]
\textbf{Question 10} 

A-$>$aAa|bAb|a|b|e

Which among the following is the correct option for the given production?

a) Left most derivation

b) Right most derivation

c) Recursive Inference

d) None of the mentioned
\vspace{5pt} \par
\begin{tcolorbox}[colback=green!5, colframe=green!40!black, title=Solution]
Answer: a

Explanation: The given form represents leftmost derivations in which at each step we replace the leftmost variable by one of its production bodies.


 
\begin{center}\includegraphics[width=0.5\textwidth]{automata-theory-questions-answers-algebric-laws-finding-patterns-derivatives-q10.png}\end{center}
\end{tcolorbox}
\end{mdframed}
\newpage

\newtopic{Properties-Non Regular Languages}

\begin{mdframed}[linewidth=1pt, roundcorner=5pt, linecolor=gray!30, backgroundcolor=gray!5]
\textbf{Question 1} 

All the regular languages can have one or more of the following descriptions:

i) DFA ii) NFA iii) e-NFA iv) Regular Expressions

Which of the following are correct?

a) i, ii, iv

b) i, ii, iii

c) i, iv

d) i, ii, iii, iv
\vspace{5pt} \par
\begin{tcolorbox}[colback=green!5, colframe=green!40!black, title=Solution]
Answer: d

Explanation: The class of languages known as the regular language has atleast four different descriptions: i) DFA ii) NFA iii) e-NFA iv) Regular Expressions
\end{tcolorbox}
\end{mdframed}
\newpage

\begin{mdframed}[linewidth=1pt, roundcorner=5pt, linecolor=gray!30, backgroundcolor=gray!5]
\textbf{Question 2} 

Which of the technique can be used to prove that a language is non regular?

a) Ardens theorem

b) Pumping Lemma

c) Ogden's Lemma

d) None of the mentioned
\vspace{5pt} \par
\begin{tcolorbox}[colback=green!5, colframe=green!40!black, title=Solution]
Answer: b

Explanation: We use the powerful technique called Pumping Lemma, for showing certain languages not to be regular. We use Ardens theorem to find out a regular expression out of a finite automaton.
\end{tcolorbox}
\end{mdframed}
\newpage

\begin{mdframed}[linewidth=1pt, roundcorner=5pt, linecolor=gray!30, backgroundcolor=gray!5]
\textbf{Question 3} 

Which of the following language regular?

a) \{a
i
b
i
|i$>$=0\}

b) \{a
i
b
i
|0$<$i$<$5\}

c) \{a
i
b
i
|i$>$=1\}

d) None of the mentioned
\vspace{5pt} \par
\begin{tcolorbox}[colback=green!5, colframe=green!40!black, title=Solution]
Answer: b

Explanation: Here, i has limits i.e. the language is finite, contains few elements and can be graphed using a deterministic finite automata. Thus, it is regular. Others can be proved non regular using Pumping lemma.
\end{tcolorbox}
\end{mdframed}
\newpage

\begin{mdframed}[linewidth=1pt, roundcorner=5pt, linecolor=gray!30, backgroundcolor=gray!5]
\textbf{Question 4} 

Which of the following are non regular?

a) The set of strings in \{a,b\}* with an even number of b's

b) The set of strings in \{a, b, c\}* where there is no c anywhere to the left of a

c) The set of strings in \{0, 1\}* that encode, in binary, an integer w that is a multiple of 3. Interpret the empty strings e as the number 0

d) None of the mentioned
\vspace{5pt} \par
\begin{tcolorbox}[colback=green!5, colframe=green!40!black, title=Solution]
Answer: d

Explanation: All of the given languages are regular and finite and thus, can be represented using respective deterministic finite automata. We can also use mealy or moore machine to represent remainders for option c.
\end{tcolorbox}
\end{mdframed}
\newpage

\begin{mdframed}[linewidth=1pt, roundcorner=5pt, linecolor=gray!30, backgroundcolor=gray!5]
\textbf{Question 5} 

If L is DFA-regular, L' is

a) Non regular

b) DFA-regular

c) Non-finite

d) None of the mentioned
\vspace{5pt} \par
\begin{tcolorbox}[colback=green!5, colframe=green!40!black, title=Solution]
Answer: b

Explanation: This is a simple example of a closure property: a property saying that the set of DFA-regular languages is closed under certain operations.
\end{tcolorbox}
\end{mdframed}
\newpage

\begin{mdframed}[linewidth=1pt, roundcorner=5pt, linecolor=gray!30, backgroundcolor=gray!5]
\textbf{Question 6} 

Which of the following options is incorrect?

a) A language L is regular if and only if ~L has finite number of equivalent classes

b) Let L be a regular language. If ~L has k equivalent classes, then any DFA that recognizes L must have atmost k states

c) A language L is NFA-regular if and only if it is DFA-regular

d) None of the mentioned
\vspace{5pt} \par
\begin{tcolorbox}[colback=green!5, colframe=green!40!black, title=Solution]
Answer: b

Explanation: Let L be a regular language. If ~L has k equivalent classes, then any DFA that recognizes L must have atleast k states.
\end{tcolorbox}
\end{mdframed}
\newpage

\begin{mdframed}[linewidth=1pt, roundcorner=5pt, linecolor=gray!30, backgroundcolor=gray!5]
\textbf{Question 7} 

Myphill Nerode  does the following:

a) Minimization of DFA

b) Tells us exactly when a language is regular

c) Minimization of DFA and tells us exactly when a language is regular

d) None of the mentioned
\vspace{5pt} \par
\begin{tcolorbox}[colback=green!5, colframe=green!40!black, title=Solution]
Answer: c

Explanation: In automata theory, the Myphill Nerode theorem provides a necessary and sufficient condition for a language to be regular. The Myphill Nerode theorem can be used to show a language L is regular by proving that the number of equivalence classes of R
L
(relation) is finite.
\end{tcolorbox}
\end{mdframed}
\newpage

\begin{mdframed}[linewidth=1pt, roundcorner=5pt, linecolor=gray!30, backgroundcolor=gray!5]
\textbf{Question 8} 

Which of the following are related to tree automaton?

a) Myphill Nerode Theorem

b) State machine

c) Courcelle's Theorem

d) All of the mentioned
\vspace{5pt} \par
\begin{tcolorbox}[colback=green!5, colframe=green!40!black, title=Solution]
Answer: d

Explanation: The myphill nerode theorem can be generalized to trees and an application of tree automata prove an algorithmic meta theorem about graphs.
\end{tcolorbox}
\end{mdframed}
\newpage

\begin{mdframed}[linewidth=1pt, roundcorner=5pt, linecolor=gray!30, backgroundcolor=gray!5]
\textbf{Question 9} 

Given languages:

i) \{a
n
b
n
|n$>$=0\}

ii) $<$div$>$
n
$<$/div$>$
n

iii) \{w$\in$\{a,b\}$\ast$| \#a(w)=\#b(w)\}, \# represents occurrences

Which of the following is/are non regular?

a) i, iii

b) i

c) iii

d) i, ii, iii
\vspace{5pt} \par
\begin{tcolorbox}[colback=green!5, colframe=green!40!black, title=Solution]
Answer: d

Explanation: There is no regular expression that can parse HTML documents. Other options are also non-regular as they cannot be drawn into finite automaton.
\end{tcolorbox}
\end{mdframed}
\newpage

\begin{mdframed}[linewidth=1pt, roundcorner=5pt, linecolor=gray!30, backgroundcolor=gray!5]
\textbf{Question 10} 

Finite state machine are not able to recognize Palindromes because:

a) Finite automata cannot deterministically find the midpoint

b) Finite automata cannot remember arbitarily large amount of data

c) Even if the mid point is known, it cannot find whether the second half matches the first

d) All of the mentioned
\vspace{5pt} \par
\begin{tcolorbox}[colback=green!5, colframe=green!40!black, title=Solution]
Answer: d

Explanation: It is the disadvantage or lack of property of a DFA that it cannot remember an arbitrarily such large amount of data which makes it incapable of accepting such languages like palindrome, reversal, etc.
\end{tcolorbox}
\end{mdframed}
\newpage

\newtopic{Pumping Lemma for Regular Language}

\begin{mdframed}[linewidth=1pt, roundcorner=5pt, linecolor=gray!30, backgroundcolor=gray!5]
\textbf{Question 1} 

Relate the following statement:

Statement: All sufficiently long words in a regular language can have a middle section of words repeated a number of times to produce a new word which also lies within the same language.

a) Turing Machine

b) Pumping Lemma

c) Arden's theorem

d) None of the mentioned
\vspace{5pt} \par
\begin{tcolorbox}[colback=green!5, colframe=green!40!black, title=Solution]
Answer: b

Explanation: Pumping lemma defines an essential property for every regular language in automata theory. It has certain rules which decide whether a language is regular or not.
\end{tcolorbox}
\end{mdframed}
\newpage

\begin{mdframed}[linewidth=1pt, roundcorner=5pt, linecolor=gray!30, backgroundcolor=gray!5]
\textbf{Question 2} 

While applying Pumping lemma over a language, we consider a string w that belong to L and fragment it into \_\_\_\_\_\_\_\_\_ parts.

a) 2

b) 5

c) 3

d) 6
\vspace{5pt} \par
\begin{tcolorbox}[colback=green!5, colframe=green!40!black, title=Solution]
Answer: c

Explanation: We select a string w such that w=xyz and |y|$>$0 and other conditions. However, there exists an integer n such that |w|$>$=n for any wÎL.
\end{tcolorbox}
\end{mdframed}
\newpage

\begin{mdframed}[linewidth=1pt, roundcorner=5pt, linecolor=gray!30, backgroundcolor=gray!5]
\textbf{Question 3} 

If we select a string w such that w$\in$L, and w=xyz. Which of the following portions cannot be an empty string?

a) x

b) y

c) z

d) all of the mentioned
\vspace{5pt} \par
\begin{tcolorbox}[colback=green!5, colframe=green!40!black, title=Solution]
Answer: b

Explanation: The lemma says, the portion y in xyz cannot be zero or empty i.e. |y|$>$0, this condition needs to be fulfilled to check the conclusion condition.
\end{tcolorbox}
\end{mdframed}
\newpage

\begin{mdframed}[linewidth=1pt, roundcorner=5pt, linecolor=gray!30, backgroundcolor=gray!5]
\textbf{Question 4} 

Let w= xyz and y refers to the middle portion and |y|$>$0.What do we call the process of repeating y 0 or more times before checking that they still belong to the language L or not?

a) Generating

b) Pumping

c) Producing

d) None of the mentioned
\vspace{5pt} \par
\begin{tcolorbox}[colback=green!5, colframe=green!40!black, title=Solution]
Answer: b

Explanation: The process of repeatation is called pumping and so, pumping is the process we perform before we check whether the pumped string belongs to L or not.
\end{tcolorbox}
\end{mdframed}
\newpage

\begin{mdframed}[linewidth=1pt, roundcorner=5pt, linecolor=gray!30, backgroundcolor=gray!5]
\textbf{Question 5} 

There exists a language L. We define a string w such that w$\in$L and w=xyz and |w| $>$=n for some constant integer n.What can be the maximum length of the substring xy i.e. |xy|$<$=?

a) n

b) |y|

c) |x|

d) none of the mentioned
\vspace{5pt} \par
\begin{tcolorbox}[colback=green!5, colframe=green!40!black, title=Solution]
Answer: a

Explanation: It is the first conditional statement of the lemma that states that |xy|$<$=n, i.e. the maximum length of the substring xy in w can be n only.
\end{tcolorbox}
\end{mdframed}
\newpage

\begin{mdframed}[linewidth=1pt, roundcorner=5pt, linecolor=gray!30, backgroundcolor=gray!5]
\textbf{Question 6} 

Fill in the blank in terms of p, where p is the maximum string length in L.

Statement: Finite languages trivially satisfy the pumping lemma by having n = \_\_\_\_\_\_

a) p*1

b) p+1

c) p-1

d) None of the mentioned
\vspace{5pt} \par
\begin{tcolorbox}[colback=green!5, colframe=green!40!black, title=Solution]
Answer: b

Explanation: Finite languages trivially satisfy the pumping lemma by having n equal to the maximum string length in l plus 1.
\end{tcolorbox}
\end{mdframed}
\newpage

\begin{mdframed}[linewidth=1pt, roundcorner=5pt, linecolor=gray!30, backgroundcolor=gray!5]
\textbf{Question 7} 

Answer in accordance to the third and last statement in pumping lemma:

For all \_\_\_\_\_\_\_ xy
i
z $\in$L

a) i$>$0

b) i$<$0

c) i$<$=0

d) i$>$=0
\vspace{5pt} \par
\begin{tcolorbox}[colback=green!5, colframe=green!40!black, title=Solution]
Answer: d

Explanation: Suppose L is a regular language . Then there is an integer n so that for any x$\in$L and |x|$>$=n, there are strings u,v,w so that

x= uvw

|uv|$<$=n

|v|$>$0

for any m$>$=0, uv
m
w $\in$L.
\end{tcolorbox}
\end{mdframed}
\newpage

\begin{mdframed}[linewidth=1pt, roundcorner=5pt, linecolor=gray!30, backgroundcolor=gray!5]
\textbf{Question 8} 

If d is a final state, which of the following is correct according to the given diagram?


 
\begin{center}\includegraphics[width=0.5\textwidth]{automata-theory-questions-answers-pumping-lemma-regular-language-q8.png}\end{center}
 

a) x=p, y=qr, z=s

b) x=p, z=qrs

c) x=pr, y=r, z=s

d) All of the mentioned
\vspace{5pt} \par
\begin{tcolorbox}[colback=green!5, colframe=green!40!black, title=Solution]
Answer: a

Explanation: The FSA accepts the string pqrs. In terms of pumping lemma, the string pqrs is broken into an x portion an a, a y portion qr and a z portion s.
\end{tcolorbox}
\end{mdframed}
\newpage

\begin{mdframed}[linewidth=1pt, roundcorner=5pt, linecolor=gray!30, backgroundcolor=gray!5]
\textbf{Question 9} 

Let w be a string and fragmented by three variable x, y, and z as per pumping lemma. What does these variables represent?

a) string count

b) string

c) string count and string

d) none of  the mentioned
\vspace{5pt} \par
\begin{tcolorbox}[colback=green!5, colframe=green!40!black, title=Solution]
Answer: a

Explanation: Given: w =xyz. Here, xyz individually represents strings or rather substrings which we  compute over conditions to check the regularity of the language.
\end{tcolorbox}
\end{mdframed}
\newpage

\begin{mdframed}[linewidth=1pt, roundcorner=5pt, linecolor=gray!30, backgroundcolor=gray!5]
\textbf{Question 10} 

Which of the following one can relate to the given statement:

Statement: If n items are put into m containers, with n$>$m, then atleast one container must contain more than one item.

a) Pumping lemma

b) Pigeon Hole principle

c) Count principle

d) None of the mentioned
\vspace{5pt} \par
\begin{tcolorbox}[colback=green!5, colframe=green!40!black, title=Solution]
Answer: b

Explanation: Pigeon hole principle states the following example: If there exists n=10 pigeons in m=9 holes, then since 10$>$9, the pigeonhole principle says that at least one hole has more than one pigeon.
\end{tcolorbox}
\end{mdframed}
\newpage

\newtopic{Applications of Pumping Lemma}

\begin{mdframed}[linewidth=1pt, roundcorner=5pt, linecolor=gray!30, backgroundcolor=gray!5]
\textbf{Question 1} 

Which kind of proof is used to prove the regularity of a language?

a) Proof by contradiction

b) Direct proof

c) Proof by induction

d) None of the mentioned
\vspace{5pt} \par
\begin{tcolorbox}[colback=green!5, colframe=green!40!black, title=Solution]
Answer: a

Explanation: We use the method of proof by contradiction in pumping lemma to prove that a language is regular or not.
\end{tcolorbox}
\end{mdframed}
\newpage

\begin{mdframed}[linewidth=1pt, roundcorner=5pt, linecolor=gray!30, backgroundcolor=gray!5]
\textbf{Question 2} 

The language of balanced paranthesis is

a) regular

b) non regular

c) may be regular

d) none of the mentioned
\vspace{5pt} \par
\begin{tcolorbox}[colback=green!5, colframe=green!40!black, title=Solution]
Answer: b

Explanation: Given n, there is a string of balanced parentheses that begins with more than p left parentheses, so that y will contain entirely of left parentheses. By repeating y, we can produce a string that does not contain the same number of left and right parentheses, and so they cannot be balanced.
\end{tcolorbox}
\end{mdframed}
\newpage

\begin{mdframed}[linewidth=1pt, roundcorner=5pt, linecolor=gray!30, backgroundcolor=gray!5]
\textbf{Question 3} 

State true or false:

Statement: Pumping lemma gives a necessary but not sufficient condition for a language to be regular.

a) true

b) false
\vspace{5pt} \par
\begin{tcolorbox}[colback=green!5, colframe=green!40!black, title=Solution]
Answer: a

Explanation: The converse of the lemma is not true. There may exists some language which satisfy all the conditions of the lemma and still be non-regular.
\end{tcolorbox}
\end{mdframed}
\newpage

\begin{mdframed}[linewidth=1pt, roundcorner=5pt, linecolor=gray!30, backgroundcolor=gray!5]
\textbf{Question 4} 

Which of the following is/are an example of pigeon hole principle?

a) Softball team

b) Sock picking

c) Hair counting

d) All of the mentioned
\vspace{5pt} \par
\begin{tcolorbox}[colback=green!5, colframe=green!40!black, title=Solution]
Answer: d

Explanation: There are several applications of pigeonhole principle:

Example: The softball team: Suppose 7 people who want to play softball(n=7 items), with a limitation of only 4 softball teams to choose from. The pigeonhole principle tells us that they cannot all play for different teams; there must be atleast one team featuring atleast two of the seven players.
\end{tcolorbox}
\end{mdframed}
\newpage

\begin{mdframed}[linewidth=1pt, roundcorner=5pt, linecolor=gray!30, backgroundcolor=gray!5]
\textbf{Question 5} 

Pigeonhole principle can be applied in the following computer science algorithms:

a) hashing algorithm

b) lossless compression algorithm

c) hashing algorithm and lossless compression algorithm

d) none of the mentioned
\vspace{5pt} \par
\begin{tcolorbox}[colback=green!5, colframe=green!40!black, title=Solution]
Answer: c

Explanation: Collisions are inevitable in a hash table because the number of possible keys exceeds the number of indices in the array.
\end{tcolorbox}
\end{mdframed}
\newpage

\begin{mdframed}[linewidth=1pt, roundcorner=5pt, linecolor=gray!30, backgroundcolor=gray!5]
\textbf{Question 6} 

If n objects are distributed over m places, and n $<$ m, then some of the places receive:

a) at least 2 objects

b) at most 2 objects

c) no object

d) none of the mentioned
\vspace{5pt} \par
\begin{tcolorbox}[colback=green!5, colframe=green!40!black, title=Solution]
Answer: c

Explanation: This is one of the alternative formulations of the pigeon hole principle. As n $<$ m, there will exist some place which will not receive any of the object.
\end{tcolorbox}
\end{mdframed}
\newpage

\begin{mdframed}[linewidth=1pt, roundcorner=5pt, linecolor=gray!30, backgroundcolor=gray!5]
\textbf{Question 7} 

Which of the following fields may have pigeonhole principle violated?

a) Discrete mathematics

b) Computer Science

c) Quantum Mechanics

d) None of the mentioned
\vspace{5pt} \par
\begin{tcolorbox}[colback=green!5, colframe=green!40!black, title=Solution]
Answer: c

Explanation: Y Aharonov proved mathematically the violation of pigeon hole principle in Quantum mechanics and proposed inferometric experiments to test it.
\end{tcolorbox}
\end{mdframed}
\newpage

\begin{mdframed}[linewidth=1pt, roundcorner=5pt, linecolor=gray!30, backgroundcolor=gray!5]
\textbf{Question 8} 

Which of the following is not an application of Pumping Lemma?

a) \{0
i
1
i
|i$>$=0\}

b) \{0
i
x|i$>$=0, x$\in$\{0, 1\}* and |x|$<$=i\}

c) \{0
n
| n is prime\}

d) None of the mentioned
\vspace{5pt} \par
\begin{tcolorbox}[colback=green!5, colframe=green!40!black, title=Solution]
Answer: d

Explanation: None of the mentioned are regular language and are an application to the technique Pumping Lemma. Each one of the mentioned can be proved non regular using the steps in Pumping lemma.
\end{tcolorbox}
\end{mdframed}
\newpage

\begin{mdframed}[linewidth=1pt, roundcorner=5pt, linecolor=gray!30, backgroundcolor=gray!5]
\textbf{Question 9} 

Which of the following can refer a language to be non regular?

a) Pumping Lemma

b) Myphill Nerode

c) Pumping Lemma and Myphill Nerode

d) None of the mentioned
\vspace{5pt} \par
\begin{tcolorbox}[colback=green!5, colframe=green!40!black, title=Solution]
Answer: c

Explanation: On the contrary, the typical way to prove that a language is to construct either a finite state machine or a regular expression for the language.
\end{tcolorbox}
\end{mdframed}
\newpage

\begin{mdframed}[linewidth=1pt, roundcorner=5pt, linecolor=gray!30, backgroundcolor=gray!5]
\textbf{Question 10} 

Which of the following is not an example of counting argument?

a) Pigeonhole principle

b) Dirichlet's drawer principle

c) Dirichlet's box principle

d) None of the mentioned
\vspace{5pt} \par
\begin{tcolorbox}[colback=green!5, colframe=green!40!black, title=Solution]
Answer: d

Explanation: Pigeon hole principle or Dirichlet's drawer principle or Dirichlet's box principle is an example of counting argument whose field is called Combinatorics.
\end{tcolorbox}
\end{mdframed}
\newpage

\newtopic{Closure Properties under Boolean Operations}

\begin{mdframed}[linewidth=1pt, roundcorner=5pt, linecolor=gray!30, backgroundcolor=gray!5]
\textbf{Question 1} 

If L1, L2 are regular and op(L1, L2) is also regular, then L1 and L2 are said to be \_\_\_\_\_\_\_\_\_\_\_\_ under an operation op.

a) open

b) closed

c) decidable

d) none of the mentioned
\vspace{5pt} \par
\begin{tcolorbox}[colback=green!5, colframe=green!40!black, title=Solution]
Answer: b

Explanation: If two regular languages are closed under an operation op, then the resultant of the languages over an operation op will also be regular.
\end{tcolorbox}
\end{mdframed}
\newpage

\begin{mdframed}[linewidth=1pt, roundcorner=5pt, linecolor=gray!30, backgroundcolor=gray!5]
\textbf{Question 2} 

Suppose a regular language L is closed under the operation halving, then the result would be:

a) 1/4 L will be regular

b) 1/2 L will be regular

c) 1/8 L will be regular

d) All of the mentioned
\vspace{5pt} \par
\begin{tcolorbox}[colback=green!5, colframe=green!40!black, title=Solution]
Answer: d

Explanation: At first stage 1/2 L will be regular and subsequently, all the options will be regular.
\end{tcolorbox}
\end{mdframed}
\newpage

\begin{mdframed}[linewidth=1pt, roundcorner=5pt, linecolor=gray!30, backgroundcolor=gray!5]
\textbf{Question 3} 

If L1$^\prime$ and L2$^\prime$ are regular languages, then L1.L2 will be

a) regular

b) non regular

c) may be regular

d) none of the mentioned
\vspace{5pt} \par
\begin{tcolorbox}[colback=green!5, colframe=green!40!black, title=Solution]
Answer: a

Explanation: Regular language is closed under complement operation. Thus, if L1$^\prime$ and L2$^\prime$ are regular so are L1 and L2. And if L1 and L2 are regular so is L1.L2.
\end{tcolorbox}
\end{mdframed}
\newpage

\begin{mdframed}[linewidth=1pt, roundcorner=5pt, linecolor=gray!30, backgroundcolor=gray!5]
\textbf{Question 4} 

If L1 and L2$^\prime$ are regular languages, L1 $\cap$ (L2$^\prime$ U L1$^\prime$)' will be

a) regular

b) non regular

c) may be regular

d) none of the mentioned
\vspace{5pt} \par
\begin{tcolorbox}[colback=green!5, colframe=green!40!black, title=Solution]
Answer: a

Explanation: If L1 is regular, so is L1$^\prime$ and if L1$^\prime$ and L2$^\prime$ are regular so is L1$^\prime$ U L2$^\prime$. Further, regular  languages are also closed under intersection operation.
\end{tcolorbox}
\end{mdframed}
\newpage

\begin{mdframed}[linewidth=1pt, roundcorner=5pt, linecolor=gray!30, backgroundcolor=gray!5]
\textbf{Question 5} 

If A and B are regular languages, !(A' U B') is:

a) regular

b) non regular

c) may be regular

d) none of the mentioned
\vspace{5pt} \par
\begin{tcolorbox}[colback=green!5, colframe=green!40!black, title=Solution]
Answer: a

Explanation: If A and B are regular languages, then A Ç B is a regular language and  A $\cap$ B is equivalent to  !(A' U B').
\end{tcolorbox}
\end{mdframed}
\newpage

\begin{mdframed}[linewidth=1pt, roundcorner=5pt, linecolor=gray!30, backgroundcolor=gray!5]
\textbf{Question 6} 

Which among the following are the boolean operations that under which regular languages are closed?

a) Union

b) Intersection

c) Complement

d) All of the mentioned
\vspace{5pt} \par
\begin{tcolorbox}[colback=green!5, colframe=green!40!black, title=Solution]
Answer: d

Explanation: Regular languages are closed under the following operations:

a) Regular expression operations

b) Boolean operations

c) Homomorphism

d) Inverse Homomorphism
\end{tcolorbox}
\end{mdframed}
\newpage

\begin{mdframed}[linewidth=1pt, roundcorner=5pt, linecolor=gray!30, backgroundcolor=gray!5]
\textbf{Question 7} 

Suppose a language L1 has 2 states and L2 has 2 states. After using the cross product construction method,	we have a machine M that accepts L1 $\cap$ L2. The total number of states in M:

a) 6

b) 4

c) 2

d) 8
\vspace{5pt} \par
\begin{tcolorbox}[colback=green!5, colframe=green!40!black, title=Solution]
Answer: b

Explanation: M is defined as: (Q, S, d, q0, F)

where Q=Q1*Q2 and F=F1*F2.
\end{tcolorbox}
\end{mdframed}
\newpage

\begin{mdframed}[linewidth=1pt, roundcorner=5pt, linecolor=gray!30, backgroundcolor=gray!5]
\textbf{Question 8} 

If L is a regular language, then (L')' U L will be :

a) L

b) L'

c) f

d) none of the mentioned
\vspace{5pt} \par
\begin{tcolorbox}[colback=green!5, colframe=green!40!black, title=Solution]
Answer: a

Explanation: (L')' is equivalent to L and L U L is subsequently equivalent to L.
\end{tcolorbox}
\end{mdframed}
\newpage

\begin{mdframed}[linewidth=1pt, roundcorner=5pt, linecolor=gray!30, backgroundcolor=gray!5]
\textbf{Question 9} 

If L is a regular language, then (((L')r)')* is:

a) regular

b) non regular

c) may be regular

d) none of the mentioned
\vspace{5pt} \par
\begin{tcolorbox}[colback=green!5, colframe=green!40!black, title=Solution]
Answer: a

Explanation: If L is regular so is its complement, if L' is regular so is its reverse, if (L')
r
 is regular so is its Kleene.
\end{tcolorbox}
\end{mdframed}
\newpage

\begin{mdframed}[linewidth=1pt, roundcorner=5pt, linecolor=gray!30, backgroundcolor=gray!5]
\textbf{Question 10} 

Which among the following is the closure property of a regular language?

a) Emptiness

b) Universality

c) Membership

d) None of the mentioned
\vspace{5pt} \par
\begin{tcolorbox}[colback=green!5, colframe=green!40!black, title=Solution]
Answer: d

Explanation: All the following mentioned are decidability properties of a regular language. The closure properties of a regular language include union, concatenation, intersection, Kleene, complement , reverse and many more operations.
\end{tcolorbox}
\end{mdframed}
\newpage

\newtopic{Context Free Grammar-Derivations and Definitions}

\begin{mdframed}[linewidth=1pt, roundcorner=5pt, linecolor=gray!30, backgroundcolor=gray!5]
\textbf{Question 1} 

The entity which generate Language is termed as:

a) Automata

b) Tokens

c) Grammar

d) Data
\vspace{5pt} \par
\begin{tcolorbox}[colback=green!5, colframe=green!40!black, title=Solution]
Answer: c

Explanation: The entity which accepts a language is termed as Automata while the one which generates it is called Grammar. Tokens are the smallest individual unit of a program.
\end{tcolorbox}
\end{mdframed}
\newpage

\begin{mdframed}[linewidth=1pt, roundcorner=5pt, linecolor=gray!30, backgroundcolor=gray!5]
\textbf{Question 2} 

Production Rule: aAb-$>$agb belongs to which of the following category?

a)  Regular Language

b) Context free Language

c) Context Sensitive Language

d) Recursively Ennumerable Language
\vspace{5pt} \par
\begin{tcolorbox}[colback=green!5, colframe=green!40!black, title=Solution]
Answer: c

Explanation: Context Sensitive Language or Type 1 or Linearly Bounded Non deterministic Language has the production rule where the production is context dependent i.e.  aAb-$>$agb.
\end{tcolorbox}
\end{mdframed}
\newpage

\begin{mdframed}[linewidth=1pt, roundcorner=5pt, linecolor=gray!30, backgroundcolor=gray!5]
\textbf{Question 3} 

Which of the following statement is false?

a) Context free language is the subset of context sensitive language

b) Regular language is the subset of context sensitive language

c) Recursively ennumerable language is the super set of regular language

d) Context sensitive language is a subset of context free language
\vspace{5pt} \par
\begin{tcolorbox}[colback=green!5, colframe=green!40!black, title=Solution]
Answer: d

Explanation: Every regular language can be produced by context free grammar and context free language can be produced by context sensitive grammar and so on.


 
\begin{center}\includegraphics[width=0.5\textwidth]{automata-theory-questions-answers-context-free-grammar-derivations-definitions-q3.png}\end{center}
\end{tcolorbox}
\end{mdframed}
\newpage

\begin{mdframed}[linewidth=1pt, roundcorner=5pt, linecolor=gray!30, backgroundcolor=gray!5]
\textbf{Question 4} 

The Grammar can be defined as: G=(V, $\Sigma$, p, S)

In the given definition, what does S represents?

a) Accepting State

b) Starting Variable

c) Sensitive Grammar

d) None of these
\vspace{5pt} \par
\begin{tcolorbox}[colback=green!5, colframe=green!40!black, title=Solution]
Answer: b

Explanation: G=(V, $\Sigma$, p, S), here V=Finite set of variables,  $\Sigma$= set of terminals, p= finite productions, S= Starting Variable.
\end{tcolorbox}
\end{mdframed}
\newpage

\begin{mdframed}[linewidth=1pt, roundcorner=5pt, linecolor=gray!30, backgroundcolor=gray!5]
\textbf{Question 5} 

Which among the following cannot be accepted by  a regular grammar?

a) L is a set of numbers divisible by 2

b) L is a set of binary complement

c) L is a set of string with odd number of 0

d) L is a set of 0
n
1
n
\vspace{5pt} \par
\begin{tcolorbox}[colback=green!5, colframe=green!40!black, title=Solution]
Answer: d

Explanation: There exists no finite automata to accept the given language i.e. 0
n
1
n
. For other options, it is possible to make a dfa or nfa representing the language set.
\end{tcolorbox}
\end{mdframed}
\newpage

\begin{mdframed}[linewidth=1pt, roundcorner=5pt, linecolor=gray!30, backgroundcolor=gray!5]
\textbf{Question 6} 

Which of the expression is appropriate?

For production p: a-$>$b where a$\in$V and b$\in$\_\_\_\_\_\_\_

a) V

b) S

c) (V+$\Sigma$)*

d) V+ $\Sigma$
\vspace{5pt} \par
\begin{tcolorbox}[colback=green!5, colframe=green!40!black, title=Solution]
Answer: c

Explanation: According to the definition, the starting variable can produce another variable or any terminal or a variable which leads to terminal.
\end{tcolorbox}
\end{mdframed}
\newpage

\begin{mdframed}[linewidth=1pt, roundcorner=5pt, linecolor=gray!30, backgroundcolor=gray!5]
\textbf{Question 7} 

For S-$>$0S1|e for $\Sigma$=\{0,1\}*, which of the following is wrong for the language produced?

a) Non regular language

b) 0
n
1
n
 | n$>$=0

c) 0
n
1
n
 | n$>$=1

d) None of the mentioned
\vspace{5pt} \par
\begin{tcolorbox}[colback=green!5, colframe=green!40!black, title=Solution]
Answer: d

Explanation: L=\{e, 01, 0011, 000111, ……0
n
1
n
 \}. As epsilon is a part of the set, thus all the options are correct implying none of them to be wrong.
\end{tcolorbox}
\end{mdframed}
\newpage

\begin{mdframed}[linewidth=1pt, roundcorner=5pt, linecolor=gray!30, backgroundcolor=gray!5]
\textbf{Question 8} 

The minimum number of productions required to produce a language consisting of palindrome strings over   $\Sigma$=\{a,b\} is

a) 3

b) 7

c) 5

d) 6
\vspace{5pt} \par
\begin{tcolorbox}[colback=green!5, colframe=green!40!black, title=Solution]
Answer: c

Explanation: The grammar which produces a palindrome set can be written as:

S-$>$ aSa | bSb | e | a | b

L=\{e, a, b, aba, abbbaabbba…..\}
\end{tcolorbox}
\end{mdframed}
\newpage

\begin{mdframed}[linewidth=1pt, roundcorner=5pt, linecolor=gray!30, backgroundcolor=gray!5]
\textbf{Question 9} 

Which of the following statement is correct?

a) All Regular grammar are context free but not vice versa

b) All context free grammar are regular grammar but not vice versa

c) Regular grammar and context free grammar are the same entity

d) None of the mentioned
\vspace{5pt} \par
\begin{tcolorbox}[colback=green!5, colframe=green!40!black, title=Solution]
Answer: a

Explanation: Regular grammar is a subset of context free grammar and thus all regular grammars are context free.
\end{tcolorbox}
\end{mdframed}
\newpage

\begin{mdframed}[linewidth=1pt, roundcorner=5pt, linecolor=gray!30, backgroundcolor=gray!5]
\textbf{Question 10} 

Are ambiguous grammar context free?

a) Yes

b) No
\vspace{5pt} \par
\begin{tcolorbox}[colback=green!5, colframe=green!40!black, title=Solution]
Answer: a

Explanation: A context free grammar G is ambiguous if there is atleast one string in L(G) which has two or more distinct leftmost derivations.
\end{tcolorbox}
\end{mdframed}
\newpage

\newtopic{The Language of a Grammar, Inferences and Ambiguity}

\begin{mdframed}[linewidth=1pt, roundcorner=5pt, linecolor=gray!30, backgroundcolor=gray!5]
\textbf{Question 1} 

Which of the following is not a notion of Context free grammars?

a) Recursive Inference

b) Derivations

c) Sentential forms

d) All of the mentioned
\vspace{5pt} \par
\begin{tcolorbox}[colback=green!5, colframe=green!40!black, title=Solution]
Answer: d

Explanation: The following are the notions to express Context free grammars:

a) Recursive Inferences

b) Derivations

c) Sentential form

d) Parse trees
\end{tcolorbox}
\end{mdframed}
\newpage

\begin{mdframed}[linewidth=1pt, roundcorner=5pt, linecolor=gray!30, backgroundcolor=gray!5]
\textbf{Question 2} 

State true or false:

Statement: The recursive inference procedure determines that string w is in the language of the variable A, A being the starting variable.

a) true

b) false
\vspace{5pt} \par
\begin{tcolorbox}[colback=green!5, colframe=green!40!black, title=Solution]
Answer: a

Explanation: We apply the productions of CFG to infer that certain strings are in the language of a certain variable.
\end{tcolorbox}
\end{mdframed}
\newpage

\begin{mdframed}[linewidth=1pt, roundcorner=5pt, linecolor=gray!30, backgroundcolor=gray!5]
\textbf{Question 3} 

Which of the following is/are the suitable approaches for inferencing?

a) Recursive Inference

b) Derivations

c) Recursive Inference and Derivations

d) None of the mentioned
\vspace{5pt} \par
\begin{tcolorbox}[colback=green!5, colframe=green!40!black, title=Solution]
Answer: c

Explanation: Two inference approaches:
\end{tcolorbox}
\end{mdframed}
\newpage

\begin{mdframed}[linewidth=1pt, roundcorner=5pt, linecolor=gray!30, backgroundcolor=gray!5]
\textbf{Question 1} 

Recursive inference, using productions from body to head
\vspace{5pt} \par
\begin{tcolorbox}[colback=green!5, colframe=green!40!black, title=Solution]
Answer not found.
\end{tcolorbox}
\end{mdframed}
\newpage

\begin{mdframed}[linewidth=1pt, roundcorner=5pt, linecolor=gray!30, backgroundcolor=gray!5]
\textbf{Question 2} 

Derivations, using productions from head to body
\vspace{5pt} \par
\begin{tcolorbox}[colback=green!5, colframe=green!40!black, title=Solution]
Answer not found.
\end{tcolorbox}
\end{mdframed}
\newpage

\begin{mdframed}[linewidth=1pt, roundcorner=5pt, linecolor=gray!30, backgroundcolor=gray!5]
\textbf{Question 4} 

If w belongs to L(G), for some CFG, then w has a parse tree, which defines the syntactic structure of w. w could be:

a) program

b) SQL-query

c) XML document

d) All of the mentioned
\vspace{5pt} \par
\begin{tcolorbox}[colback=green!5, colframe=green!40!black, title=Solution]
Answer: d

Explanation: Parse trees are an alternative representation to derivations and recursive inferences. There can be several parse trees for the same string.
\end{tcolorbox}
\end{mdframed}
\newpage

\begin{mdframed}[linewidth=1pt, roundcorner=5pt, linecolor=gray!30, backgroundcolor=gray!5]
\textbf{Question 5} 

Is the following statement correct?

Statement: Recursive inference and derivation are equivalent.

a) Yes

b) No
\vspace{5pt} \par
\begin{tcolorbox}[colback=green!5, colframe=green!40!black, title=Solution]
Answer: a

Explanation: Yes, they are equivalent. Both the terminologies represent the two approaches of recursive inferencing.
\end{tcolorbox}
\end{mdframed}
\newpage

\begin{mdframed}[linewidth=1pt, roundcorner=5pt, linecolor=gray!30, backgroundcolor=gray!5]
\textbf{Question 6} 

A-$>$aA| a| b

The number of steps to form aab:

a) 2

b) 3

c) 4

d) 5
\vspace{5pt} \par
\begin{tcolorbox}[colback=green!5, colframe=green!40!black, title=Solution]
Answer: b

Explanation: A-$>$aA=$>$aaA=$>$aab
\end{tcolorbox}
\end{mdframed}
\newpage

\begin{mdframed}[linewidth=1pt, roundcorner=5pt, linecolor=gray!30, backgroundcolor=gray!5]
\textbf{Question 7} 

An expression is mentioned as follows. Figure out number of incorrect notations or symbols, such that a change in those could make the expression correct.

L(G)=\{w in T*|S→*w\}

a) 0 Errors

b) 1 Error

c) 2 Error

d) Invalid Expression
\vspace{5pt} \par
\begin{tcolorbox}[colback=green!5, colframe=green!40!black, title=Solution]
Answer: a

Explanation: For the given expression, L(G)=\{w in T*|S→*w\}, If G(V, T, P, S) is a CFG, the language of G, denoted by L(G), is the set of terminal strings that have derivations from the start symbol.
\end{tcolorbox}
\end{mdframed}
\newpage

\begin{mdframed}[linewidth=1pt, roundcorner=5pt, linecolor=gray!30, backgroundcolor=gray!5]
\textbf{Question 8} 

The language accepted by Push down Automaton:

a) Recursive Language

b) Context free language

c) Linearly Bounded language

d) All of the mentioned
\vspace{5pt} \par
\begin{tcolorbox}[colback=green!5, colframe=green!40!black, title=Solution]
Answer: b

Explanation: Push down automata accepts context free language.
\end{tcolorbox}
\end{mdframed}
\newpage

\begin{mdframed}[linewidth=1pt, roundcorner=5pt, linecolor=gray!30, backgroundcolor=gray!5]
\textbf{Question 9} 

Which among the following is the correct option for the given grammar?

G-$>$X111|G1,X-$>$X0|00

a) \{0
a
1
b
|a=2,b=3\}

b) \{0
a
1
b
|a=1,b=5\}

c) \{0
a
1
b
|a=b\}

d) More than one of the mentioned is correct
\vspace{5pt} \par
\begin{tcolorbox}[colback=green!5, colframe=green!40!black, title=Solution]
Answer: a

Explanation: Using the recursive approach, we can conclude that option a is the correct answer, and its not possible for a grammar to have more than one language.
\end{tcolorbox}
\end{mdframed}
\newpage

\begin{mdframed}[linewidth=1pt, roundcorner=5pt, linecolor=gray!30, backgroundcolor=gray!5]
\textbf{Question 10} 

Which of the following the given language belongs to?

L=\{a
m
b
m
c
m
| m$>$=1\}

a) Context free language

b) Regular language

c) Context free language \& Regular language

d) None of the mentioned
\vspace{5pt} \par
\begin{tcolorbox}[colback=green!5, colframe=green!40!black, title=Solution]
Answer: d

Explanation: The given language is neither accepted by a finite automata or a push down automata. Thus, it is neither a context free language nor a regular language.
\end{tcolorbox}
\end{mdframed}
\newpage

\begin{mdframed}[linewidth=1pt, roundcorner=5pt, linecolor=gray!30, backgroundcolor=gray!5]
\textbf{Question 11} 

Choose the correct option:

Statement: There exists two inference approaches: Recursive Inference \& Derivation

a) true

b) partially true

c) false

d) none of the mentioned
\vspace{5pt} \par
\begin{tcolorbox}[colback=green!5, colframe=green!40!black, title=Solution]
Answer: a

Explanation: We apply the productions of a CFG to infer that certain strings are in a language of certain variable.
\end{tcolorbox}
\end{mdframed}
\newpage

\begin{mdframed}[linewidth=1pt, roundcorner=5pt, linecolor=gray!30, backgroundcolor=gray!5]
\textbf{Question 12} 

Choose the correct option:

Statement 1: Recursive Inference, using productions from head to body.

Statement 2: Derivations, using productions from body to head.

a) Statement 1 is true and Statement 2 is true

b) Statement 1 and Statement 2, both are false

c) Statement 1 is true and Statement 2 is false

d) Statement 2 is true and Statement 1 is true
\vspace{5pt} \par
\begin{tcolorbox}[colback=green!5, colframe=green!40!black, title=Solution]
Answer: b

Explanation: Both statements are false. Recursive Inference, using productions from body to head. Derivations, using productions from head to body.
\end{tcolorbox}
\end{mdframed}
\newpage

\begin{mdframed}[linewidth=1pt, roundcorner=5pt, linecolor=gray!30, backgroundcolor=gray!5]
\textbf{Question 13} 

Which of the following statements are correct for a concept called inherent ambiguity in CFL?

a) Every CFG for L is ambiguous

b) Every CFG for L is unambiguous

c) Every CFG is also regular

d) None of the mentioned
\vspace{5pt} \par
\begin{tcolorbox}[colback=green!5, colframe=green!40!black, title=Solution]
Answer: a

Explanation: A CFL L is said to be inherently ambiguous if every CFG for L is ambiguous.
\end{tcolorbox}
\end{mdframed}
\newpage

\begin{mdframed}[linewidth=1pt, roundcorner=5pt, linecolor=gray!30, backgroundcolor=gray!5]
\textbf{Question 14} 

Which of the theorem defines the existence of Parikhs theorem?

a) Parikh's theorem

b) Jacobi theorem

c) AF+BG theorem

d) None of the mentioned
\vspace{5pt} \par
\begin{tcolorbox}[colback=green!5, colframe=green!40!black, title=Solution]
Answer: a

Explanation: Rohit Parikh in 1961 proved in his MIT research paper that some context free language can only have ambiguous grammars.
\end{tcolorbox}
\end{mdframed}
\newpage

\newtopic{Sentential Forms}

\begin{mdframed}[linewidth=1pt, roundcorner=5pt, linecolor=gray!30, backgroundcolor=gray!5]
\textbf{Question 1} 

State true or false:

Statement: Every right-linear grammar generates a regular language.

a) True

b) False
\vspace{5pt} \par
\begin{tcolorbox}[colback=green!5, colframe=green!40!black, title=Solution]
Answer: a

Explanation: A CFG is said to right linear if each production body has at most one variable, and that variable is at the right end. That is, all productions of a right linear grammar are of the form A-$>$wB or A-$>$w, where A and B are variables while w is some terminal.
\end{tcolorbox}
\end{mdframed}
\newpage

\begin{mdframed}[linewidth=1pt, roundcorner=5pt, linecolor=gray!30, backgroundcolor=gray!5]
\textbf{Question 2} 

What the does the given CFG defines?

S-$>$aSbS|bSaS|e and w denotes terminal

a) wwr

b) wSw

c) Equal number of a's and b's

d) None of the mentioned
\vspace{5pt} \par
\begin{tcolorbox}[colback=green!5, colframe=green!40!black, title=Solution]
Answer: c

Explanation: Using the derivation approach, we can conclude that the given grammar produces a language with a set of string which have equal number of a's and b's.
\end{tcolorbox}
\end{mdframed}
\newpage

\begin{mdframed}[linewidth=1pt, roundcorner=5pt, linecolor=gray!30, backgroundcolor=gray!5]
\textbf{Question 3} 

If L1 and L2 are context free languages, which of the following is context free?

a) L1*

b) L2UL1

c) L1.L2

d) All of the mentioned
\vspace{5pt} \par
\begin{tcolorbox}[colback=green!5, colframe=green!40!black, title=Solution]
Answer: d

Explanation: The following is a theorem which states the closure property of context free languages which includes Kleene operation, Union operation and Dot operation.
\end{tcolorbox}
\end{mdframed}
\newpage

\begin{mdframed}[linewidth=1pt, roundcorner=5pt, linecolor=gray!30, backgroundcolor=gray!5]
\textbf{Question 4} 

For the given Regular expression, the minimum number of variables including starting variable required to derive its grammar is:

(011+1)*(01)*

a) 4

b) 3

c) 5

d) 6
\vspace{5pt} \par
\begin{tcolorbox}[colback=green!5, colframe=green!40!black, title=Solution]
Answer: c

Explanation: The grammar can be written as:

S-$>$BC

B-$>$AB|$\epsilon$

A-$>$011|1

C-$>$DC|$\epsilon$

D-$>$01
\end{tcolorbox}
\end{mdframed}
\newpage

\begin{mdframed}[linewidth=1pt, roundcorner=5pt, linecolor=gray!30, backgroundcolor=gray!5]
\textbf{Question 5} 

For the given Regular expression, the minimum number of terminals required to derive its grammar is:

(011+1)*(01)*

a) 4

b) 3

c) 5

d) 6
\vspace{5pt} \par
\begin{tcolorbox}[colback=green!5, colframe=green!40!black, title=Solution]
Answer: b

Explanation: The grammar can be written as:

S-$>$BC

B-$>$AB|$\epsilon$

A-$>$011|1

C-$>$DC|$\epsilon$

D-$>$01
\end{tcolorbox}
\end{mdframed}
\newpage

\begin{mdframed}[linewidth=1pt, roundcorner=5pt, linecolor=gray!30, backgroundcolor=gray!5]
\textbf{Question 6} 

A grammar G=(V, T, P, S) is \_\_\_\_\_\_\_\_\_\_ if every production taken one of the two forms:

B-$>$aC

B-$>$a

a) Ambiguous

b) Regular

c) Non Regular

d) None of the mentioned
\vspace{5pt} \par
\begin{tcolorbox}[colback=green!5, colframe=green!40!black, title=Solution]
Answer: b

Explanation: The following format of grammar is of Regular grammar and is a part of Context free grammar i.e. like a specific form whose finite automata can be generated.
\end{tcolorbox}
\end{mdframed}
\newpage

\begin{mdframed}[linewidth=1pt, roundcorner=5pt, linecolor=gray!30, backgroundcolor=gray!5]
\textbf{Question 7} 

Which among the following is a CFG for the given Language:

L=\{x$\in$\{0,1\}*|number of zeroes in x=number of one's in x\}

a) S-$>$e|0S1|1S0|SS

b) S-$>$0B|1A|e  A-$>$0S  B-$>$1S

c) All of the mentioned

d) None of the mentioned
\vspace{5pt} \par
\begin{tcolorbox}[colback=green!5, colframe=green!40!black, title=Solution]
Answer: c

Explanation: We can build context free grammar through different approaches, recursively defining the variables and terminals inorder to fulfil the conditions.
\end{tcolorbox}
\end{mdframed}
\newpage

\begin{mdframed}[linewidth=1pt, roundcorner=5pt, linecolor=gray!30, backgroundcolor=gray!5]
\textbf{Question 8} 

Which of the following languages are most suitable for implement context free languages ?

a) C

b) Perl

c) Assembly Language

d) None of the mentioned
\vspace{5pt} \par
\begin{tcolorbox}[colback=green!5, colframe=green!40!black, title=Solution]
Answer: a

Explanation: The advantage of using high level programming language like C and Pascal is that they allow us to write statements that look more like English.
\end{tcolorbox}
\end{mdframed}
\newpage

\begin{mdframed}[linewidth=1pt, roundcorner=5pt, linecolor=gray!30, backgroundcolor=gray!5]
\textbf{Question 9} 

Which among the following is the correct grammar for the given language?

L=\{x$\in$\{0,1\}*|number of zeroes in x¹number of one's in x\}

a) S-$>$ 0|SS|1SS|SS1|S1S

b) S-$>$ 1|0S|0SS|SS0|S0S

c) S-$>$ 0|0S|1SS|SS1|S1S

d) None of the mentioned
\vspace{5pt} \par
\begin{tcolorbox}[colback=green!5, colframe=green!40!black, title=Solution]
Answer: c

Explanation: L=\{0, 1, 00, 11, 001, 010,…\}

The grammar can be framed as:  S-$>$ 0|0S|1SS|SS1|S1S
\end{tcolorbox}
\end{mdframed}
\newpage

\begin{mdframed}[linewidth=1pt, roundcorner=5pt, linecolor=gray!30, backgroundcolor=gray!5]
\textbf{Question 10} 

L=\{0
i
1
j
2
k
 | j$>$i+k\}

Which of the following satisfies the language?

a) 0111100

b) 011100

c) 0001100

d) 0101010
\vspace{5pt} \par
\begin{tcolorbox}[colback=green!5, colframe=green!40!black, title=Solution]
Answer: a

Explanation: It is just required to put the value in the variables in the question and check if it satisfies or not.
\end{tcolorbox}
\end{mdframed}
\newpage

\newtopic{Construction and Yield of a Parse Tree}

\begin{mdframed}[linewidth=1pt, roundcorner=5pt, linecolor=gray!30, backgroundcolor=gray!5]
\textbf{Question 1} 

The most suitable data structure used to represent the derivations in compiler:

a) Queue

b) Linked List

c) Tree

d) Hash Tables
\vspace{5pt} \par
\begin{tcolorbox}[colback=green!5, colframe=green!40!black, title=Solution]
Answer: c

Explanation: The tree, known as "Parse tree" when used in a compiler, is the data structure of choice to represent the source program.
\end{tcolorbox}
\end{mdframed}
\newpage

\begin{mdframed}[linewidth=1pt, roundcorner=5pt, linecolor=gray!30, backgroundcolor=gray!5]
\textbf{Question 2} 

Which of the following statement is false in context of tree terminology?

a) Root with no children is called a leaf

b) A node can have three children

c) Root has no parent

d) Trees are collection of nodes, with a parent child relationship
\vspace{5pt} \par
\begin{tcolorbox}[colback=green!5, colframe=green!40!black, title=Solution]
Answer: a

Explanation: A node has atmost one parent, drawn above the node, and zero or more children drawn below. Lines connect parents to children. There is one node, one root, that has no parent; this node appears to be at the top of the tree. Nodes with no children are called leaves. Nodes that are not leaves are called interior nodes.
\end{tcolorbox}
\end{mdframed}
\newpage

\begin{mdframed}[linewidth=1pt, roundcorner=5pt, linecolor=gray!30, backgroundcolor=gray!5]
\textbf{Question 3} 

In which order are the children of any node ordered?

a) From the left

b) From the right

c) Arbitrarily

d) None of the mentioned
\vspace{5pt} \par
\begin{tcolorbox}[colback=green!5, colframe=green!40!black, title=Solution]
Answer: a

Explanation: The children of a node are ordered from the left and drawn so. If N is to the left of node M, then all the descendents of N are considered  to be to the left of all the descendents of M.
\end{tcolorbox}
\end{mdframed}
\newpage

\begin{mdframed}[linewidth=1pt, roundcorner=5pt, linecolor=gray!30, backgroundcolor=gray!5]
\textbf{Question 4} 

Which among the following is the root of the parse tree?

a) Production P

b) Terminal T

c) Variable V

d) Starting Variable S
\vspace{5pt} \par
\begin{tcolorbox}[colback=green!5, colframe=green!40!black, title=Solution]
Answer: d

Explanation: The root is labelled by the start symbol. All the leaves are either labelled by a a terminal or with e.
\end{tcolorbox}
\end{mdframed}
\newpage

\begin{mdframed}[linewidth=1pt, roundcorner=5pt, linecolor=gray!30, backgroundcolor=gray!5]
\textbf{Question 5} 

For the expression E*(E) where * and brackets are the operation, number of nodes in the respective parse tree are:

a) 6

b) 7

c) 5

d) 2
\vspace{5pt} \par
\begin{tcolorbox}[colback=green!5, colframe=green!40!black, title=Solution]
Answer: b

Explanation:
 
\begin{center}\includegraphics[width=0.5\textwidth]{automata-theory-questions-answers-construction-yield-parse-tree-q5.png}\end{center}
\end{tcolorbox}
\end{mdframed}
\newpage

\begin{mdframed}[linewidth=1pt, roundcorner=5pt, linecolor=gray!30, backgroundcolor=gray!5]
\textbf{Question 6} 

The number of leaves in a parse tree with expression E*(E) where * and () are operators

a) 5

b) 2

c) 4

d) 3
\vspace{5pt} \par
\begin{tcolorbox}[colback=green!5, colframe=green!40!black, title=Solution]
Answer: a

Explanation: 
 
\begin{center}\includegraphics[width=0.5\textwidth]{automata-theory-questions-answers-construction-yield-parse-tree-q5.png}\end{center}
\end{tcolorbox}
\end{mdframed}
\newpage

\newtopic{Inferences to Trees and Trees to Derivations}

\begin{mdframed}[linewidth=1pt, roundcorner=5pt, linecolor=gray!30, backgroundcolor=gray!5]
\textbf{Question 1} 

A symbol X is \_\_\_\_\_\_\_\_ if there exists : S-$>$* aXb

a) reachable

b) generating

c) context free

d) none of the mentioned
\vspace{5pt} \par
\begin{tcolorbox}[colback=green!5, colframe=green!40!black, title=Solution]
Answer: a

Explanation: A symbol X is generating if there exists : X-$>$*w for some w that belongs to T*.

Also, a symbol can never be context free.
\end{tcolorbox}
\end{mdframed}
\newpage

\begin{mdframed}[linewidth=1pt, roundcorner=5pt, linecolor=gray!30, backgroundcolor=gray!5]
\textbf{Question 2} 

A symbol X is called to be useful if and only if its is:

a) generating

b) reachable

c) both generating and reachable

d) none of the mentioned
\vspace{5pt} \par
\begin{tcolorbox}[colback=green!5, colframe=green!40!black, title=Solution]
Answer: c

Explanation: For a symbol X to be useful, it has to be both reachable and generating i.e.

S-$>$* aXb -$>$ * w where w belongs to T*.
\end{tcolorbox}
\end{mdframed}
\newpage

\begin{mdframed}[linewidth=1pt, roundcorner=5pt, linecolor=gray!30, backgroundcolor=gray!5]
\textbf{Question 3} 

Which of the following is false for a grammar G in Chomsky Normal Form:

a) G has no useless symbols

b) G has no unit productions

c) G has no epsilon productions

d) None of the mentioned
\vspace{5pt} \par
\begin{tcolorbox}[colback=green!5, colframe=green!40!black, title=Solution]
Answer: d

Explanation: G, a CFG is said to be in Chomsky normal form if all its productions are in one of the following form:

A-$>$BC or A-$>$a
\end{tcolorbox}
\end{mdframed}
\newpage

\begin{mdframed}[linewidth=1pt, roundcorner=5pt, linecolor=gray!30, backgroundcolor=gray!5]
\textbf{Question 4} 

Given Checklist:

i) G has no useless symbols

ii) G has no unit productions

iii) G has no epsilon productions

iv) Normal form for production is violated

Is it possible for the grammar G to be in CNF with the following checklist?

a) Yes

b) No
\vspace{5pt} \par
\begin{tcolorbox}[colback=green!5, colframe=green!40!black, title=Solution]
Answer: b

Explanation: The grammar is not in CNF if it violates the normal form of the productions which is strictly restricted.
\end{tcolorbox}
\end{mdframed}
\newpage

\begin{mdframed}[linewidth=1pt, roundcorner=5pt, linecolor=gray!30, backgroundcolor=gray!5]
\textbf{Question 5} 

State true or false:

Statement: A CNF parse tree's string yield (w) can no longer be 2h-1.

a) true

b) false
\vspace{5pt} \par
\begin{tcolorbox}[colback=green!5, colframe=green!40!black, title=Solution]
Answer: a

Explanation: It is the parse tree theorem which states:

Given: Suppose we have a parse tree for a string w, according to a CNF grammar, G=(V, T, P, S). Let h be the height of the parse tree. Now, Implication: |w|$<$=2h-1.
\end{tcolorbox}
\end{mdframed}
\newpage

\begin{mdframed}[linewidth=1pt, roundcorner=5pt, linecolor=gray!30, backgroundcolor=gray!5]
\textbf{Question 6} 

If |w|$>$=2
h
, then its parse tree's height is at least \_\_\_\_\_

a) h

b) h+1

c) h-1

d) 2
h
\vspace{5pt} \par
\begin{tcolorbox}[colback=green!5, colframe=green!40!black, title=Solution]
Answer: b

Explanation: It is the basic implication of Parse tree theorem (assuming CNF). If the height of the parse tree is h, then |w| $<$=2
h-1
.
\end{tcolorbox}
\end{mdframed}
\newpage

\begin{mdframed}[linewidth=1pt, roundcorner=5pt, linecolor=gray!30, backgroundcolor=gray!5]
\textbf{Question 7} 

If w belongs to L(G), for some CFG, then w has a parse tree, which tell us the \_\_\_\_\_\_\_\_ structure of w.

a) semantic

b) syntactic

c) lexical

d) all of the mentioned
\vspace{5pt} \par
\begin{tcolorbox}[colback=green!5, colframe=green!40!black, title=Solution]
Answer: b

Explanation: A parse tree or concrete syntactic tree is an ordered, rooted tree that represents the syntactic structure of a string according to some context free grammar.
\end{tcolorbox}
\end{mdframed}
\newpage

\begin{mdframed}[linewidth=1pt, roundcorner=5pt, linecolor=gray!30, backgroundcolor=gray!5]
\textbf{Question 8} 

Which of the following are distinct to parse trees?

a) abstract parse trees

b) sentence diagrams

c) both abstract parse trees and sentence diagrams

d) none of the mentioned
\vspace{5pt} \par
\begin{tcolorbox}[colback=green!5, colframe=green!40!black, title=Solution]
Answer: c

Explanation: Both of the mentioned are different from parse trees. Sentence diagrams are pictorial representations of grammatical structure of a sentence.
\end{tcolorbox}
\end{mdframed}
\newpage

\begin{mdframed}[linewidth=1pt, roundcorner=5pt, linecolor=gray!30, backgroundcolor=gray!5]
\textbf{Question 9} 

Choose the correct option:

Statement: Unambiguity is the ideal structure of a language.

a) true

b) partially true

c) false

d) cant be said
\vspace{5pt} \par
\begin{tcolorbox}[colback=green!5, colframe=green!40!black, title=Solution]
Answer: a

Explanation: Ideally, there should be only one parse tree for each string, i.e. the language should be unambiguous.
\end{tcolorbox}
\end{mdframed}
\newpage

\begin{mdframed}[linewidth=1pt, roundcorner=5pt, linecolor=gray!30, backgroundcolor=gray!5]
\textbf{Question 10} 

Is the given statement correct?

Statement: The mere existence of several derivations is not an issue, its is the existence of several parse trees that ruins a grammar.

a) Yes

b) No
\vspace{5pt} \par
\begin{tcolorbox}[colback=green!5, colframe=green!40!black, title=Solution]
Answer: a

Explanation: It is also true that multiple leftmost or rightmost derivations do cause ambiguity. Unfortunately, it is not possible to remove the ambiguity always.
\end{tcolorbox}
\end{mdframed}
\newpage

\newtopic{Ambiguous Grammar}

\begin{mdframed}[linewidth=1pt, roundcorner=5pt, linecolor=gray!30, backgroundcolor=gray!5]
\textbf{Question 1} 

A CFG is ambiguous if

a) It has more than one rightmost derivations

b) It has more than one leftmost derivations

c) No parse tree can be generated for the CFG

d) None of the mentioned
\vspace{5pt} \par
\begin{tcolorbox}[colback=green!5, colframe=green!40!black, title=Solution]
Answer: b

Explanation: A context free grammar is ambiguous if it has more than one parse tree generated or more than one leftmost derivations. An unambiguous grammar is a context free grammar for which every valid string has a unique leftmost derivation.
\end{tcolorbox}
\end{mdframed}
\newpage

\begin{mdframed}[linewidth=1pt, roundcorner=5pt, linecolor=gray!30, backgroundcolor=gray!5]
\textbf{Question 2} 

Which of the following are always unambiguous?

a) Deterministic Context free grammars

b) Non-Deterministic Regular grammars

c) Context sensitive grammar

d) None of the mentioned
\vspace{5pt} \par
\begin{tcolorbox}[colback=green!5, colframe=green!40!black, title=Solution]
Answer: a

Explanation: Deterministic CFGs are always unambiguous , and are an important subclass of unambiguous CFGs; there are non-deterministic unambiguous CFGs, however.
\end{tcolorbox}
\end{mdframed}
\newpage

\begin{mdframed}[linewidth=1pt, roundcorner=5pt, linecolor=gray!30, backgroundcolor=gray!5]
\textbf{Question 3} 

A CFG is not closed under

a) Dot operation

b) Union Operation

c) Concatenation

d) Iteration
\vspace{5pt} \par
\begin{tcolorbox}[colback=green!5, colframe=green!40!black, title=Solution]
Answer: d

Explanation: The closure property of a context free grammar does not include iteration or kleene or star operation.
\end{tcolorbox}
\end{mdframed}
\newpage

\begin{mdframed}[linewidth=1pt, roundcorner=5pt, linecolor=gray!30, backgroundcolor=gray!5]
\textbf{Question 4} 

Which of the following is an real-world programming language ambiguity?

a) dangling else problem

b) halting problem

c) maze problem

d) none of the mentioned
\vspace{5pt} \par
\begin{tcolorbox}[colback=green!5, colframe=green!40!black, title=Solution]
Answer: a

Explanation: Dangling else problem: In many languages,the else in an if-then-else statement is optional, which results into nested conditionals being ambiguous, at least in terms of the CFG.
\end{tcolorbox}
\end{mdframed}
\newpage

\begin{mdframed}[linewidth=1pt, roundcorner=5pt, linecolor=gray!30, backgroundcolor=gray!5]
\textbf{Question 5} 

Which of the following is a parser for an ambiguous grammar?

a) GLR parser

b) Chart parser

c) All of the mentioned

d) None of the mentioned
\vspace{5pt} \par
\begin{tcolorbox}[colback=green!5, colframe=green!40!black, title=Solution]
Answer: c

Explanation: GLR parser: a type of parser for non deterministic and ambiguous grammar

Chart parser: aa type of parser for ambiguous grammar.
\end{tcolorbox}
\end{mdframed}
\newpage

\begin{mdframed}[linewidth=1pt, roundcorner=5pt, linecolor=gray!30, backgroundcolor=gray!5]
\textbf{Question 6} 

A language that admits only ambiguous grammar:

a) Inherent Ambiguous language

b) Inherent Unambiguous language

c) Context free language

d) Context Sensitive language
\vspace{5pt} \par
\begin{tcolorbox}[colback=green!5, colframe=green!40!black, title=Solution]
Answer: a

Explanation: A context free language for which no unambiguous grammar exists, is called Inherent ambiguous language.
\end{tcolorbox}
\end{mdframed}
\newpage

\begin{mdframed}[linewidth=1pt, roundcorner=5pt, linecolor=gray!30, backgroundcolor=gray!5]
\textbf{Question 7} 

Which of the following is an example of inherent ambiguous language?

a) \{a
n
|n$>$1\}

b) \{a
n
b
n
c
m
d
m
| n,m $>$ 0\}

c) \{0
n
1
n
|n$>$0\}

d) None of the mentioned
\vspace{5pt} \par
\begin{tcolorbox}[colback=green!5, colframe=green!40!black, title=Solution]
Answer: b

Explanation: This set is context-free, since the union of two context-free languages is always context free.
\end{tcolorbox}
\end{mdframed}
\newpage

\begin{mdframed}[linewidth=1pt, roundcorner=5pt, linecolor=gray!30, backgroundcolor=gray!5]
\textbf{Question 8} 

State true or false:

Statement: R-$>$R|T   T-$>$$\epsilon$ is an ambiguous grammar

a) true

b) false
\vspace{5pt} \par
\begin{tcolorbox}[colback=green!5, colframe=green!40!black, title=Solution]
Answer: a

Explanation: The production can be either itself or an empty string. Thus the empty string has more than one leftmost derivations, depending on how many times R-$>$R is being used.
\end{tcolorbox}
\end{mdframed}
\newpage

\begin{mdframed}[linewidth=1pt, roundcorner=5pt, linecolor=gray!30, backgroundcolor=gray!5]
\textbf{Question 9} 

In context to ambiguity, the number of times the following programming statement can be interpreted as:

Statement: if R then if T then P else V

a) 2

b) 3

c) 4

d) 1
\vspace{5pt} \par
\begin{tcolorbox}[colback=green!5, colframe=green!40!black, title=Solution]
Answer: a

Explanation:  Dangling else problem

if R then (if T then P else V) and if R then (if T then P) else V are the two ways in which the given if else statement can be parsed.
\end{tcolorbox}
\end{mdframed}
\newpage

\begin{mdframed}[linewidth=1pt, roundcorner=5pt, linecolor=gray!30, backgroundcolor=gray!5]
\textbf{Question 10} 

CFGs can be parsed in polynomial time using\_\_\_\_\_\_\_\_\_\_

a) LR parser

b) CYK algorithm

c) SLR parser

d) None of the mentioned
\vspace{5pt} \par
\begin{tcolorbox}[colback=green!5, colframe=green!40!black, title=Solution]
Answer: b

Explanation: CYK algorithm parses the CFG in polynomial time while LR parsers do the same in linear time. DCFGs are accepted by DPDAs and parsed using LR parsers or CYK algorithm.
\end{tcolorbox}
\end{mdframed}
\newpage

\newtopic{PDA-Acceptance by Final State}

\begin{mdframed}[linewidth=1pt, roundcorner=5pt, linecolor=gray!30, backgroundcolor=gray!5]
\textbf{Question 1} 

A push down automaton employs \_\_\_\_\_\_\_\_ data structure.

a) Queue

b) Linked List

c) Hash Table

d) Stack
\vspace{5pt} \par
\begin{tcolorbox}[colback=green!5, colframe=green!40!black, title=Solution]
Answer: d

Explanation: A push down automata uses a stack to carry out its operations. They are more capable than the finite automatons but less than the turing model.
\end{tcolorbox}
\end{mdframed}
\newpage

\begin{mdframed}[linewidth=1pt, roundcorner=5pt, linecolor=gray!30, backgroundcolor=gray!5]
\textbf{Question 2} 

State true or false:

Statement: The operations of PDA never work on elements, other than the top.

a) true

b) false
\vspace{5pt} \par
\begin{tcolorbox}[colback=green!5, colframe=green!40!black, title=Solution]
Answer: a

Explanation: The term pushdown refers to the fact that the elements are pushed down in the stack and as per the LIFO principle, the operation is always performed on the top element of the stack.
\end{tcolorbox}
\end{mdframed}
\newpage

\begin{mdframed}[linewidth=1pt, roundcorner=5pt, linecolor=gray!30, backgroundcolor=gray!5]
\textbf{Question 3} 

Which of the following allows stacked values to be sub-stacks rather than just finite symbols?

a) Push Down Automaton

b) Turing Machine

c) Nested Stack Automaton

d) None of the mentioned
\vspace{5pt} \par
\begin{tcolorbox}[colback=green!5, colframe=green!40!black, title=Solution]
Answer: c

Explanation: In computational theory, a nested stack automaton is a finite automaton which makes use of stack containing data which can be additional stacks.
\end{tcolorbox}
\end{mdframed}
\newpage

\begin{mdframed}[linewidth=1pt, roundcorner=5pt, linecolor=gray!30, backgroundcolor=gray!5]
\textbf{Question 4} 

A non deterministic two way, nested stack automaton has n-tuple definition. State the value of n.

a) 5

b) 8

c) 4

d) 10
\vspace{5pt} \par
\begin{tcolorbox}[colback=green!5, colframe=green!40!black, title=Solution]
Answer: d

Explanation: The 10-tuple can be stated as: NSA=  $<$Q,$\Sigma$,$\Gamma$,$\delta$,q0,Z0,F,[,],]$>$.
\end{tcolorbox}
\end{mdframed}
\newpage

\begin{mdframed}[linewidth=1pt, roundcorner=5pt, linecolor=gray!30, backgroundcolor=gray!5]
\textbf{Question 5} 

Push down automata accepts \_\_\_\_\_\_\_\_\_ languages.

a) Type 3

b) Type 2

c) Type 1

d) Type 0
\vspace{5pt} \par
\begin{tcolorbox}[colback=green!5, colframe=green!40!black, title=Solution]
Answer: b

Explanation: Push down automata is for Context free languages and they are termed as Type 2 languages according to Chomsky hierarchy.
\end{tcolorbox}
\end{mdframed}
\newpage

\begin{mdframed}[linewidth=1pt, roundcorner=5pt, linecolor=gray!30, backgroundcolor=gray!5]
\textbf{Question 6} 

The class of languages  not accepted by non deterministic, nonerasing stack automata is \_\_\_\_\_\_\_

a) NSPACE(n2)

b) NL

c) CSL

d) All of the mentioned
\vspace{5pt} \par
\begin{tcolorbox}[colback=green!5, colframe=green!40!black, title=Solution]
Answer: d

Explanation: NSPACE or non deterministic space is the computational resource describing the memory space for a non deterministic turing machine.
\end{tcolorbox}
\end{mdframed}
\newpage

\begin{mdframed}[linewidth=1pt, roundcorner=5pt, linecolor=gray!30, backgroundcolor=gray!5]
\textbf{Question 7} 

A push down automaton with only symbol allowed on the stack along with fixed symbol.

a) Embedded PDA

b) Nested Stack automata

c) DPDA

d) Counter Automaton
\vspace{5pt} \par
\begin{tcolorbox}[colback=green!5, colframe=green!40!black, title=Solution]
Answer: d

Explanation: This class of automata can recognize a set of context free languages like \{anbn|n belongs to N\}
\end{tcolorbox}
\end{mdframed}
\newpage

\begin{mdframed}[linewidth=1pt, roundcorner=5pt, linecolor=gray!30, backgroundcolor=gray!5]
\textbf{Question 8} 

Which of the operations are eligible in PDA?

a) Push

b) Delete

c) Insert

d) Add
\vspace{5pt} \par
\begin{tcolorbox}[colback=green!5, colframe=green!40!black, title=Solution]
Answer: a

Explanation: Push and pop are the operations we perform to operate a stack. A stack follows the LIFO principle, which states its rule as: Last In First Out.
\end{tcolorbox}
\end{mdframed}
\newpage

\begin{mdframed}[linewidth=1pt, roundcorner=5pt, linecolor=gray!30, backgroundcolor=gray!5]
\textbf{Question 9} 

A string is accepted by a PDA when

a) Stack is not empty

b) Acceptance state

c) All of the mentioned

d) None of the mentioned
\vspace{5pt} \par
\begin{tcolorbox}[colback=green!5, colframe=green!40!black, title=Solution]
Answer: b

Explanation: When we reach the acceptance state and find the stack to be empty, we say, the string has been accepted by the push down automata.
\end{tcolorbox}
\end{mdframed}
\newpage

\begin{mdframed}[linewidth=1pt, roundcorner=5pt, linecolor=gray!30, backgroundcolor=gray!5]
\textbf{Question 10} 

The following move of a PDA is on the basis of:

a) Present state

b) Input Symbol

c) Present state and Input Symbol

d) None of the mentioned
\vspace{5pt} \par
\begin{tcolorbox}[colback=green!5, colframe=green!40!black, title=Solution]
Answer: c

Explanation: The next operation is performed by PDA considering three factors: present state,symbol on the top of the stack and the input symbol.
\end{tcolorbox}
\end{mdframed}
\newpage

\newtopic{PDA-Acceptance by Empty Stack}

\begin{mdframed}[linewidth=1pt, roundcorner=5pt, linecolor=gray!30, backgroundcolor=gray!5]
\textbf{Question 1} 

If two sets, R and T has no elements in common i.e. RÇT=Æ, then the sets are called

a) Complement

b) Union

c) Disjoint

d) Connected
\vspace{5pt} \par
\begin{tcolorbox}[colback=green!5, colframe=green!40!black, title=Solution]
Answer: c

Explanation: Two sets are called disjoint if they have no elements in common i.e.  RÇT=Æ.
\end{tcolorbox}
\end{mdframed}
\newpage

\begin{mdframed}[linewidth=1pt, roundcorner=5pt, linecolor=gray!30, backgroundcolor=gray!5]
\textbf{Question 2} 

Which among the following is not a part of the Context free grammar tuple?

a) End symbol

b) Start symbol

c) Variable

d) Production
\vspace{5pt} \par
\begin{tcolorbox}[colback=green!5, colframe=green!40!black, title=Solution]
Answer: a

Explanation: The tuple definition of context free grammar is: (V, T, P, S) where V=set of variables, T=set of terminals, P=production, S= Starting Variable.
\end{tcolorbox}
\end{mdframed}
\newpage

\begin{mdframed}[linewidth=1pt, roundcorner=5pt, linecolor=gray!30, backgroundcolor=gray!5]
\textbf{Question 3} 

A context free grammar is a \_\_\_\_\_\_\_\_\_\_\_

a) A rule system for parsing formal languages

b) Regular grammar

c) Context sensitive grammar

d) None of the mentioned
\vspace{5pt} \par
\begin{tcolorbox}[colback=green!5, colframe=green!40!black, title=Solution]
Answer: a

Explanation: Context free grammar or CFG defines a set of rules for generating valid strings in a formal language, which can be programming languages, data formats, or other structured languages. Regular grammars are simpler than CFGs. Context-sensitive grammars are more powerful than CFGs and can handle complex structures based on surrounding context.
\end{tcolorbox}
\end{mdframed}
\newpage

\begin{mdframed}[linewidth=1pt, roundcorner=5pt, linecolor=gray!30, backgroundcolor=gray!5]
\textbf{Question 4} 

The closure property of context free grammar includes :

a) Kleene

b) Concatenation

c) Union

d) All of the mentioned
\vspace{5pt} \par
\begin{tcolorbox}[colback=green!5, colframe=green!40!black, title=Solution]
Answer: d

Explanation: Context free grammars are closed under kleene operation, union and concatenation too.
\end{tcolorbox}
\end{mdframed}
\newpage

\begin{mdframed}[linewidth=1pt, roundcorner=5pt, linecolor=gray!30, backgroundcolor=gray!5]
\textbf{Question 5} 

Which of the following automata takes stack as auxiliary storage?

a) Finite automata

b) Push down automata

c) Turing machine

d) All of the mentioned
\vspace{5pt} \par
\begin{tcolorbox}[colback=green!5, colframe=green!40!black, title=Solution]
Answer: b

Explanation: Pushdown Automaton uses stack as an auxiliary storage for its operations. Turing machines use Queue for the same.
\end{tcolorbox}
\end{mdframed}
\newpage

\begin{mdframed}[linewidth=1pt, roundcorner=5pt, linecolor=gray!30, backgroundcolor=gray!5]
\textbf{Question 6} 

Which of the following automata takes queue as an auxiliary storage?

a) Finite automata

b) Push down automata

c) Turing machine

d) All of the mentioned
\vspace{5pt} \par
\begin{tcolorbox}[colback=green!5, colframe=green!40!black, title=Solution]
Answer: c

Explanation: Pushdown Automaton uses stack as an auxiliary storage for its operations. Turing machines use Queue for the same.
\end{tcolorbox}
\end{mdframed}
\newpage

\begin{mdframed}[linewidth=1pt, roundcorner=5pt, linecolor=gray!30, backgroundcolor=gray!5]
\textbf{Question 7} 

A context free grammar can be recognized by

a) Push down automata

b) 2 way linearly bounded automata

c) Push down automata \& 2 way linearly bounded automata

d) None of the mentioned
\vspace{5pt} \par
\begin{tcolorbox}[colback=green!5, colframe=green!40!black, title=Solution]
Answer: c

Explanation: A linearly bounded automata is a restricted non deterministic turing machine which is capable of accepting ant context free grammar.
\end{tcolorbox}
\end{mdframed}
\newpage

\begin{mdframed}[linewidth=1pt, roundcorner=5pt, linecolor=gray!30, backgroundcolor=gray!5]
\textbf{Question 8} 

A null production can be referred to as:

a) String

b) Symbol

c) Word

d) All of the mentioned
\vspace{5pt} \par
\begin{tcolorbox}[colback=green!5, colframe=green!40!black, title=Solution]
Answer: a

Explanation: Null production is always taken as a string in computational theory.
\end{tcolorbox}
\end{mdframed}
\newpage

\begin{mdframed}[linewidth=1pt, roundcorner=5pt, linecolor=gray!30, backgroundcolor=gray!5]
\textbf{Question 9} 

The context free grammar which generates a Regular Language is termed as:

a) Context Regular Grammar

b) Regular Grammar

c) Context Sensitive Grammar

d) None of the mentioned
\vspace{5pt} \par
\begin{tcolorbox}[colback=green!5, colframe=green!40!black, title=Solution]
Answer: b

Explanation: Regular grammar is a subset of Context free grammar. The CFGs which produces a language for which a finite automaton can be created is called Regular grammar.
\end{tcolorbox}
\end{mdframed}
\newpage

\begin{mdframed}[linewidth=1pt, roundcorner=5pt, linecolor=gray!30, backgroundcolor=gray!5]
\textbf{Question 10} 

NPDA stands for

a) Non-Deterministic Push Down Automata

b) Null-Push Down Automata

c) Nested Push Down Automata

d) All of the mentioned
\vspace{5pt} \par
\begin{tcolorbox}[colback=green!5, colframe=green!40!black, title=Solution]
Answer: a

Explanation: NPDA stands for non-deterministic push down automata whereas DPDA stands for deterministic push down automata.
\end{tcolorbox}
\end{mdframed}
\newpage

\newtopic{From Grammars to Push Down Automata}

\begin{mdframed}[linewidth=1pt, roundcorner=5pt, linecolor=gray!30, backgroundcolor=gray!5]
\textbf{Question 1} 

The production of the form A-$>$B , where A and B are non terminals is called

a) Null production

b) Unit production

c) Greibach Normal Form

d) Chomsky Normal Form
\vspace{5pt} \par
\begin{tcolorbox}[colback=green!5, colframe=green!40!black, title=Solution]
Answer: b

Explanation: A-$>$$\epsilon$ is termed as Null production while A-$>$B is termed as Unit production.
\end{tcolorbox}
\end{mdframed}
\newpage

\begin{mdframed}[linewidth=1pt, roundcorner=5pt, linecolor=gray!30, backgroundcolor=gray!5]
\textbf{Question 2} 

Halting states are of two types. They are:

a) Accept and Reject

b) Reject and Allow

c) Start and Reject

d) None of the mentioned
\vspace{5pt} \par
\begin{tcolorbox}[colback=green!5, colframe=green!40!black, title=Solution]
Answer: a

Explanation: Halting states are the new tuple members introduced in turing machine and is of two types: Accept Halting State and Reject Halting State.
\end{tcolorbox}
\end{mdframed}
\newpage

\begin{mdframed}[linewidth=1pt, roundcorner=5pt, linecolor=gray!30, backgroundcolor=gray!5]
\textbf{Question 3} 

A push down automata can be represented as:

PDA= $\epsilon$-NFA +[stack]
State true or false:

a) true

b) false
\vspace{5pt} \par
\begin{tcolorbox}[colback=green!5, colframe=green!40!black, title=Solution]
Answer: a

Explanation:


 
\begin{center}\includegraphics[width=0.5\textwidth]{automata-theory-basic-questions-answers-q3.png}\end{center}
\end{tcolorbox}
\end{mdframed}
\newpage

\begin{mdframed}[linewidth=1pt, roundcorner=5pt, linecolor=gray!30, backgroundcolor=gray!5]
\textbf{Question 4} 

A pushdown automata can be defined as: (Q, $\Sigma$, G, q0, z0, A, d)

What does the symbol z0 represents?

a) an element of G

b) initial stack symbol

c) top stack alphabet

d) all of the mentioned
\vspace{5pt} \par
\begin{tcolorbox}[colback=green!5, colframe=green!40!black, title=Solution]
Answer: d

Explanation: z0 is the initial stack symbol, is an element of G. Other symbols like d represents the transition function of the machine.
\end{tcolorbox}
\end{mdframed}
\newpage

\begin{mdframed}[linewidth=1pt, roundcorner=5pt, linecolor=gray!30, backgroundcolor=gray!5]
\textbf{Question 5} 

Which of the following correctly recognize the symbol ‘|-‘ in context to PDA?

a) Moves

b) transition function

c) or/not symbol

d) none of the mentioned
\vspace{5pt} \par
\begin{tcolorbox}[colback=green!5, colframe=green!40!black, title=Solution]
Answer: a

Explanation: Using this notation, we can define moves and further acceptance of a string by the machine.
\end{tcolorbox}
\end{mdframed}
\newpage

\begin{mdframed}[linewidth=1pt, roundcorner=5pt, linecolor=gray!30, backgroundcolor=gray!5]
\textbf{Question 6} 

Which among the following is true for the given statement?

Statement :If there are strings R and T in a language L so that R is prefix of T and R is not equivalent to T.

a) No DPDA can accept L by empty stack

b) DPDA can accept L by an empty stack

c) L is regular

d) None of the mentioned
\vspace{5pt} \par
\begin{tcolorbox}[colback=green!5, colframe=green!40!black, title=Solution]
Answer: a

Explanation: If M is a DPDA accepting L by an empty stsck, R and T are distinct strings in L, and R is a prefix of T, then the sequence of moves M must make in order to accept R leaves the stack empty, since R$\in$L. But then T cannot be accepted, since M cant move with an empty stack.
\end{tcolorbox}
\end{mdframed}
\newpage

\begin{mdframed}[linewidth=1pt, roundcorner=5pt, linecolor=gray!30, backgroundcolor=gray!5]
\textbf{Question 7} 

Which of the following can be accepted by a DPDA?

a) The set of even length palindrome over \{a,b\}

b) The set of odd length palindrome over \{a,b\}

c) \{xx
c
| where c stands for the complement,\{0,1\}\}

d) None of the mentioned
\vspace{5pt} \par
\begin{tcolorbox}[colback=green!5, colframe=green!40!black, title=Solution]
Answer: d

Explanation: Theorem: The language pal of palindromes over the alphabet \{0,1\} cannot be accepted by any finite automaton , and it is therefore not regular.
\end{tcolorbox}
\end{mdframed}
\newpage

\begin{mdframed}[linewidth=1pt, roundcorner=5pt, linecolor=gray!30, backgroundcolor=gray!5]
\textbf{Question 8} 

For a counter automaton, with the symbols A and Z0, the string on the stack is always in the form of \_\_\_\_\_\_\_\_\_\_

a) A

b) A
n
Z0, n$>$=0

c) Z0A
n
, n$>$=0

d) None of the mentioned
\vspace{5pt} \par
\begin{tcolorbox}[colback=green!5, colframe=green!40!black, title=Solution]
Answer: b

Explanation:The possible change in the stack contents is a change in the number of A's on the stack.
\end{tcolorbox}
\end{mdframed}
\newpage

\begin{mdframed}[linewidth=1pt, roundcorner=5pt, linecolor=gray!30, backgroundcolor=gray!5]
\textbf{Question 9} 

State true or false:

Statement: Counter Automaton can exist for the language L=\{0
i
1
i
|i$>$=0\}

a) true

b) false
\vspace{5pt} \par
\begin{tcolorbox}[colback=green!5, colframe=green!40!black, title=Solution]
Answer: a

Explanation: The PDA works as follows. Instead of saving excess 0's or 1's on the stack, we save *'s and use two different states to indicate which symbol there is currently a surplus of. The state q0 is the initial state and the only accepting state.
\end{tcolorbox}
\end{mdframed}
\newpage

\begin{mdframed}[linewidth=1pt, roundcorner=5pt, linecolor=gray!30, backgroundcolor=gray!5]
\textbf{Question 10} 

Let $\Sigma$=\{0,1\}* and the grammar G be:

S-$>$$\epsilon$

S-$>$SS

S-$>$0S1|1S0

State which of the following is true for the given

a) Language of all and only Balanced strings

b) It contains equal number of 0's and 1's

c) Ambiguous Grammar

d) All of the mentioned
\vspace{5pt} \par
\begin{tcolorbox}[colback=green!5, colframe=green!40!black, title=Solution]
Answer: d

Explanation: A string is said to be balanced if it consist of equal number of 0's and 1's.
\end{tcolorbox}
\end{mdframed}
\newpage

\newtopic{From PDA to Grammars}

\begin{mdframed}[linewidth=1pt, roundcorner=5pt, linecolor=gray!30, backgroundcolor=gray!5]
\textbf{Question 1} 

The instantaneous PDA is has the following elements

a) State

b) Unconsumed input

c) Stack content

d) All of the mentioned
\vspace{5pt} \par
\begin{tcolorbox}[colback=green!5, colframe=green!40!black, title=Solution]
Answer: d

Explanation: The instantaneous description of a PDA is represented by 3 tuple:

(q,w,s)

where q is the state, w is the unconsumed input and s is the stack content.
\end{tcolorbox}
\end{mdframed}
\newpage

\begin{mdframed}[linewidth=1pt, roundcorner=5pt, linecolor=gray!30, backgroundcolor=gray!5]
\textbf{Question 2} 

The moves in the PDA is technically termed as:

a) Turnstile

b) Shifter

c) Router

d) None of the mentioned
\vspace{5pt} \par
\begin{tcolorbox}[colback=green!5, colframe=green!40!black, title=Solution]
Answer: a

Explanation: A turnstile notation is used for connecting pairs od ID's taht represents one or many moves of a PDA.
\end{tcolorbox}
\end{mdframed}
\newpage

\begin{mdframed}[linewidth=1pt, roundcorner=5pt, linecolor=gray!30, backgroundcolor=gray!5]
\textbf{Question 3} 

Which of the following option resembles the given PDA?


 
\begin{center}\includegraphics[width=0.5\textwidth]{automata-theory-questions-answers-from-pda-grammars-q3.png}\end{center}
 

a) \{0
n
1
n
|n$>$=0\}

b) \{0
n
1
2n
|n$>$=0\}

c) \{0
2n
1
n
|n$>$=0\}

d) None of the mentioned
\vspace{5pt} \par
\begin{tcolorbox}[colback=green!5, colframe=green!40!black, title=Solution]
Answer: a

Explanation: None.
\end{tcolorbox}
\end{mdframed}
\newpage

\begin{mdframed}[linewidth=1pt, roundcorner=5pt, linecolor=gray!30, backgroundcolor=gray!5]
\textbf{Question 4} 

Which of the following correctly resembles the given state diagram?


 
\begin{center}\includegraphics[width=0.5\textwidth]{automata-theory-questions-answers-from-pda-grammars-q4.png}\end{center}
 

a) \{ww
r
|w=(a+b)*\}

b) $\epsilon$ is called the initial stack symbol

c) All of the mentioned

d) None of the mentioned
\vspace{5pt} \par
\begin{tcolorbox}[colback=green!5, colframe=green!40!black, title=Solution]
Answer: a

Explanation: Initially we put a special symbol ‘\#' into the empty stack. At state q1, the w is being read. In state q2, each 0 or 1 is popped when it matches the input. If any other input is given, the PDA will go to a dead state. When we reach that special symbol ‘\#', we go to the accepting state q3.
\end{tcolorbox}
\end{mdframed}
\newpage

\begin{mdframed}[linewidth=1pt, roundcorner=5pt, linecolor=gray!30, backgroundcolor=gray!5]
\textbf{Question 5} 

Which of the following assertion is false?

a) If L is a language accepted by PDA1 by final state, there exist a PDA2 that accepts L by empty stack i.e. L=L(PDA1)=L(PDA2)

b) If L is a CFL then there exists a push down automata P accepting CF; by empty stack i.e. L=M(P)

c) Let L is a language accepted by PDA1 then there exist a CFG X such that L(X)=M(P)

d) All of the mentioned
\vspace{5pt} \par
\begin{tcolorbox}[colback=green!5, colframe=green!40!black, title=Solution]
Answer: d

Explanation: All the assertions mentioned are theorems or corollary.
\end{tcolorbox}
\end{mdframed}
\newpage

\begin{mdframed}[linewidth=1pt, roundcorner=5pt, linecolor=gray!30, backgroundcolor=gray!5]
\textbf{Question 6} 

A push down automata  can represented using:

a) Transition graph

b) Transition table

c) ID

d) All of the mentioned
\vspace{5pt} \par
\begin{tcolorbox}[colback=green!5, colframe=green!40!black, title=Solution]
Answer: d

Explanation: Yes, a PDA can be represented using a transition diagram, transition table and an instantaneous description.
\end{tcolorbox}
\end{mdframed}
\newpage

\begin{mdframed}[linewidth=1pt, roundcorner=5pt, linecolor=gray!30, backgroundcolor=gray!5]
\textbf{Question 7} 

State true or false:

Statement: Every context free grammar can be transformed into an equvalent non deterministic push down automata.

a) true

b) false
\vspace{5pt} \par
\begin{tcolorbox}[colback=green!5, colframe=green!40!black, title=Solution]
Answer: a

Explanation: Push down automata is the automaton machine for all the context free grammar or Type 2 languages.
\end{tcolorbox}
\end{mdframed}
\newpage

\begin{mdframed}[linewidth=1pt, roundcorner=5pt, linecolor=gray!30, backgroundcolor=gray!5]
\textbf{Question 8} 

Which of the following statement is false?

a) For non deterministic PDA, equivalence is undecidable

b) For deterministic PDA, equivalence is decidable

c) For deterministic PDA, equivalence is undecidable

d) None of the mentioned
\vspace{5pt} \par
\begin{tcolorbox}[colback=green!5, colframe=green!40!black, title=Solution]
Answer: c

Explanation: Geraud proved the equivalence problem decidable for Deterministic PDA .
\end{tcolorbox}
\end{mdframed}
\newpage

\begin{mdframed}[linewidth=1pt, roundcorner=5pt, linecolor=gray!30, backgroundcolor=gray!5]
\textbf{Question 9} 

Which of the following are the actions that operates on stack top?

a) Pushing

b) Popping

c) Replacing

d) All of the mentioned
\vspace{5pt} \par
\begin{tcolorbox}[colback=green!5, colframe=green!40!black, title=Solution]
Answer: d

Explanation: Push, pop and replace are all the basic and only operations that takes place on stack top.
\end{tcolorbox}
\end{mdframed}
\newpage

\begin{mdframed}[linewidth=1pt, roundcorner=5pt, linecolor=gray!30, backgroundcolor=gray!5]
\textbf{Question 10} 

A push down automata is said to be \_\_\_\_\_\_\_\_\_ if it has atmost one transition around all configurations.

a) Finite

b) Non regular

c) Non-deterministic

d) Deterministic
\vspace{5pt} \par
\begin{tcolorbox}[colback=green!5, colframe=green!40!black, title=Solution]
Answer: d

Explanation: DPDA or Deterministic Push down automata has atmost one transition applicable to each configuration.
\end{tcolorbox}
\end{mdframed}
\newpage

\newtopic{Deterministic PDA}

\begin{mdframed}[linewidth=1pt, roundcorner=5pt, linecolor=gray!30, backgroundcolor=gray!5]
\textbf{Question 1} 

The transition a Push down automaton makes is additionally dependent upon the:

a) stack

b) input tape

c) terminals

d) none of the mentioned
\vspace{5pt} \par
\begin{tcolorbox}[colback=green!5, colframe=green!40!black, title=Solution]
Answer: a

Explanation: A PDA is a finite machine which has an additional stack storage. Its transitions are based not only on input and the correct state but also on the stack.
\end{tcolorbox}
\end{mdframed}
\newpage

\begin{mdframed}[linewidth=1pt, roundcorner=5pt, linecolor=gray!30, backgroundcolor=gray!5]
\textbf{Question 2} 

A PDA machine configuration (p, w, y) can be correctly represented as:

a) (current state, unprocessed input, stack content)

b) (unprocessed input, stack content, current state)

c) (current state, stack content, unprocessed input)

d) none of the mentioned
\vspace{5pt} \par
\begin{tcolorbox}[colback=green!5, colframe=green!40!black, title=Solution]
Answer: a

Explanation: A machine configuration is an element of K×$\Sigma$*×$\Gamma$*.

(p,w,$\gamma$) = (current state, unprocessed input, stack content).
\end{tcolorbox}
\end{mdframed}
\newpage

\begin{mdframed}[linewidth=1pt, roundcorner=5pt, linecolor=gray!30, backgroundcolor=gray!5]
\textbf{Question 3} 

|-* is the \_\_\_\_\_\_\_\_\_\_ closure of |-

a) symmetric and reflexive

b) transitive and reflexive

c) symmetric and transitive

d) none of the mentioned
\vspace{5pt} \par
\begin{tcolorbox}[colback=green!5, colframe=green!40!black, title=Solution]
Answer: b

Explanation: A string w is accepted by a PDA if and only if (s,w, e) |-* (f, e, e)
\end{tcolorbox}
\end{mdframed}
\newpage

\begin{mdframed}[linewidth=1pt, roundcorner=5pt, linecolor=gray!30, backgroundcolor=gray!5]
\textbf{Question 4} 

With reference of a DPDA, which among the following do we perform from the start state with an empty stack?

a) process the whole string

b) end in final state

c) end with an empty stack

d) all of the mentioned
\vspace{5pt} \par
\begin{tcolorbox}[colback=green!5, colframe=green!40!black, title=Solution]
Answer: d

Explanation: The empty stack in the end is our requirement relative to finite state automatons.
\end{tcolorbox}
\end{mdframed}
\newpage

\begin{mdframed}[linewidth=1pt, roundcorner=5pt, linecolor=gray!30, backgroundcolor=gray!5]
\textbf{Question 5} 

A DPDA is a PDA in which:

a) No state p has two outgoing transitions

b) More than one state can have two or more outgoing transitions

c) Atleast one state has more than one transitions

d) None of the mentioned
\vspace{5pt} \par
\begin{tcolorbox}[colback=green!5, colframe=green!40!black, title=Solution]
Answer: a

Explanation: A Deterministic Push Down Automata is a Push Down Automata in which no state p has two or more transitions.
\end{tcolorbox}
\end{mdframed}
\newpage

\begin{mdframed}[linewidth=1pt, roundcorner=5pt, linecolor=gray!30, backgroundcolor=gray!5]
\textbf{Question 6} 

State true or false:

Statement: For every CFL, G, there exists a PDA M such that L(G) = L(M) and vice versa.

a) true

b) false
\vspace{5pt} \par
\begin{tcolorbox}[colback=green!5, colframe=green!40!black, title=Solution]
Answer: a

Explanation: There exists two lemma's such that:

a) Given a grammar G, construct the PDA and show the equivalence

b) Given a PDA, construct a grammar and show the equivalence
\end{tcolorbox}
\end{mdframed}
\newpage

\begin{mdframed}[linewidth=1pt, roundcorner=5pt, linecolor=gray!30, backgroundcolor=gray!5]
\textbf{Question 7} 

If the PDA does not stop on an accepting state and the stack is not empty, the string is:

a) rejected

b) goes into loop forever

c) all of the mentioned

d) none of the mentioned
\vspace{5pt} \par
\begin{tcolorbox}[colback=green!5, colframe=green!40!black, title=Solution]
Answer: a

Explanation: To accept a string, PDA needs to halt at an accepting state and with a stack empty, else it is called rejected. Given a PDA M, we can construct a PDA M' that accepts the same language as M, by both acceptance criteria.
\end{tcolorbox}
\end{mdframed}
\newpage

\begin{mdframed}[linewidth=1pt, roundcorner=5pt, linecolor=gray!30, backgroundcolor=gray!5]
\textbf{Question 8} 

A language accepted by Deterministic Push down automata is closed under which of the following?

a) Complement

b) Union

c) All of the mentioned

d) None of the mentioned
\vspace{5pt} \par
\begin{tcolorbox}[colback=green!5, colframe=green!40!black, title=Solution]
Answer: a

Explanation: Deterministic Context free languages(one accepted by PDA by final state), are drastically different from the context free languages. For example they are closed under complementation and not union.
\end{tcolorbox}
\end{mdframed}
\newpage

\begin{mdframed}[linewidth=1pt, roundcorner=5pt, linecolor=gray!30, backgroundcolor=gray!5]
\textbf{Question 9} 

Which of the following is a simulator for non deterministic automata?

a) JFLAP

b) Gedit

c) FAUTO

d) None of the mentioned
\vspace{5pt} \par
\begin{tcolorbox}[colback=green!5, colframe=green!40!black, title=Solution]
Answer: a

Explanation: JFLAP is a software for experimenting with formal topics including NFA, NPDA, multi-tape turing machines and L-systems.
\end{tcolorbox}
\end{mdframed}
\newpage

\begin{mdframed}[linewidth=1pt, roundcorner=5pt, linecolor=gray!30, backgroundcolor=gray!5]
\textbf{Question 10} 

Finite-state acceptors for the nested words can be:

a) nested word automata

b) push down automata

c) ndfa

d) none of the mentioned
\vspace{5pt} \par
\begin{tcolorbox}[colback=green!5, colframe=green!40!black, title=Solution]
Answer: a

Explanation: The linear encodings of languages accepted by finite nested word automata gives the class of ‘visibly pushdown automata'.
\end{tcolorbox}
\end{mdframed}
\newpage

\newtopic{Regular Languages and D-PDA}

\begin{mdframed}[linewidth=1pt, roundcorner=5pt, linecolor=gray!30, backgroundcolor=gray!5]
\textbf{Question 1} 

Which of the following is analogous to the following?

:NFA and NPDA

a) Regular language and Context Free language

b) Regular language and Context Sensitive language

c) Context free language and Context Sensitive language

d) None of the mentioned
\vspace{5pt} \par
\begin{tcolorbox}[colback=green!5, colframe=green!40!black, title=Solution]
Answer: a

Explanation: All regular languages can be accepted by a non deterministic finite automata and all context free languages can be accepted by a non deterministic push down automata.
\end{tcolorbox}
\end{mdframed}
\newpage

\begin{mdframed}[linewidth=1pt, roundcorner=5pt, linecolor=gray!30, backgroundcolor=gray!5]
\textbf{Question 2} 

Let T=\{p, q, r, s, t\}. The number of strings in S* of length 4 such that no symbols can be repeated.

a) 120

b) 625

c) 360

d) 36
\vspace{5pt} \par
\begin{tcolorbox}[colback=green!5, colframe=green!40!black, title=Solution]
Answer: b

Explanation: Using the permutation rule, we can calculate that there will be total of 625 permutations on 5 elements taking 4 as the length.
\end{tcolorbox}
\end{mdframed}
\newpage

\begin{mdframed}[linewidth=1pt, roundcorner=5pt, linecolor=gray!30, backgroundcolor=gray!5]
\textbf{Question 3} 

Which of the following relates to Chomsky hierarchy?

a) Regular$<$CFL$<$CSL$<$Unrestricted

b) CFL$<$CSL$<$Unrestricted$<$Regular

c) CSL$<$Unrestricted$<$CF$<$Regular

d) None of the mentioned
\vspace{5pt} \par
\begin{tcolorbox}[colback=green!5, colframe=green!40!black, title=Solution]
Answer: a

Explanation: The chomsky hierarchy lays down the following order: Regular$<$CFL$<$CSL$<$Unrestricted
\end{tcolorbox}
\end{mdframed}
\newpage

\begin{mdframed}[linewidth=1pt, roundcorner=5pt, linecolor=gray!30, backgroundcolor=gray!5]
\textbf{Question 4} 

A language is accepted by a push down automata if it is:

a) regular

b) context free

c) regular and context free

d) none of the mentioned
\vspace{5pt} \par
\begin{tcolorbox}[colback=green!5, colframe=green!40!black, title=Solution]
Answer: c

Explanation: All the regular languages are the subset to context free languages and thus can be accepted using push down automata.
\end{tcolorbox}
\end{mdframed}
\newpage

\begin{mdframed}[linewidth=1pt, roundcorner=5pt, linecolor=gray!30, backgroundcolor=gray!5]
\textbf{Question 5} 

Which of the following is an incorrect regular expression identity?

a) R+f=R

b) eR=e

c) Rf=f

d) None of the mentioned
\vspace{5pt} \par
\begin{tcolorbox}[colback=green!5, colframe=green!40!black, title=Solution]
Answer: b

Explanation: e is the identity for concatenation. Thus, eR=R.
\end{tcolorbox}
\end{mdframed}
\newpage

\begin{mdframed}[linewidth=1pt, roundcorner=5pt, linecolor=gray!30, backgroundcolor=gray!5]
\textbf{Question 6} 

Which of the following strings do not belong the given regular expression?

(a)*(a+cba)

a) aa

b) aaa

c) acba

d) acbacba
\vspace{5pt} \par
\begin{tcolorbox}[colback=green!5, colframe=green!40!black, title=Solution]
Answer: d

Explanation: The string acbacba is unacceptable by the regular expression (a)*(a+cba).
\end{tcolorbox}
\end{mdframed}
\newpage

\begin{mdframed}[linewidth=1pt, roundcorner=5pt, linecolor=gray!30, backgroundcolor=gray!5]
\textbf{Question 7} 

Which of the following regular expression allows strings on \{a,b\}* with length n where n is a multiple of 4.

a) (a+b+ab+ba+aa+bb+aba+bab+abab+baba)*

b) (bbbb+aaaa)*

c) ((a+b)(a+b)(a+b)(a+b))*

d) None of the mentioned
\vspace{5pt} \par
\begin{tcolorbox}[colback=green!5, colframe=green!40!black, title=Solution]
Answer: c

Explanation: Other mentioned options do not many of the combinations while option c seems most reliable.
\end{tcolorbox}
\end{mdframed}
\newpage

\begin{mdframed}[linewidth=1pt, roundcorner=5pt, linecolor=gray!30, backgroundcolor=gray!5]
\textbf{Question 8} 

Which of the following strings is not generated by the given grammar:

S-$>$SaSbS|e

a) aabb

b) abab

c) abaabb

d) None of the mentioned
\vspace{5pt} \par
\begin{tcolorbox}[colback=green!5, colframe=green!40!black, title=Solution]
Answer: d

Explanation: All the given options are generated by the given grammar. Using the methods of left and right derivations, it is simpler to look for string which a grammar can generate.
\end{tcolorbox}
\end{mdframed}
\newpage

\begin{mdframed}[linewidth=1pt, roundcorner=5pt, linecolor=gray!30, backgroundcolor=gray!5]
\textbf{Question 9} 

abb*c denotes which of the following?

a) \{abnc|n=0\}

b) \{abnc|n=1\}

c) \{anbc|n=0\}

d) \{abcn|n$>$0\}
\vspace{5pt} \par
\begin{tcolorbox}[colback=green!5, colframe=green!40!black, title=Solution]
Answer: b

Explanation: Here, the first mentioned b is fixed while the other can be zero or can be repeated any number of time.
\end{tcolorbox}
\end{mdframed}
\newpage

\begin{mdframed}[linewidth=1pt, roundcorner=5pt, linecolor=gray!30, backgroundcolor=gray!5]
\textbf{Question 10} 

The following denotion belongs to which type of language:

G=(V, T, P, S)

a) Regular grammar

b) Context free grammar

c) Context Sensitive grammar

d) All of the mentioned
\vspace{5pt} \par
\begin{tcolorbox}[colback=green!5, colframe=green!40!black, title=Solution]
Answer: b

Explanation: Ant formal  grammar is represented using a 4-tuple definition where V= finite set of variables, T= set of terminal characters, P=set of productions and S= Starting Variable with certain conditions based on the type of formal grammar.
\end{tcolorbox}
\end{mdframed}
\newpage

\newtopic{DPDA and Context Free Languages}

\begin{mdframed}[linewidth=1pt, roundcorner=5pt, linecolor=gray!30, backgroundcolor=gray!5]
\textbf{Question 1} 

Context free grammar is called Type 2 grammar because of \_\_\_\_\_\_\_\_\_\_\_\_\_\_ hierarchy.

a) Greibach

b) Backus

c) Chomsky

d) None of the mentioned
\vspace{5pt} \par
\begin{tcolorbox}[colback=green!5, colframe=green!40!black, title=Solution]
Answer: c

Explanation: Chomsky hierarchy decide four type of language :Type 3- Regular Language, Type 2-Context free language, Type 1-Context Sensitive Language, Type 0- Unrestricted or Recursively Ennumerable language.
\end{tcolorbox}
\end{mdframed}
\newpage

\begin{mdframed}[linewidth=1pt, roundcorner=5pt, linecolor=gray!30, backgroundcolor=gray!5]
\textbf{Question 2} 

a→b

Restriction: Length of b must be atleast as much length of a.

Which of the following is correct for the given assertion?

a) Greibach Normal form

b) Context Sensitive Language

c) Chomsky Normal form

d) Recursively Ennumerable language
\vspace{5pt} \par
\begin{tcolorbox}[colback=green!5, colframe=green!40!black, title=Solution]
Answer: b

Explanation: A context-sensitive grammar (CSG) is a formal grammar in which the left-hand sides and right-hand sides of any production rules may be surrounded by a context of terminal and non terminal symbols. Context-sensitive grammars are more general than context-free grammars, in the sense that there are some languages that cannot be described by context-free grammars, but can be described by CSG.
\end{tcolorbox}
\end{mdframed}
\newpage

\begin{mdframed}[linewidth=1pt, roundcorner=5pt, linecolor=gray!30, backgroundcolor=gray!5]
\textbf{Question 3} 

From the definition of context free grammars,

G=(V, T, P, S)

What is the solution of VÇT?

a) Null

b) Not Null

c) Cannot be determined, depends on the language

d) None of the mentioned
\vspace{5pt} \par
\begin{tcolorbox}[colback=green!5, colframe=green!40!black, title=Solution]
Answer: a

Explanation: V is the set of non terminal symbols while T is the st of terminal symbols, their intersection would always be null.
\end{tcolorbox}
\end{mdframed}
\newpage

\begin{mdframed}[linewidth=1pt, roundcorner=5pt, linecolor=gray!30, backgroundcolor=gray!5]
\textbf{Question 4} 

If P is the production, for the given statement, state true or false.

P: V-$>$(V$\Sigma$T)* represents that the left hand side production rule has no right or left context.

a) true

b) false
\vspace{5pt} \par
\begin{tcolorbox}[colback=green!5, colframe=green!40!black, title=Solution]
Answer: a

Explanation: Here the production P is from the definition of Context free grammar and thus, has no right or left context.
\end{tcolorbox}
\end{mdframed}
\newpage

\begin{mdframed}[linewidth=1pt, roundcorner=5pt, linecolor=gray!30, backgroundcolor=gray!5]
\textbf{Question 5} 

There exists a Context free grammar such that:

X-$>$aX

Which among the following is correct with respect to the given assertion?

a) Left Recursive Grammar

b) Right Recursive Grammar

c) Non Recursive Grammar

d) None of the mentioned
\vspace{5pt} \par
\begin{tcolorbox}[colback=green!5, colframe=green!40!black, title=Solution]
Answer: b

Explanation: The grammar with right recursive production is known as Right recursive grammar. Right recursive production is of the form X-$>$aX where a is a terminal and X is a non terminal.
\end{tcolorbox}
\end{mdframed}
\newpage

\begin{mdframed}[linewidth=1pt, roundcorner=5pt, linecolor=gray!30, backgroundcolor=gray!5]
\textbf{Question 6} 

If the partial derivation tree contains the root as the starting variable, the form is known as:

a) Chomsky hierarchy

b) Sentential form

c) Root form

d) None of the mentioned
\vspace{5pt} \par
\begin{tcolorbox}[colback=green!5, colframe=green!40!black, title=Solution]
Answer: b

Explanation: Example: For any grammar, productions be:

S-$>$AB

A-$>$aaA| \textasciicircum{}

B-$>$Bb| \textasciicircum{}

The partial derivation tree can be drawn as:


 
\begin{center}\includegraphics[width=0.5\textwidth]{automata-theory-questions-answers-dpda-context-free-languages-q6.png}\end{center}
 

Since it has the root as S, this can be said to be in sentential form.
\end{tcolorbox}
\end{mdframed}
\newpage

\begin{mdframed}[linewidth=1pt, roundcorner=5pt, linecolor=gray!30, backgroundcolor=gray!5]
\textbf{Question 7} 

Find a regular expression for a grammar which generates a language which states :

L contains a set of strings starting wth an a and ending with a b, with something in the middle.

a) a(a*Ub*)b

b) a*(aUb)b*

c) a(a*b*)b

d) None of the mentioned
\vspace{5pt} \par
\begin{tcolorbox}[colback=green!5, colframe=green!40!black, title=Solution]
Answer: a

Explanation: The grammar for the same language can be stated as :

(1) S → aMb

(2) M → A

(3) M → B

(4) A → e

(5) A → aA

(6) B → e

(7) B → bB
\end{tcolorbox}
\end{mdframed}
\newpage

\begin{mdframed}[linewidth=1pt, roundcorner=5pt, linecolor=gray!30, backgroundcolor=gray!5]
\textbf{Question 8} 

Which of the following is the correct representation of grammar for the given regular expression?

a(aUb)*b

a) 


    (1) S → aMb

    (2) M → e

    (3) M → aM

    (4) M → bM




b) 


    (1) S → aMb

    (2) M → Mab

    (3) M → aM

    (4) M → bM




c) 


    (1) S → aMb

    (2) M → e

    (3) M → aMb

    (4) M → bMa




d) None of the mentioned
\vspace{5pt} \par
\begin{tcolorbox}[colback=green!5, colframe=green!40!black, title=Solution]
Answer: a

Explanation:

The basic idea of grammar formalisms is to capture the structure of string by

a) using special symbols to stand for substrings of a particular structure

b) using rules to specify how the substrings are combined to form new substrings.
\end{tcolorbox}
\end{mdframed}
\newpage

\begin{mdframed}[linewidth=1pt, roundcorner=5pt, linecolor=gray!30, backgroundcolor=gray!5]
\textbf{Question 9} 

A CFG consist of the following elements:

a) a set of terminal symbols

b) a set of non terminal symbols

c) a set of productions

d) all of the mentioned
\vspace{5pt} \par
\begin{tcolorbox}[colback=green!5, colframe=green!40!black, title=Solution]
Answer: d

Explanation: A CFG consists of:

a) a set of terminals, which are characters of alphabets that appear in the string generated by the grammar.

b) a set of non terminals, which are placeholders for patterns of terminal symbols that can be generated by the nonterminal symbols.

c) a set of productions, which are set of rules to transit from one state to other forming up the string

d) a start symbol, a special non terminal symbol that appears in the initial string generated in the grammar.
\end{tcolorbox}
\end{mdframed}
\newpage

\begin{mdframed}[linewidth=1pt, roundcorner=5pt, linecolor=gray!30, backgroundcolor=gray!5]
\textbf{Question 10} 

A CFG for a program describing strings of letters with the word "main" somewhere in the string:

a)  


    
 -$>$ 
 m a i n 


    
 -$>$ 
 
 | epsilon

    
 -$>$ A | B | ... | Z | a | b ... | z




b) 


    
 --$>$ 
 m a i n 


    
 --$>$ 
 
 

    
 --$>$ A | B | ... | Z | a | b ... | z




c)  


    
 --$>$ 
 m a i n 


     
 --$>$ 
 | epsilon

     
 --$>$ A | B | ... | Z | a | b ... | z




d) None of the mentioned
\vspace{5pt} \par
\begin{tcolorbox}[colback=green!5, colframe=green!40!black, title=Solution]
Answer: a

Explanation: None.
\end{tcolorbox}
\end{mdframed}
\newpage

\newtopic{DPDA and Ambiguous Grammars}

\begin{mdframed}[linewidth=1pt, roundcorner=5pt, linecolor=gray!30, backgroundcolor=gray!5]
\textbf{Question 1} 

CFGs are more powerful than:

a) DFA

b) NDFA

c) Mealy Machine

d) All of the mentioned
\vspace{5pt} \par
\begin{tcolorbox}[colback=green!5, colframe=green!40!black, title=Solution]
Answer: d

Explanation:

Context-free grammars are strictly more powerful than regular expressions:

1) Any language that can be generated using regular expressions can be generated by a context-free grammar.

2) There are languages that can be generated by a context-free grammar that cannot be generated by any regular expression.

As a corollary, CFGs are strictly more powerful than DFAs and NDFAs.
\end{tcolorbox}
\end{mdframed}
\newpage

\begin{mdframed}[linewidth=1pt, roundcorner=5pt, linecolor=gray!30, backgroundcolor=gray!5]
\textbf{Question 2} 

State true or false:

S-$>$ 0S1|01

Statement: No regular expression exists for the given grammar.

a) true

b) false
\vspace{5pt} \par
\begin{tcolorbox}[colback=green!5, colframe=green!40!black, title=Solution]
Answer: a

Explanation: The grammar generates a language L such that L=\{0
n
1
n
|n$>$=1\} which is not regular. Thus, no regular expression exists for the same.
\end{tcolorbox}
\end{mdframed}
\newpage

\begin{mdframed}[linewidth=1pt, roundcorner=5pt, linecolor=gray!30, backgroundcolor=gray!5]
\textbf{Question 3} 

For the given set of code, the grammar representing real numbers in Pascal has error in one of the six lines. Fetch the error.

 (1)
 -$>$ 
 
 
 

 (2)
 -$>$ 
 
 | epsilon

 (3)
 -$>$  
 
 | epsilon

 (4)
 -$>$ ‘E' 
 
 
 | epsilon

 (5)
 -$>$ + | – | epsilon

 (6)
 -$>$ 0 | 1 | 2 | 3 | 4 | 5 | 6 | 7 | 8 | 9

a) 3

b) 4

c) 2

d) No errors
\vspace{5pt} \par
\begin{tcolorbox}[colback=green!5, colframe=green!40!black, title=Solution]
Answer: a

Explanation:


 –$>$ 
 
 
 


 –$>$ 
 
 | epsilon


 –$>$ ‘.' 
 
 | epsilon


 –$>$ ‘E' 
 
 
 | epsilon


 –$>$ + | – | epsilon


 –$>$ 0 | 1 | 2 | 3 | 4 | 5 | 6 | 7 | 8 | 9
\end{tcolorbox}
\end{mdframed}
\newpage

\begin{mdframed}[linewidth=1pt, roundcorner=5pt, linecolor=gray!30, backgroundcolor=gray!5]
\textbf{Question 4} 

Which among the following is incorrect with reference to a derivation tree?

a) Every vertex has a label which is a terminal or a variable.

b) The root has a label which can be a terminal.

c) The label of the internal vertex is a variable.

d) None of the mentioned
\vspace{5pt} \par
\begin{tcolorbox}[colback=green!5, colframe=green!40!black, title=Solution]
Answer: b

Explanation: The root or interms of the grammar, starting variable can not be a terminal.
\end{tcolorbox}
\end{mdframed}
\newpage

\begin{mdframed}[linewidth=1pt, roundcorner=5pt, linecolor=gray!30, backgroundcolor=gray!5]
\textbf{Question 5} 

Let G=(V, T, P, S)

where a production can be written as:

S-$>$aAS|a

A-$>$SbA|ba|SS

Which of the following string is produced by the grammar?

a) aabbaab

b) aabbaa

c) baabab

d) None of the mentioned
\vspace{5pt} \par
\begin{tcolorbox}[colback=green!5, colframe=green!40!black, title=Solution]
Answer: b

Explanation: The step wise grammar translation can be written as:

 aAS-$>$aSbaA-$>$aabAS-$>$aabbaa
\end{tcolorbox}
\end{mdframed}
\newpage

\begin{mdframed}[linewidth=1pt, roundcorner=5pt, linecolor=gray!30, backgroundcolor=gray!5]
\textbf{Question 6} 

Statement 1: Ambiguity is the property of grammar but not the language.

 Statement 2: Same language can have more than one grammar.

Which of the following options are correct with respect to the given statements?

a) Statement 1 is true but statement 2 is false

b) Statement 1 is false but statement 2 is true

c) Both the statements are true

d) Both the statements are false
\vspace{5pt} \par
\begin{tcolorbox}[colback=green!5, colframe=green!40!black, title=Solution]
Answer: c

Explanation: One language can more than one grammar. Some can be ambiguous and some cannot.
\end{tcolorbox}
\end{mdframed}
\newpage

\begin{mdframed}[linewidth=1pt, roundcorner=5pt, linecolor=gray!30, backgroundcolor=gray!5]
\textbf{Question 7} 

Which of the following are non essential while simplifying a grammar?

a) Removal of useless symbols

b) Removal of unit productions

c) Removal of null production

d) None of the mentioned
\vspace{5pt} \par
\begin{tcolorbox}[colback=green!5, colframe=green!40!black, title=Solution]
Answer: d

Explanation: Here are some process used to simplify a CFG but to produce an equivalent grammar:

a) Removal of useless symbols(non terminal) b) Removal of Unit productions and c) Removal of Null productions.
\end{tcolorbox}
\end{mdframed}
\newpage

\begin{mdframed}[linewidth=1pt, roundcorner=5pt, linecolor=gray!30, backgroundcolor=gray!5]
\textbf{Question 8} 

Which of the following are context free language?

a) L=\{a
i
b
i
|i$>$=0\}

b) L=\{ww
r
| w is a string and r represents reverse\}

c) All of the mentioned

d) one of the mentioned
\vspace{5pt} \par
\begin{tcolorbox}[colback=green!5, colframe=green!40!black, title=Solution]
Answer: a

Explanation: None.
\end{tcolorbox}
\end{mdframed}
\newpage

\begin{mdframed}[linewidth=1pt, roundcorner=5pt, linecolor=gray!30, backgroundcolor=gray!5]
\textbf{Question 9} 

The language L =\{a
i
2b
i
|i$>$=0\} is:

a) recursive

b) deterministic CFL

c) regular

d) Two of the mentioned is correct
\vspace{5pt} \par
\begin{tcolorbox}[colback=green!5, colframe=green!40!black, title=Solution]
Answer: d

Explanation: The language is recursive and every recursive language is a CFL.
\end{tcolorbox}
\end{mdframed}
\newpage

\begin{mdframed}[linewidth=1pt, roundcorner=5pt, linecolor=gray!30, backgroundcolor=gray!5]
\textbf{Question 10} 

L-$>$rLt|tLr|t|r

The given grammar produces a language which is:

a) All palindrome

b) All even palindromes

c) All odd palindromes

d) Strings with same begin and end symbols
\vspace{5pt} \par
\begin{tcolorbox}[colback=green!5, colframe=green!40!black, title=Solution]
Answer: c

Explanation:  As there exists no production for the palindrome set, even palindromes like abba, aabbaa, baaaaaab, etc will not be generated.
\end{tcolorbox}
\end{mdframed}
\newpage

\newtopic{CFG-Eliminating Useless Symbols}

\begin{mdframed}[linewidth=1pt, roundcorner=5pt, linecolor=gray!30, backgroundcolor=gray!5]
\textbf{Question 1} 

Suppose A-$>$xBz and B-$>$y, then the simplified grammar would be:

a) A-$>$xyz

b) A-$>$xBz|xyz

c) A-$>$xBz|B|y

d) none of the mentioned
\vspace{5pt} \par
\begin{tcolorbox}[colback=green!5, colframe=green!40!black, title=Solution]
Answer: a

Explanation: For the first step, substitute B in first production as it only produces terminal and remove B production as it has already been utilized.

We get A-$>$xBz|xyz and now, as B has no production, we eliminate the terms which hold the variable B, thus the answer remain A-$>$xyz.
\end{tcolorbox}
\end{mdframed}
\newpage

\begin{mdframed}[linewidth=1pt, roundcorner=5pt, linecolor=gray!30, backgroundcolor=gray!5]
\textbf{Question 2} 

Given Grammar: S-$>$A, A-$>$aA, A-$>$e, B-$>$bA

Which among the following productions are Useless productions?

a) S-$>$A

b) A-$>$aA

c) A-$>$e

d) B-$>$bA
\vspace{5pt} \par
\begin{tcolorbox}[colback=green!5, colframe=green!40!black, title=Solution]
Answer: d

Explanation: Some derivations are not reachable from the starting variable. As B is not reachable from the starting variable, it is a useless symbol and thus, can be eliminated.
\end{tcolorbox}
\end{mdframed}
\newpage

\begin{mdframed}[linewidth=1pt, roundcorner=5pt, linecolor=gray!30, backgroundcolor=gray!5]
\textbf{Question 3} 

Given:

S-$>$…-$>$xAy-$>$…-$>$w

if \_\_\_\_\_\_\_\_\_\_\_\_, then A is useful, else useless symbol.

a) A is a non terminal

b) A is a terminal

c) w Î L

d) w Ë L
\vspace{5pt} \par
\begin{tcolorbox}[colback=green!5, colframe=green!40!black, title=Solution]
Answer: c

Explanation: Whatever operation we perform in intermediate stages, if the string produced belongs to the language, A is termed as useful and if not, not a useful variable.
\end{tcolorbox}
\end{mdframed}
\newpage

\begin{mdframed}[linewidth=1pt, roundcorner=5pt, linecolor=gray!30, backgroundcolor=gray!5]
\textbf{Question 4} 

Given:

S-$>$aSb

S-$>$e

S-$>$ A

A-$>$aA

B-$>$C

C-$>$D

The ratio of number of useless variables to number of useless production is:

a) 1

b) 3/4

c) 2/3

d) 0
\vspace{5pt} \par
\begin{tcolorbox}[colback=green!5, colframe=green!40!black, title=Solution]
Answer: a

Explanation: A, B, C, D are the useless symbols in the given grammar as they never tend to lead to a terminal. The productions S-$>$ A, A-$>$aA, B-$>$C, C-$>$D are also termed as useless production as they will never produce a string to the grammar.
\end{tcolorbox}
\end{mdframed}
\newpage

\begin{mdframed}[linewidth=1pt, roundcorner=5pt, linecolor=gray!30, backgroundcolor=gray!5]
\textbf{Question 5} 

Given grammar G:

S-$>$aS|A|C

A-$>$a

B-$>$aa

C-$>$aCb

Find the set of variables thet can produce strings only with the set of terminals.

a) \{C\}

b) \{A,B\}

c) \{A,B,S\}

d) None of the mentioned
\vspace{5pt} \par
\begin{tcolorbox}[colback=green!5, colframe=green!40!black, title=Solution]
Answer: c

Explanation: First step: Make a set of variables that directly end up with a terminal

             Second step: Modify the set with variables that produce the elements of above

                          generated set.

The rest variables are termed as useless.
\end{tcolorbox}
\end{mdframed}
\newpage

\begin{mdframed}[linewidth=1pt, roundcorner=5pt, linecolor=gray!30, backgroundcolor=gray!5]
\textbf{Question 6} 

Given grammar:

S-$>$aS|A

A-$>$a

B-$>$aa

Find the number of variables reachable from the Starting Variable?

a) 0

b) 1

c) 2

d) None of the mentioned
\vspace{5pt} \par
\begin{tcolorbox}[colback=green!5, colframe=green!40!black, title=Solution]
Answer: b

Explanation: Use a dependency graph to find which variable is reachable and which is not.


 
\begin{center}\includegraphics[width=0.5\textwidth]{automata-theory-assessment-questions-answers-q6.png}\end{center}
\end{tcolorbox}
\end{mdframed}
\newpage

\begin{mdframed}[linewidth=1pt, roundcorner=5pt, linecolor=gray!30, backgroundcolor=gray!5]
\textbf{Question 7} 

Inorder to simplify a context free grammar, we can skip the following operation:

a) Removal of null production

b) Removal of useless symbols

c) Removal of unit productions

d) None of the mentioned
\vspace{5pt} \par
\begin{tcolorbox}[colback=green!5, colframe=green!40!black, title=Solution]
Answer: d

Explanation: Inorder to simplify the grammar all of the process including the removal of null productions, unit productions and useless symbols is necessary.
\end{tcolorbox}
\end{mdframed}
\newpage

\begin{mdframed}[linewidth=1pt, roundcorner=5pt, linecolor=gray!30, backgroundcolor=gray!5]
\textbf{Question 8} 

Given a Grammar G:

S-$>$aA

A-$>$a

A-$>$B

B-$>$A

B-$>$bb

Which among the following will be the simplified grammar?

a) S-$>$aA|aB, A-$>$a, B-$>$bb

b) S-$>$aA|aB, A-$>$B, B-$>$bb

c) S-$>$aA|aB, A-$>$a, B-$>$A

d) None of the emntioned
\vspace{5pt} \par
\begin{tcolorbox}[colback=green!5, colframe=green!40!black, title=Solution]
Answer: a

Explanation: Step 1: Substitute A-$>$B

                     Step 2: Remove B-$>$B

                     Step 3: Substitute B-$>$A

                     Step 4: Remove Repeated productions
\end{tcolorbox}
\end{mdframed}
\newpage

\begin{mdframed}[linewidth=1pt, roundcorner=5pt, linecolor=gray!30, backgroundcolor=gray!5]
\textbf{Question 9} 

Simplify the given grammar:

A-$>$ a| aaA| abBc

B-$>$ abba| b


a) A-$>$ a| aaA| ababbAc| abbc

b) A-$>$ a| aaA| ababbAc| abbc, B-$>$ abba|b

c) A-$>$ a| aaA| abbc, B-$>$abba

d) None of the mentioned
\vspace{5pt} \par
\begin{tcolorbox}[colback=green!5, colframe=green!40!black, title=Solution]
Answer: a

Explanation: Using the substitution rules, we can simply eradicate what is useless and thus produce the simplified result i.e. A-$>$ a| aaA| ababbAc| abbc.
\end{tcolorbox}
\end{mdframed}
\newpage

\begin{mdframed}[linewidth=1pt, roundcorner=5pt, linecolor=gray!30, backgroundcolor=gray!5]
\textbf{Question 10} 

In context to the process of removing useless symbols, which of the following is correct?

a) We remove the Nullable variables

b) We eliminate the unit productions

c) We eliminate products which yield no terminals

d) All of the mentioned
\vspace{5pt} \par
\begin{tcolorbox}[colback=green!5, colframe=green!40!black, title=Solution]
Answer: c

Explanation: In the  process of removal of useless symbols, we want to remove productions that can never take part in any derivation.
\end{tcolorbox}
\end{mdframed}
\newpage

\newtopic{Eliminating Epsilon Productions}

\begin{mdframed}[linewidth=1pt, roundcorner=5pt, linecolor=gray!30, backgroundcolor=gray!5]
\textbf{Question 1} 

The use of variable dependency graph is in:

a) Removal of useless variables

b) Removal of null productions

c) Removal of unit productions

d) None of the mentioned
\vspace{5pt} \par
\begin{tcolorbox}[colback=green!5, colframe=green!40!black, title=Solution]
Answer: a

Explanation: We use the concept of dependency graph inorder to check, whether any of the variable is reachable from the starting variable or not.
\end{tcolorbox}
\end{mdframed}
\newpage

\begin{mdframed}[linewidth=1pt, roundcorner=5pt, linecolor=gray!30, backgroundcolor=gray!5]
\textbf{Question 2} 

The variable which produces an epsilon is called:

a) empty variable

b) nullable

c) terminal

d) all of the mentioned
\vspace{5pt} \par
\begin{tcolorbox}[colback=green!5, colframe=green!40!black, title=Solution]
Answer: b

Explanation: Any variable A for which the derivation: A-$>$*e is possible is called Nullable.
\end{tcolorbox}
\end{mdframed}
\newpage

\begin{mdframed}[linewidth=1pt, roundcorner=5pt, linecolor=gray!30, backgroundcolor=gray!5]
\textbf{Question 3} 

Statement:

For A-$>$ e ,A can be erased. So whenever it appears on the left side of a production, replace with another production without the A.

State true or false:

a) true

b) false
\vspace{5pt} \par
\begin{tcolorbox}[colback=green!5, colframe=green!40!black, title=Solution]
Answer: b

Explanation:  A can be erased. So whenever it appears on the right side of the production, replace with another production without the A.
\end{tcolorbox}
\end{mdframed}
\newpage

\begin{mdframed}[linewidth=1pt, roundcorner=5pt, linecolor=gray!30, backgroundcolor=gray!5]
\textbf{Question 4} 

Simplify the given grammar:

S-$>$aXb

X-$>$aXb | e

a) S-$>$aXb | ab, X-$>$ aXb | ab

b) S-$>$X | ab, X-$>$ aXb | ab

c) S-$>$aXb | ab, X-$>$ S | ab

d) None of the mentioned
\vspace{5pt} \par
\begin{tcolorbox}[colback=green!5, colframe=green!40!black, title=Solution]
Answer: a

Explanation: As X is nullable, we replace every right hand side presence of X with e and produce the simplified result.
\end{tcolorbox}
\end{mdframed}
\newpage

\begin{mdframed}[linewidth=1pt, roundcorner=5pt, linecolor=gray!30, backgroundcolor=gray!5]
\textbf{Question 5} 

Consider the following grammar:

A-$>$e

B-$>$aAbC

B-$>$bAbA

A-$>$bB

The number of productions added on the removal of the nullable in the given grammar:

a) 3

b) 4

c) 2

d) 0
\vspace{5pt} \par
\begin{tcolorbox}[colback=green!5, colframe=green!40!black, title=Solution]
Answer: b

Explanation: The modified grammar aftyer the removal of nullable can be shown as:

B-$>$aAbC| abC

B-$>$bAbA| bbA|  bAb| bb

A-$>$bB
\end{tcolorbox}
\end{mdframed}
\newpage

\begin{mdframed}[linewidth=1pt, roundcorner=5pt, linecolor=gray!30, backgroundcolor=gray!5]
\textbf{Question 6} 

Let G=(V, T, P, S) be a CFG such that \_\_\_\_\_\_\_\_\_\_\_\_\_. Then there exists an equivalent grammar G' having no e productions.

a) e $\in$ L(G)

b) w $\notin$ L(G)

c) e $\notin$ L(G)

d) w $\in$ L(G)
\vspace{5pt} \par
\begin{tcolorbox}[colback=green!5, colframe=green!40!black, title=Solution]
Answer: c

Explanation: Theorem: Let G = (V, T, S, P) be a CFG such that  e $\notin$ L(G). Then there exists an equivalent grammar G' having no e-productions.
\end{tcolorbox}
\end{mdframed}
\newpage

\begin{mdframed}[linewidth=1pt, roundcorner=5pt, linecolor=gray!30, backgroundcolor=gray!5]
\textbf{Question 7} 

For each production in P of the form:

A-$>$ x1x2x3…xn

put into P' that production as well as all those generated by replacing null variables with e in all possible combinations. If all x(i) are nullable,

a) A-$>$e is put into P'

b) A-$>$e is not put into P'

c) e is a member of G'

d) None of the mentioned
\vspace{5pt} \par
\begin{tcolorbox}[colback=green!5, colframe=green!40!black, title=Solution]
Answer: b

Explanation: It is an exception that A-$>$e is not put into P' if all x(i) are nullable variables.
\end{tcolorbox}
\end{mdframed}
\newpage

\begin{mdframed}[linewidth=1pt, roundcorner=5pt, linecolor=gray!30, backgroundcolor=gray!5]
\textbf{Question 8} 

For the given grammar G:

S-$>$ABaC

A-$>$BC

B-$>$b| e

C-$>$D| e

D-$>$ d

Remove the e productions and generate the number of productions from S in the modified or simplified grammar.

a) 6

b) 7

c) 5

d) 8
\vspace{5pt} \par
\begin{tcolorbox}[colback=green!5, colframe=green!40!black, title=Solution]
Answer: d

Explanation: The grammar after the removal of epsilon production can be shown as:

S-$>$ABaC| AaC| ABa| Aa| a| aC| Ba| BaC

A-$>$BC| B| C

B-$>$b

C-$>$D

D-$>$ d
\end{tcolorbox}
\end{mdframed}
\newpage

\begin{mdframed}[linewidth=1pt, roundcorner=5pt, linecolor=gray!30, backgroundcolor=gray!5]
\textbf{Question 9} 

Consider G=(\{S,A,B,E\}, \{a,b,c\},P,S), where P consists of S →AB, A →a, B →b and E →c.

Number of productions in P' after removal of useless symbols:

a) 4

b) 3

c) 2

d) 5
\vspace{5pt} \par
\begin{tcolorbox}[colback=green!5, colframe=green!40!black, title=Solution]
Answer: a

Explanation:

P'= S-$>$AB, A-$>$a, B-$>$ b,

V'=\{S, A, B\},

$\Sigma$'=\{a, b\}
\end{tcolorbox}
\end{mdframed}
\newpage

\begin{mdframed}[linewidth=1pt, roundcorner=5pt, linecolor=gray!30, backgroundcolor=gray!5]
\textbf{Question 10} 

Given grammar G:

S-$>$aS| AB

A-$>$ e

B-$>$ e

D-$>$ b

Reduce the grammar, removing all the e productions:

a) S-$>$aS| AB| A| B, D-$>$ b

b) S-$>$aS| AB| A| B| a, D-$>$ b

c) S-$>$aS| AB| A| B

d) None of the mentioned
\vspace{5pt} \par
\begin{tcolorbox}[colback=green!5, colframe=green!40!black, title=Solution]
Answer: b

Explanation: We will replace all the nullables wherever they appear in the right hand side of any production. D will not be erased as we are just removing nullable variables not completely simplifying the grammar.
\end{tcolorbox}
\end{mdframed}
\newpage

\newtopic{Eliminating Unit Productions}

\begin{mdframed}[linewidth=1pt, roundcorner=5pt, linecolor=gray!30, backgroundcolor=gray!5]
\textbf{Question 1} 

Which of the following is the format of unit production?

a) A-$>$B

b) A-$>$b

c) B-$>$Aa

d) None of the mentioned
\vspace{5pt} \par
\begin{tcolorbox}[colback=green!5, colframe=green!40!black, title=Solution]
Answer: a

Explanation: Any production of the format A-$>$ B where A and B belongs to the V set, is called Unit production.
\end{tcolorbox}
\end{mdframed}
\newpage

\begin{mdframed}[linewidth=1pt, roundcorner=5pt, linecolor=gray!30, backgroundcolor=gray!5]
\textbf{Question 2} 

Given Grammar G:

S-$>$aA

A-$>$a| A

B-$>$B

The number of productions to be removed immediately as Unit productions:

a) 0

b) 1

c) 2

d) 3
\vspace{5pt} \par
\begin{tcolorbox}[colback=green!5, colframe=green!40!black, title=Solution]
Answer: c

Explanation: The productions in the format A-$>$ A are removed immediately as they produce self and that is not a terminal or will not lead to a string. Hence, it is removed immediately.
\end{tcolorbox}
\end{mdframed}
\newpage

\begin{mdframed}[linewidth=1pt, roundcorner=5pt, linecolor=gray!30, backgroundcolor=gray!5]
\textbf{Question 3} 

Given grammar:

S-$>$aA

A-$>$a

A-$>$B

B-$>$ A

B-$>$bb

Which of the following is the production of B after simplification by removal of unit productions?

a) A

b) bb

c) aA

d) A| bb
\vspace{5pt} \par
\begin{tcolorbox}[colback=green!5, colframe=green!40!black, title=Solution]
Answer: b

Explanation: The simplified grammar can be presented as follows:

S-$>$aA| aB

A-$>$a

B-$>$ bb
\end{tcolorbox}
\end{mdframed}
\newpage

\begin{mdframed}[linewidth=1pt, roundcorner=5pt, linecolor=gray!30, backgroundcolor=gray!5]
\textbf{Question 4} 

If grammar G is unambiguous, G' produced after the removal of Unit production will be:

a) ambiguous

b) unambiguous

c) finite

d) cannot be said
\vspace{5pt} \par
\begin{tcolorbox}[colback=green!5, colframe=green!40!black, title=Solution]
Answer: b

Explanation: With the simplification of Context free grammars, undesirable properties are introduced. It says, if grammar G, before simplification is unambiguous, after simplification will also be unambiguous.
\end{tcolorbox}
\end{mdframed}
\newpage

\begin{mdframed}[linewidth=1pt, roundcorner=5pt, linecolor=gray!30, backgroundcolor=gray!5]
\textbf{Question 5} 

If C is A-derivable, C-$>$B is a production, and B ¹ A, then B is

a) nullable

b) Non-derivable

c) A-derivable

d) None of the mentioned
\vspace{5pt} \par
\begin{tcolorbox}[colback=green!5, colframe=green!40!black, title=Solution]
Answer: c

Explanation:

If A-$>$ B is a production, B is called A- derivable.

If C is A-derivable, C-$>$B is a production, and B ¹ A, then B is A -derivable.

No other variables are A-derivable.
\end{tcolorbox}
\end{mdframed}
\newpage

\begin{mdframed}[linewidth=1pt, roundcorner=5pt, linecolor=gray!30, backgroundcolor=gray!5]
\textbf{Question 6} 

A can be A-$>$ derivable if and only if \_\_\_\_\_\_\_\_\_\_

a) A-$>$ A is actually a production

b) A-$>$B, B-$>$ A exists

c) All of the mentioned

d) None of the mentioned
\vspace{5pt} \par
\begin{tcolorbox}[colback=green!5, colframe=green!40!black, title=Solution]
Answer: a

Explanation: The format says: If A-$>$B is a production, B is called A-derivable.Thus A to be A-derivable, a production : A-$>$ A need to exist.
\end{tcolorbox}
\end{mdframed}
\newpage

\begin{mdframed}[linewidth=1pt, roundcorner=5pt, linecolor=gray!30, backgroundcolor=gray!5]
\textbf{Question 7} 

Given Grammar:

T-$>$ T+R| R

R-$>$ R*V| V

V-$>$(T)| u

When unit productions are deleted we are left with

T-$>$ T+R| \_\_\_\_\_\_\_|(T)| u

R-$>$R*V|(T)| u

V-$>$ (T)| u

Fill in the blank:

a) T*V

b) T+V

c) R*T

d) R*V
\vspace{5pt} \par
\begin{tcolorbox}[colback=green!5, colframe=green!40!black, title=Solution]
Answer: d

Explanation: The grammar produced after the elimination of unit production is:

T-$>$ T+R| R*V| (T)| u

R-$>$R*V|(T)| u

V-$>$ (T)| u
\end{tcolorbox}
\end{mdframed}
\newpage

\begin{mdframed}[linewidth=1pt, roundcorner=5pt, linecolor=gray!30, backgroundcolor=gray!5]
\textbf{Question 8} 

Given grammar G:

S-$>$ ABA, A-$>$aA|e, B-$>$ bB|e

Eliminate e and unit productions. State the number of productions the starting variable holds?

a) 6

b) 7

c) 9

d) 5
\vspace{5pt} \par
\begin{tcolorbox}[colback=green!5, colframe=green!40!black, title=Solution]
Answer: b

Explanation: After reduction the grammar looks like:

S-$>$ABA| AB| BA| AA| Aa| a| bB| b

A-$>$aA| a

B-$>$bB| b
\end{tcolorbox}
\end{mdframed}
\newpage

\begin{mdframed}[linewidth=1pt, roundcorner=5pt, linecolor=gray!30, backgroundcolor=gray!5]
\textbf{Question 9} 

Given grammar G:

S-$>$ A| B| C

A-$>$ aAa| B

B-$>$ bB|bb

C-$>$aCaa|D

D-$>$baD|abD|aa

Eliminate e and unit productions and state the number of variables left?

a) 5

b) 4

c) 3

d) 2
\vspace{5pt} \par
\begin{tcolorbox}[colback=green!5, colframe=green!40!black, title=Solution]
Answer: a

Explanation: The reduced production:

 S-$>$aAa| bB|bb aCaa| baD| abD| aa, A-$>$aAa| bB| bb, B-$>$bB| bb, C-$>$aCaa| baD| abD| aa, D-$>$ baD| abD| aa
\end{tcolorbox}
\end{mdframed}
\newpage

\begin{mdframed}[linewidth=1pt, roundcorner=5pt, linecolor=gray!30, backgroundcolor=gray!5]
\textbf{Question 10} 

Which of the following variables in the given grammar is called live variable?

S-$>$AB

A-$>$a

a) S

b) A

c) B

d) None of the mentioned
\vspace{5pt} \par
\begin{tcolorbox}[colback=green!5, colframe=green!40!black, title=Solution]
Answer: b

Explanation: Any variable A for which there is a production A-$>$ x with x E $\Sigma$* is called live.
\end{tcolorbox}
\end{mdframed}
\newpage

\newtopic{Chomsky Normal Form}

\begin{mdframed}[linewidth=1pt, roundcorner=5pt, linecolor=gray!30, backgroundcolor=gray!5]
\textbf{Question 1} 

The format: A-$>$aB refers to which of the following?

a) Chomsky Normal Form

b) Greibach Normal Form

c) Backus Naur Form

d) None of the mentioned
\vspace{5pt} \par
\begin{tcolorbox}[colback=green!5, colframe=green!40!black, title=Solution]
Answer: b

Explanation: A context free grammar is in Greibach Normal Form if the right hand sides of all the production rules start with a terminal, optionally followed by some variables.
\end{tcolorbox}
\end{mdframed}
\newpage

\begin{mdframed}[linewidth=1pt, roundcorner=5pt, linecolor=gray!30, backgroundcolor=gray!5]
\textbf{Question 2} 

Which of the following does not have left recursions?

a) Chomsky Normal Form

b) Greibach Normal Form

c) Backus Naur Form

d) All  of the mentioned
\vspace{5pt} \par
\begin{tcolorbox}[colback=green!5, colframe=green!40!black, title=Solution]
Answer: b

Explanation: The normal form is of the format:

A-$>$aB where the right hand side production tends to begin with a terminal symbo, thus having no left recursions.
\end{tcolorbox}
\end{mdframed}
\newpage

\begin{mdframed}[linewidth=1pt, roundcorner=5pt, linecolor=gray!30, backgroundcolor=gray!5]
\textbf{Question 3} 

Every grammar in Chomsky Normal Form is:

a) regular

b) context sensitive

c) context free

d) all of the mentioned
\vspace{5pt} \par
\begin{tcolorbox}[colback=green!5, colframe=green!40!black, title=Solution]
Answer: c

Explanation: Conversely, every context frr grammar can be converted into Chomsky Normal form and to other forms.
\end{tcolorbox}
\end{mdframed}
\newpage

\begin{mdframed}[linewidth=1pt, roundcorner=5pt, linecolor=gray!30, backgroundcolor=gray!5]
\textbf{Question 4} 

Which of the production rule can be accepted by Chomsky grammar?

a) A-$>$BC

b) A-$>$a

c) S-$>$e

d) All of the mentioned
\vspace{5pt} \par
\begin{tcolorbox}[colback=green!5, colframe=green!40!black, title=Solution]
Answer: d

Explanation: in CNF, the production rules are of the form:

A-$>$BC

A-$>$ a

S-$>$e
\end{tcolorbox}
\end{mdframed}
\newpage

\begin{mdframed}[linewidth=1pt, roundcorner=5pt, linecolor=gray!30, backgroundcolor=gray!5]
\textbf{Question 5} 

Given grammar G:

(1)S-$>$AS

(2)S-$>$AAS

(3)A-$>$SA

(4)A-$>$aa

Which of the following productions denies the format of Chomsky Normal Form?

a) 2,4

b) 1,3

c) 1, 2, 3, 4

d) 2, 3, 4
\vspace{5pt} \par
\begin{tcolorbox}[colback=green!5, colframe=green!40!black, title=Solution]
Answer: a

Explanation: The correct format: A-$>$BC, A-$>$a, X-$>$e.
\end{tcolorbox}
\end{mdframed}
\newpage

\begin{mdframed}[linewidth=1pt, roundcorner=5pt, linecolor=gray!30, backgroundcolor=gray!5]
\textbf{Question 6} 

Which of the following grammars are in Chomsky Normal Form:

a) S-$>$AB|BC|CD, A-$>$0, B-$>$1, C-$>$2, D-$>$3

b) S-$>$AB, S-$>$BCA|0|1|2|3

c) S-$>$ABa, A-$>$aab, B-$>$Ac

d) All of the mentioned
\vspace{5pt} \par
\begin{tcolorbox}[colback=green!5, colframe=green!40!black, title=Solution]
Answer: a

Explanation: We can eliminate the options on the basis of the format we are aware of: A-$>$BC, B-$>$b and so on.
\end{tcolorbox}
\end{mdframed}
\newpage

\begin{mdframed}[linewidth=1pt, roundcorner=5pt, linecolor=gray!30, backgroundcolor=gray!5]
\textbf{Question 7} 

With reference to the process of conversion of a context free grammar to CNF, the number of variables to be introduced for the terminals are:

S-$>$ABa

A-$>$aab

B-$>$Ac

a) 3

b) 4

c) 2

d) 5
\vspace{5pt} \par
\begin{tcolorbox}[colback=green!5, colframe=green!40!black, title=Solution]
Answer: a

Explanation: According to the number of terminals present in the grammar, we need the corresponding that number of terminal variables while conversion.
\end{tcolorbox}
\end{mdframed}
\newpage

\begin{mdframed}[linewidth=1pt, roundcorner=5pt, linecolor=gray!30, backgroundcolor=gray!5]
\textbf{Question 8} 

In which of the following, does the CNF conversion find its use?

a) CYK Algorithm

b) Bottom up parsing

c) Preprocessing step in some algorithms

d) All of the mentioned
\vspace{5pt} \par
\begin{tcolorbox}[colback=green!5, colframe=green!40!black, title=Solution]
Answer: d

Explanation: Besides the theoretical significance of CNF, it conversion scheme is helpful in algorithms as a preprocessing step, CYK algorithms and the bottom up parsing of context free grammars.
\end{tcolorbox}
\end{mdframed}
\newpage

\begin{mdframed}[linewidth=1pt, roundcorner=5pt, linecolor=gray!30, backgroundcolor=gray!5]
\textbf{Question 9} 

Let G be a grammar. When the production in G satisfy certain restrictions, then G is said to be in \_\_\_\_\_\_\_\_\_\_\_

a) restricted form

b) parsed form

c) normal form

d) all of the mentioned
\vspace{5pt} \par
\begin{tcolorbox}[colback=green!5, colframe=green!40!black, title=Solution]
Answer: c

Explanation: When the production in G satisfy certain restrictions, then G is said to be in ‘normal form'.
\end{tcolorbox}
\end{mdframed}
\newpage

\begin{mdframed}[linewidth=1pt, roundcorner=5pt, linecolor=gray!30, backgroundcolor=gray!5]
\textbf{Question 10} 

Let G be a grammar: S-$>$AB|e, A-$>$a, B-$>$b

Is the given grammar in CNF?

a) Yes

b) No
\vspace{5pt} \par
\begin{tcolorbox}[colback=green!5, colframe=green!40!black, title=Solution]
Answer: a

Explanation: e is allowed in CNF only if the starting variable does not occur on the right hand side of the derivation.
\end{tcolorbox}
\end{mdframed}
\newpage

\newtopic{Pumping Lemma for Context Free Language}

\begin{mdframed}[linewidth=1pt, roundcorner=5pt, linecolor=gray!30, backgroundcolor=gray!5]
\textbf{Question 1} 

Which of the following is called Bar-Hillel lemma?

a) Pumping lemma for regular language

b) Pumping lemma for context free languages

c) Pumping lemma for context sensitive languages

d) None of the mentioned
\vspace{5pt} \par
\begin{tcolorbox}[colback=green!5, colframe=green!40!black, title=Solution]
Answer: b

Explanation: In automata theory, the pumping lemma for context free languages, also kmown as the  Bar-Hillel lemma, represents a property of all context free languages.
\end{tcolorbox}
\end{mdframed}
\newpage

\begin{mdframed}[linewidth=1pt, roundcorner=5pt, linecolor=gray!30, backgroundcolor=gray!5]
\textbf{Question 2} 

Which of the expressions correctly is an requirement of the pumping lemma for the context free languages?

a) uv
n
wx
n
y

b) uv
n
w
n
x
n
y

c) uv
2n
wx
2n
y

d) All of the mentioned
\vspace{5pt} \par
\begin{tcolorbox}[colback=green!5, colframe=green!40!black, title=Solution]
Answer: b

Explanation: Let L be a CFL. Then there is an integer n so that for any u that belong to language L satisfying |t| $>$=n, there are strings u, v, w, x, y and z satisfying

t=uvwxy

|vx|$>$0

|vwx|$<$=n
For any m$>$=0, uv
n
wx
n
y $\in$ L
\end{tcolorbox}
\end{mdframed}
\newpage

\begin{mdframed}[linewidth=1pt, roundcorner=5pt, linecolor=gray!30, backgroundcolor=gray!5]
\textbf{Question 3} 

Let L be a CFL. Then there is an integer n so that for any u that belong to language L satisfying

|t|$>$=n, there are strings u, v, w, x, y and z satisfying

t=uvwxy.

Let p be the number of variables in CNF form of the context free grammar. The value of n in terms of p is

a) 2p

b) 2p

c) 2p+1

d) p
2
\vspace{5pt} \par
\begin{tcolorbox}[colback=green!5, colframe=green!40!black, title=Solution]
Answer: c

Explanation: This inequation has been derived from derivation tree for t which must have height at least p+2(It has more than 2
p
 leaf nodes, and therefore its height is $>$p+1).
\end{tcolorbox}
\end{mdframed}
\newpage

\begin{mdframed}[linewidth=1pt, roundcorner=5pt, linecolor=gray!30, backgroundcolor=gray!5]
\textbf{Question 4} 

Which of the following gives a positive result to the pumping lemma restrictions and requirements?

a) \{a
i
b
i
c
i
|i$>$=0\}

b) \{0
i
1
i
|i$>$=0\}

c) \{ss|s$\in$\{a,b\}*\}

d) None of the mentioned
\vspace{5pt} \par
\begin{tcolorbox}[colback=green!5, colframe=green!40!black, title=Solution]
Answer: b

Explanation: A positive result to the pumping lemma shows that the language is a CFL and ist contradiction or negative result shows that the given language is not a Context Free language.
\end{tcolorbox}
\end{mdframed}
\newpage

\begin{mdframed}[linewidth=1pt, roundcorner=5pt, linecolor=gray!30, backgroundcolor=gray!5]
\textbf{Question 5} 

Using pumping lemma, which of the following cannot be proved as ‘not a CFL'?

a) \{a
i
b
i
c
i
|i$>$=0\}

b) \{ss|s$\in$\{a,b\}*\}

c) The set legal C programs

d) None of the mentioned
\vspace{5pt} \par
\begin{tcolorbox}[colback=green!5, colframe=green!40!black, title=Solution]
Answer: d

Explanation: There are few rules in C that are context dependent. For example, declaration of a variable before it can be used.
\end{tcolorbox}
\end{mdframed}
\newpage

\begin{mdframed}[linewidth=1pt, roundcorner=5pt, linecolor=gray!30, backgroundcolor=gray!5]
\textbf{Question 6} 

State true or false:

Statement: We cannot use Ogden's lemma when pumping lemma fails.

a) true

b) false
\vspace{5pt} \par
\begin{tcolorbox}[colback=green!5, colframe=green!40!black, title=Solution]
Answer: b

Explanation: Although the pumping lemma provides some information about v and x that are pumped, it says little about the location of these substrings in the string t. It can be used whenever the pumping lemma fails. Example: \{a
p
b
q
c
r
d
s
|p=0 or q=r=s\}, etc.
\end{tcolorbox}
\end{mdframed}
\newpage

\begin{mdframed}[linewidth=1pt, roundcorner=5pt, linecolor=gray!30, backgroundcolor=gray!5]
\textbf{Question 7} 

Which of the following cannot be filled in the blank below?

Statement: There are CFLs L1 nad L2 so that \_\_\_\_\_\_\_\_\_\_\_is not a CFL.

a) L1$\cap$L2

b) L1$^\prime$

c) L1*

d) None of the mentioned
\vspace{5pt} \par
\begin{tcolorbox}[colback=green!5, colframe=green!40!black, title=Solution]
Answer: c

Explanation: A set of context free language is closed under the following operations:

a) Union

b) Concatenation

c) Kleene
\end{tcolorbox}
\end{mdframed}
\newpage

\begin{mdframed}[linewidth=1pt, roundcorner=5pt, linecolor=gray!30, backgroundcolor=gray!5]
\textbf{Question 8} 

The pumping lemma is often used to prove that a language is:

a) Context free

b) Not context free

c) Regular

d) None of the mentioned
\vspace{5pt} \par
\begin{tcolorbox}[colback=green!5, colframe=green!40!black, title=Solution]
Answer: b

Explanation: The pumping lemma is often used to prove that a given language L is non-context-free, by showing that arbitrarily long strings s are in L that cannot be "pumped" without producing strings outside L.
\end{tcolorbox}
\end{mdframed}
\newpage

\begin{mdframed}[linewidth=1pt, roundcorner=5pt, linecolor=gray!30, backgroundcolor=gray!5]
\textbf{Question 9} 

What is the pumping length of string of length x?

a) x+1

b) x

c) x-1

d) x2
\vspace{5pt} \par
\begin{tcolorbox}[colback=green!5, colframe=green!40!black, title=Solution]
Answer: a

Explanation: There exists a property of all strings in the language that are of length p, where p is the constant-called the pumping length .For a finite language L, p is equal to the maximum string length  in L plus 1.
\end{tcolorbox}
\end{mdframed}
\newpage

\begin{mdframed}[linewidth=1pt, roundcorner=5pt, linecolor=gray!30, backgroundcolor=gray!5]
\textbf{Question 10} 

Which of the following does not obey pumping lemma for context free languages ?

a) Finite languages

b) Context free languages

c) Unrestricted languages

d) None of the mentioned
\vspace{5pt} \par
\begin{tcolorbox}[colback=green!5, colframe=green!40!black, title=Solution]
Answer: c

Explanation: Finite languages (which are regular hence context free ) obey pumping lemma where as unrestricted languages like recursive languages do not obey pumping lemma for context free languages.
\end{tcolorbox}
\end{mdframed}
\newpage

\newtopic{CFL- Closure Properties}

\begin{mdframed}[linewidth=1pt, roundcorner=5pt, linecolor=gray!30, backgroundcolor=gray!5]
\textbf{Question 1} 

The context free languages are closed under:

a) Intersection

b) Complement

c) Kleene

d) None of the mentioned
\vspace{5pt} \par
\begin{tcolorbox}[colback=green!5, colframe=green!40!black, title=Solution]
Answer: c

Explanation: Context free languages are closed under the following operation: union, kleene and concatenation. For regular languages, we can add intersection and complement to the list.
\end{tcolorbox}
\end{mdframed}
\newpage

\begin{mdframed}[linewidth=1pt, roundcorner=5pt, linecolor=gray!30, backgroundcolor=gray!5]
\textbf{Question 2} 

Given Grammar G1:

S-$>$aSb

S-$>$e

Grammar G2:

R-$>$cRd

R-$>$e

If L(G)=L(G1) U L(G2), the number of productions the new starting variable would have:

a) 2

b) 3

c) 4

d) 1
\vspace{5pt} \par
\begin{tcolorbox}[colback=green!5, colframe=green!40!black, title=Solution]
Answer: a

Explanation:

T-$>$S|R

S-$>$aSb

S-$>$e

R-$>$cRd

R-$>$e
\end{tcolorbox}
\end{mdframed}
\newpage

\begin{mdframed}[linewidth=1pt, roundcorner=5pt, linecolor=gray!30, backgroundcolor=gray!5]
\textbf{Question 3} 

Context free languages are not closed under:

a) Intersection

b) Intersection with Regular Language

c) Complement

d) All of the mentioned
\vspace{5pt} \par
\begin{tcolorbox}[colback=green!5, colframe=green!40!black, title=Solution]
Answer: d

Explanation: It is a theorem which states that, Context free languages are not closed under operations like intersection and complement.
\end{tcolorbox}
\end{mdframed}
\newpage

\begin{mdframed}[linewidth=1pt, roundcorner=5pt, linecolor=gray!30, backgroundcolor=gray!5]
\textbf{Question 4} 

Which of the following is incorrect?

There exists algorithms to decide if:

a) String w is in CFL L

b) CFL L is empty

c) CFL L is infinite

d) All of the mentioned
\vspace{5pt} \par
\begin{tcolorbox}[colback=green!5, colframe=green!40!black, title=Solution]
Answer: d

Explanation: These properties are termed as decision properties of a CFL and include a set of problems like infiniteness problem, emptiness problem and membership problem.
\end{tcolorbox}
\end{mdframed}
\newpage

\begin{mdframed}[linewidth=1pt, roundcorner=5pt, linecolor=gray!30, backgroundcolor=gray!5]
\textbf{Question 5} 

If the start symbol is one of those symbols which produce no terminal through any sequence, the CFL is said to be

a) nullable

b) empty

c) eliminated

d) none of the mentioned
\vspace{5pt} \par
\begin{tcolorbox}[colback=green!5, colframe=green!40!black, title=Solution]
Answer: b

Explanation: In the process of removing useless symbols, if the starting symbol is also a part, the CFL can be then termed as empty; otherwise not.
\end{tcolorbox}
\end{mdframed}
\newpage

\begin{mdframed}[linewidth=1pt, roundcorner=5pt, linecolor=gray!30, backgroundcolor=gray!5]
\textbf{Question 6} 

Using the pumping constant n, If there is a string in the language of length between \_\_\_\_\_ and \_\_\_\_ then the language is infite else not.

a) n, 2n-1

b) 2n, n

c) n+1, 3n+6

d) 0, n+1
\vspace{5pt} \par
\begin{tcolorbox}[colback=green!5, colframe=green!40!black, title=Solution]
Answer: a

Explanation: If there is a string in the language of length between n and 2n-1 then the language is infite else not. The idea is essentially the same for regular languages.
\end{tcolorbox}
\end{mdframed}
\newpage

\begin{mdframed}[linewidth=1pt, roundcorner=5pt, linecolor=gray!30, backgroundcolor=gray!5]
\textbf{Question 7} 

Which of the following is/are CFL not closed under?

a) Reverse

b) Homomorphism

c) Inverse Homomorphism

d) All of the mentioned
\vspace{5pt} \par
\begin{tcolorbox}[colback=green!5, colframe=green!40!black, title=Solution]
Answer: d

Explanation: CFL is closed under union, kleene and concatenation along with the properties reversal,homomorphism and inverse homomorphism but not difference and intersection.
\end{tcolorbox}
\end{mdframed}
\newpage

\begin{mdframed}[linewidth=1pt, roundcorner=5pt, linecolor=gray!30, backgroundcolor=gray!5]
\textbf{Question 8} 

If L1 and L2 are context free languages, L1-L2 are context free:

a) always

b) sometimes

c) never

d) none of the mentioned
\vspace{5pt} \par
\begin{tcolorbox}[colback=green!5, colframe=green!40!black, title=Solution]
Answer: c

Explanation: Context free languages are not closed under difference, intersection and complement operations.
\end{tcolorbox}
\end{mdframed}
\newpage

\begin{mdframed}[linewidth=1pt, roundcorner=5pt, linecolor=gray!30, backgroundcolor=gray!5]
\textbf{Question 9} 

A\_\_\_\_\_\_\_\_\_\_\_ is context free grammar with atmost one non terminal in the right handside of the production.

a) linear grammar

b) linear bounded grammar

c) regular grammar

d) none of the mentioned
\vspace{5pt} \par
\begin{tcolorbox}[colback=green!5, colframe=green!40!black, title=Solution]
Answer: a

Explanation: A simple linear grammar is G with N = \{S\}, $\Sigma$ = \{a, b\}, P with start symbol S and rules

S → aSb

S → $\epsilon$
\end{tcolorbox}
\end{mdframed}
\newpage

\begin{mdframed}[linewidth=1pt, roundcorner=5pt, linecolor=gray!30, backgroundcolor=gray!5]
\textbf{Question 10} 

There is a linear grammar that generates a context free grammar

a) always

b) never

c) sometimes

d) none of the mentioned
\vspace{5pt} \par
\begin{tcolorbox}[colback=green!5, colframe=green!40!black, title=Solution]
Answer: c

Explanation: Linear grammar is a subset of context free grammar which has atmost one non terminal symbol in the right hand side of the production.Thus, there exists some languages which are generated by Linear grammars.
\end{tcolorbox}
\end{mdframed}
\newpage

\newtopic{CFL- Other Normal Forms}

\begin{mdframed}[linewidth=1pt, roundcorner=5pt, linecolor=gray!30, backgroundcolor=gray!5]
\textbf{Question 1} 

The following format of grammatical notation is accepted by which of the following:

AB-$>$CD

A-$>$BC or

A-$>$B or

A-$>$a

where A, B, C, D are non terminal symbols and a is a terminal symbol.

a) Greibach Normal Form

b) Chomsky Nrmal Form

c) Kuroda Normal Form

d) None of the mentioned
\vspace{5pt} \par
\begin{tcolorbox}[colback=green!5, colframe=green!40!black, title=Solution]
Answer: c

Explanation: Linearly Bounded grammar or Kuroda Normal Form allows the following format of grammatical analysis:

AB-$>$CD

A-$>$BC or

A-$>$B or

A-$>$a
\end{tcolorbox}
\end{mdframed}
\newpage

\begin{mdframed}[linewidth=1pt, roundcorner=5pt, linecolor=gray!30, backgroundcolor=gray!5]
\textbf{Question 2} 

Every Kuroda Normal form grammar generates \_\_\_\_\_\_\_\_\_\_\_

a) Context free grammar

b) Context sensitive grammar

c) Unrestricted grammar

d) None of the mentioned
\vspace{5pt} \par
\begin{tcolorbox}[colback=green!5, colframe=green!40!black, title=Solution]
Answer: b

Explanation: Every context sensitive grammar which does not produce an empty string can be generated by a grammar in Kuroda Normal form.
\end{tcolorbox}
\end{mdframed}
\newpage

\begin{mdframed}[linewidth=1pt, roundcorner=5pt, linecolor=gray!30, backgroundcolor=gray!5]
\textbf{Question 3} 

Which of the following can generate Unrestricted grammars?

a) Pentonnen Normal form

b) Floyd Normal form

c) Greibach Normal form

d) None of the mentioned
\vspace{5pt} \par
\begin{tcolorbox}[colback=green!5, colframe=green!40!black, title=Solution]
Answer: a

Explanation: Pentonnen Normal form(for Unrestricted grammars) is a special case where there is a slight modification in the format of Kuroda Normal form.

AB-$>$AD

A-$>$BC

A-$>$a
\end{tcolorbox}
\end{mdframed}
\newpage

\begin{mdframed}[linewidth=1pt, roundcorner=5pt, linecolor=gray!30, backgroundcolor=gray!5]
\textbf{Question 4} 

Given a grammar in GNF and a derivable string in the grammar with the length n, any \_\_\_\_\_\_\_\_\_\_\_will halt at depth n.

a) top-down parser

b) bottom-up parser

c) multitape turing machine

d) none of the mentioned
\vspace{5pt} \par
\begin{tcolorbox}[colback=green!5, colframe=green!40!black, title=Solution]
Answer: a

Explanation: Given a grammar in GNF and a derivable string in the grammar with the length n, any top-down parser will halt at depth n. As the parameter ‘depth' is mentioned, we will use a top-down parser. Example-LL parser.
\end{tcolorbox}
\end{mdframed}
\newpage

\begin{mdframed}[linewidth=1pt, roundcorner=5pt, linecolor=gray!30, backgroundcolor=gray!5]
\textbf{Question 5} 

Which of the following grammars is similar to Floyd Normal form?

a) Backus Naur Form

b) Kuroda Normal Form

c) Greibach Normal Form

d) Chomsky Normal Form
\vspace{5pt} \par
\begin{tcolorbox}[colback=green!5, colframe=green!40!black, title=Solution]
Answer: a

Explanation: Donald Knuth implied a BNF" syntax in which all definitions have such a form may be said to be in "Floyd Normal Form".

A-$>$B|C

A-$>$BC

A-$>$a
\end{tcolorbox}
\end{mdframed}
\newpage

\begin{mdframed}[linewidth=1pt, roundcorner=5pt, linecolor=gray!30, backgroundcolor=gray!5]
\textbf{Question 6} 

Which among the following can parse a context free grammar?

a) top down parser

b) bottom up parser

c) CYK algorithm

d) all of the mentioned
\vspace{5pt} \par
\begin{tcolorbox}[colback=green!5, colframe=green!40!black, title=Solution]
Answer: d

Explanation: We use certain algorithms to parse a context free grammar which include the most popular CYK algorithm which employs the concept of bottom up parsing and dynamic parsing.
\end{tcolorbox}
\end{mdframed}
\newpage

\begin{mdframed}[linewidth=1pt, roundcorner=5pt, linecolor=gray!30, backgroundcolor=gray!5]
\textbf{Question 7} 

The standard version of CYK algorithm operates only on context free grammars in the following form:

a) Greibach Normal form

b) Chomsky Normal form

c) Backus Naur form

d) All of the mentioned
\vspace{5pt} \par
\begin{tcolorbox}[colback=green!5, colframe=green!40!black, title=Solution]
Answer: b

Explanation: It requires the presence of a context free grammar into Chomsky Normal form to operate. However, every context free grammar can be converted into CNF for keeping the sense of grammar equivalent.
\end{tcolorbox}
\end{mdframed}
\newpage

\begin{mdframed}[linewidth=1pt, roundcorner=5pt, linecolor=gray!30, backgroundcolor=gray!5]
\textbf{Question 8} 

The \_\_\_\_\_\_\_\_\_\_ running time  of CYK is O(n
3
 .|G|)

where n is the length of the parse string and |G| is the size of the context free grammar G.

a) worst

b) best

c) average

d) none of the mentioned
\vspace{5pt} \par
\begin{tcolorbox}[colback=green!5, colframe=green!40!black, title=Solution]
Answer: a

Explanation: This is the worst case running time of CYK and and this makes it one of the most efficient algorithms for recognizing general context free languages in practice.
\end{tcolorbox}
\end{mdframed}
\newpage

\begin{mdframed}[linewidth=1pt, roundcorner=5pt, linecolor=gray!30, backgroundcolor=gray!5]
\textbf{Question 9} 

Which of the following is true for Valiants algorithm?

a) an extension of CYK

b) deals with efficient multiplication algorithms

c) matrices with 0-1 entries

d) all of the mentioned
\vspace{5pt} \par
\begin{tcolorbox}[colback=green!5, colframe=green!40!black, title=Solution]
Answer: d

Explanation: Valiants algorithm is actually an extention of CYK which even computes the same parsing table yet he showed another method can be utilized fro performing this operation.
\end{tcolorbox}
\end{mdframed}
\newpage

\begin{mdframed}[linewidth=1pt, roundcorner=5pt, linecolor=gray!30, backgroundcolor=gray!5]
\textbf{Question 10} 

Which among the following is a correct option in format for representing symbol and expression in Backus normal form?

a) $<$symbol$>$ -$>$expression

b) $<$symbol$>$::=\_expression\_

c) $<$symbol$>$=$<$expression$>$

d) all of the mentioned
\vspace{5pt} \par
\begin{tcolorbox}[colback=green!5, colframe=green!40!black, title=Solution]
Answer: b

Explanation: $<$symbol$>$::=\_expression\_ is the correct representation where $<$symbol$>$ is a non terminal, and expression consist of one or more sequence of symbols, more sequence are separated by |, indicating a choice.
\end{tcolorbox}
\end{mdframed}
\newpage

\newtopic{Intersection with Regular Languages}

\begin{mdframed}[linewidth=1pt, roundcorner=5pt, linecolor=gray!30, backgroundcolor=gray!5]
\textbf{Question 1} 

Which of the following is not a negative property of Context free languages?

a) Intersection

b) Complement

c) Intersection and Complement

d) None of the mentioned
\vspace{5pt} \par
\begin{tcolorbox}[colback=green!5, colframe=green!40!black, title=Solution]
Answer: c

Explanation: Context free languages are not closed under complement and intersection. Thus, are called Negative properties.
\end{tcolorbox}
\end{mdframed}
\newpage

\begin{mdframed}[linewidth=1pt, roundcorner=5pt, linecolor=gray!30, backgroundcolor=gray!5]
\textbf{Question 2} 

The intersection of context free language and regular language is \_\_\_\_\_\_\_\_\_

a) regular language

b) context free language

c) context sensitive language

d) non of the mentioned
\vspace{5pt} \par
\begin{tcolorbox}[colback=green!5, colframe=green!40!black, title=Solution]
Answer: b

Explanation: If a language L1 is regular and L2 is a context free language, then L1 intersection L2 will result into a context free language.
\end{tcolorbox}
\end{mdframed}
\newpage

\begin{mdframed}[linewidth=1pt, roundcorner=5pt, linecolor=gray!30, backgroundcolor=gray!5]
\textbf{Question 3} 

Which of the following is regular?

a) a
100
b
100

b) (a+b)*-\{a
100
b
100
\}

c) a
100
b
100
 and (a+b)*-\{a
100
b
100
\}

d) None of the mentioned
\vspace{5pt} \par
\begin{tcolorbox}[colback=green!5, colframe=green!40!black, title=Solution]
Answer: c

Explanation: As the language seems to be finite, a dfa can be constructed for the same, thus is regular.
\end{tcolorbox}
\end{mdframed}
\newpage

\begin{mdframed}[linewidth=1pt, roundcorner=5pt, linecolor=gray!30, backgroundcolor=gray!5]
\textbf{Question 4} 

Which of the following is not context free?

a) \{w: nA=nB=nC\}

b) \{a*b*c*\}

c) \{a
100
b
100
\}

d) All of the mentioned
\vspace{5pt} \par
\begin{tcolorbox}[colback=green!5, colframe=green!40!black, title=Solution]
Answer: d

Explanation: \{a*b*c*\} and (c) are regular languages while option (a) is not context free language.
\end{tcolorbox}
\end{mdframed}
\newpage

\begin{mdframed}[linewidth=1pt, roundcorner=5pt, linecolor=gray!30, backgroundcolor=gray!5]
\textbf{Question 5} 

Which of the following can be used to prove a language is not context free?

a) Ardens theorem

b) Power Construction method

c) Regular Closure

d) None of the mentioned
\vspace{5pt} \par
\begin{tcolorbox}[colback=green!5, colframe=green!40!black, title=Solution]
Answer: c

Explanation: We can use the properties of regular closure to prove that a language is not a context free language. Example: Intersection of context free language and regular language is a context free language. Proof by contradiction helps here.
\end{tcolorbox}
\end{mdframed}
\newpage

\begin{mdframed}[linewidth=1pt, roundcorner=5pt, linecolor=gray!30, backgroundcolor=gray!5]
\textbf{Question 6} 

Which of the following are valid membership algorithms?

a) CYK algorithm

b) Exhaustive search parser

c) CYK algorithm and Exhaustive search parser

d) None of the mentioned
\vspace{5pt} \par
\begin{tcolorbox}[colback=green!5, colframe=green!40!black, title=Solution]
Answer: c

Explanation: CYK algorithm is a parsing algorithm for context free grammars, which employs bottom up parsing and dynamic programming.
\end{tcolorbox}
\end{mdframed}
\newpage

\begin{mdframed}[linewidth=1pt, roundcorner=5pt, linecolor=gray!30, backgroundcolor=gray!5]
\textbf{Question 7} 

Which of the following belong to the steps to prove emptiness?

a) Remove useless variable

b) Check if a start variable S is useless

c) Remove useless variable and Check if a start variable S is useless

d) None of the mentioned
\vspace{5pt} \par
\begin{tcolorbox}[colback=green!5, colframe=green!40!black, title=Solution]
Answer: c

Explanation: The empty-language question can be stated as: For context free grammar G find if L(G) =f?
\end{tcolorbox}
\end{mdframed}
\newpage

\begin{mdframed}[linewidth=1pt, roundcorner=5pt, linecolor=gray!30, backgroundcolor=gray!5]
\textbf{Question 8} 

Which of the following is true for CYK Algorithm?

a) Triangular Table

b) Circular Chart

c) Linked List

d) None of the mentioned
\vspace{5pt} \par
\begin{tcolorbox}[colback=green!5, colframe=green!40!black, title=Solution]
Answer: a

Explanation: A triangular table is constructed to facilitate the solution of membership problem using bottom up parsing and dynamic programming.
\end{tcolorbox}
\end{mdframed}
\newpage

\begin{mdframed}[linewidth=1pt, roundcorner=5pt, linecolor=gray!30, backgroundcolor=gray!5]
\textbf{Question 9} 

Which of the following steps are wrong with respect to infiniteness problem?

a) Remove useless variables

b) Remove unit and epsilon production

c) Create dependency graph for variables

d) If there is a loop in the dependency graph the language is finite else infinite
\vspace{5pt} \par
\begin{tcolorbox}[colback=green!5, colframe=green!40!black, title=Solution]
Answer: d

Explanation: If we are able to detect a loop in the formed dependency graph, then the language in infinite.
\end{tcolorbox}
\end{mdframed}
\newpage

\begin{mdframed}[linewidth=1pt, roundcorner=5pt, linecolor=gray!30, backgroundcolor=gray!5]
\textbf{Question 10} 

State true or false:

Statement: Every context free language can be generated by a grammar which contains no useless non terminals.

a) true

b) false
\vspace{5pt} \par
\begin{tcolorbox}[colback=green!5, colframe=green!40!black, title=Solution]
Answer: a

Explanation: At first, we detect useless symbols and discard them. Inorder to find whether a symbol is useless, just make it the starting symbol and check for emptiness.
\end{tcolorbox}
\end{mdframed}
\newpage

\newtopic{Turing Machine – Notation and Transition Diagrams}

\begin{mdframed}[linewidth=1pt, roundcorner=5pt, linecolor=gray!30, backgroundcolor=gray!5]
\textbf{Question 1} 

A turing machine is a

a) real machine

b) abstract machine

c) hypothetical machine

d) more than one option is correct
\vspace{5pt} \par
\begin{tcolorbox}[colback=green!5, colframe=green!40!black, title=Solution]
Answer: d

Explanation: A turing machine is abstract or hypothetical machine thought by mathematician Alan Turing in 1936 capable of simulating any algorithm, however complicated it is.
\end{tcolorbox}
\end{mdframed}
\newpage

\begin{mdframed}[linewidth=1pt, roundcorner=5pt, linecolor=gray!30, backgroundcolor=gray!5]
\textbf{Question 2} 

A turing machine operates over:

a) finite memory tape

b) infinite memory tape

c) depends on the algorithm

d) none of the mentioned
\vspace{5pt} \par
\begin{tcolorbox}[colback=green!5, colframe=green!40!black, title=Solution]
Answer: b

Explanation: The turing machine operates on an infinite memory tape divided into cells. The machine positions its head over the cell and reads the symbol.
\end{tcolorbox}
\end{mdframed}
\newpage

\begin{mdframed}[linewidth=1pt, roundcorner=5pt, linecolor=gray!30, backgroundcolor=gray!5]
\textbf{Question 3} 

Which of the functions are not performed by the turing machine after reading a symbol?

a) writes the symbol

b) moves the tape one cell left/right

c) proceeds with next instruction or halts

d) none of the mentioned
\vspace{5pt} \par
\begin{tcolorbox}[colback=green!5, colframe=green!40!black, title=Solution]
Answer: d

Explanation: After the read head  reads the symbol from the input tape, it performs the following functions:

a) writes a symbol(some model allow symbol erasure/no writing)

b) moves the tape left or right (some models allows no motion)

c) proceeds with subsequent instruction or goes either into accepting halting state or rejecting halting state.
\end{tcolorbox}
\end{mdframed}
\newpage

\begin{mdframed}[linewidth=1pt, roundcorner=5pt, linecolor=gray!30, backgroundcolor=gray!5]
\textbf{Question 4} 

‘a' in a-machine is :

a) Alan

b) arbitrary

c) automatic

d) None of the mentioned
\vspace{5pt} \par
\begin{tcolorbox}[colback=green!5, colframe=green!40!black, title=Solution]
Answer: c

Explanation: The turing machine was invented by Alan turing in 1936. He named it as a-machine(automatic machine).
\end{tcolorbox}
\end{mdframed}
\newpage

\begin{mdframed}[linewidth=1pt, roundcorner=5pt, linecolor=gray!30, backgroundcolor=gray!5]
\textbf{Question 5} 

Which of the problems were not answered when the turing machine was invented?

a) Does a machine exists that can determine whether any arbitrary machine on its tape is circular.

b) Does a machine exists that can determine whether any arbitrary machine on its tape is ever prints a symbol

c) Hilbert Entscheidungs problem

d) None of the mentioned
\vspace{5pt} \par
\begin{tcolorbox}[colback=green!5, colframe=green!40!black, title=Solution]
Answer: d

Explanation: Invention of turing machine answered a lot of questions which included problems like decision problem, etc.) . Alan was able to prove the properties of computation using such model.
\end{tcolorbox}
\end{mdframed}
\newpage

\begin{mdframed}[linewidth=1pt, roundcorner=5pt, linecolor=gray!30, backgroundcolor=gray!5]
\textbf{Question 6} 

The ability for a system of instructions to simulate a Turing Machine is called \_\_\_\_\_\_\_\_\_

a) Turing Completeness

b) Simulation

c) Turing Halting

d) None of the mentioned
\vspace{5pt} \par
\begin{tcolorbox}[colback=green!5, colframe=green!40!black, title=Solution]
Answer: a

Explanation: Turing Completeness the ability for a system of instructions to simulate a Turing machine. A programming language that is Turing complete is theoretically capable of expressing all tasks accomplishable by computers; nearly all programming languages are Turing complete.
\end{tcolorbox}
\end{mdframed}
\newpage

\begin{mdframed}[linewidth=1pt, roundcorner=5pt, linecolor=gray!30, backgroundcolor=gray!5]
\textbf{Question 7} 

Turing machine can be represented using the following tools:

a) Transition graph

b) Transition table

c) Queue and Input tape

d) All of the mentioned
\vspace{5pt} \par
\begin{tcolorbox}[colback=green!5, colframe=green!40!black, title=Solution]
Answer: d

Explanation: We can represent a turing machine, graphically, tabularly and diagramatically.
\end{tcolorbox}
\end{mdframed}
\newpage

\begin{mdframed}[linewidth=1pt, roundcorner=5pt, linecolor=gray!30, backgroundcolor=gray!5]
\textbf{Question 8} 

Which of the following is false for an abstract machine?

a) Turing machine

b) theoretical model of computer

c) assumes a discrete time paradigm

d) all of the mentioned
\vspace{5pt} \par
\begin{tcolorbox}[colback=green!5, colframe=green!40!black, title=Solution]
Answer: d

Explanation: A n abstract machine also known as abstract computer, is a theoretical model of computer or hardware system in automata theory. Abstraction in computing process usually assumes a discrete time paradigm.
\end{tcolorbox}
\end{mdframed}
\newpage

\begin{mdframed}[linewidth=1pt, roundcorner=5pt, linecolor=gray!30, backgroundcolor=gray!5]
\textbf{Question 9} 

Fill in the blank with the most appropriate option.

Statement: In theory of computation, abstract machines are often used in \_\_\_\_\_\_\_\_\_\_\_ regarding computability or to analyze the complexity of an algorithm.

a) thought experiments

b) principle

c) hypothesis

d) all of the mentioned
\vspace{5pt} \par
\begin{tcolorbox}[colback=green!5, colframe=green!40!black, title=Solution]
Answer: d

Explanation: A thought experiment considers some hypothesis, theory or principle for the purpose of thinking through its consequences.
\end{tcolorbox}
\end{mdframed}
\newpage

\begin{mdframed}[linewidth=1pt, roundcorner=5pt, linecolor=gray!30, backgroundcolor=gray!5]
\textbf{Question 10} 

State true or false:

Statement: RAM model allows random access to indexed memory locations.

a) true

b) false
\vspace{5pt} \par
\begin{tcolorbox}[colback=green!5, colframe=green!40!black, title=Solution]
Answer: a

Explanation: In computer science, Random access machine is an abstract machine in the general class of register machines. Random access machine should not be confused with Random access memory.
\end{tcolorbox}
\end{mdframed}
\newpage

\newtopic{The Language of Turing Machine-1}

\begin{mdframed}[linewidth=1pt, roundcorner=5pt, linecolor=gray!30, backgroundcolor=gray!5]
\textbf{Question 1} 

A turing machine that is able to simulate other turing machines:

a) Nested Turing machines

b) Universal Turing machine

c) Counter machine

d) None of the mentioned
\vspace{5pt} \par
\begin{tcolorbox}[colback=green!5, colframe=green!40!black, title=Solution]
Answer: b

Explanation: A more mathematically oriented  definition with the same universal nature was introduced by church and turing together called the Church-Turing thesis(formal theory of computation).
\end{tcolorbox}
\end{mdframed}
\newpage

\begin{mdframed}[linewidth=1pt, roundcorner=5pt, linecolor=gray!30, backgroundcolor=gray!5]
\textbf{Question 2} 

Which of the problems are unsolvable?

a) Halting problem

b) Boolean Satisfiability problem

c) Halting problem \& Boolean Satisfiability problem

d) None of the mentioned
\vspace{5pt} \par
\begin{tcolorbox}[colback=green!5, colframe=green!40!black, title=Solution]
Answer: c

Explanation: Alan turing proved in 1936 that a general algorithm to solve the halting problem for all possible program-input pairs cannot exist.
\end{tcolorbox}
\end{mdframed}
\newpage

\begin{mdframed}[linewidth=1pt, roundcorner=5pt, linecolor=gray!30, backgroundcolor=gray!5]
\textbf{Question 3} 

Which of the following a turing machine does not consist of?

a) input tape

b) head

c) state register

d) none of the mentioned
\vspace{5pt} \par
\begin{tcolorbox}[colback=green!5, colframe=green!40!black, title=Solution]
Answer: d

Explanation: A state register is one which stores the state of the turing machine, one of the finitely many. Among these is the special start state with which the state register is initialized.
\end{tcolorbox}
\end{mdframed}
\newpage

\begin{mdframed}[linewidth=1pt, roundcorner=5pt, linecolor=gray!30, backgroundcolor=gray!5]
\textbf{Question 4} 

The value of n if turing machine is defined using n-tuples:

a) 6

b) 7

c) 8

d) 5
\vspace{5pt} \par
\begin{tcolorbox}[colback=green!5, colframe=green!40!black, title=Solution]
Answer: b

Explanation:

The 7-tuple definition of turing machine: (Q, S, G, d, q0, B, F)

where Q= The finite set of states of finite control

S= The finite set of input symbols

G= The complete set of tape symbols

d= The transition function

q0= The start state, a member of Q, in which the finite control is found initially.

B= The blank symbol

F= The set of final or accepting states, a subset of Q.
\end{tcolorbox}
\end{mdframed}
\newpage

\begin{mdframed}[linewidth=1pt, roundcorner=5pt, linecolor=gray!30, backgroundcolor=gray!5]
\textbf{Question 5} 

If  d is not defined on the current state and the current tape symbol, then the machine \_\_\_\_\_\_

a) does not halts

b) halts

c) goes into loop forever

d) none of the mentioned
\vspace{5pt} \par
\begin{tcolorbox}[colback=green!5, colframe=green!40!black, title=Solution]
Answer: b

Explanation: If we reach h
A
 or h
R
, we say TM halts. Once it has halted, it cannot move further, since  d is not defined at any pair (h
A
,X) or (h
R
,X) where h
A
 = accept halting state and h
R
 = reject halting state.
\end{tcolorbox}
\end{mdframed}
\newpage

\begin{mdframed}[linewidth=1pt, roundcorner=5pt, linecolor=gray!30, backgroundcolor=gray!5]
\textbf{Question 6} 

Statement: Instantaneous descriptions can be designed for a Turing machine.

State true or false:

a) true

b) false
\vspace{5pt} \par
\begin{tcolorbox}[colback=green!5, colframe=green!40!black, title=Solution]
Answer: a

Explanation: Inorder to describe formally what a Turing machine does, we need to develop a notation for configurations or Instantaneous descriptions(ID).
\end{tcolorbox}
\end{mdframed}
\newpage

\begin{mdframed}[linewidth=1pt, roundcorner=5pt, linecolor=gray!30, backgroundcolor=gray!5]
\textbf{Question 7} 

Which of the following are the models equivalent to Turing machine?

a) Multi tape turing machine

b) Multi track turing machine

c) Register machine

d) All of the mentioned
\vspace{5pt} \par
\begin{tcolorbox}[colback=green!5, colframe=green!40!black, title=Solution]
Answer: d

Explanation: Many machines that might be thought to have more computational capability than a simple UTM can be shown to have  no more power. They might compute faster or use less memory but cannot compute more powerfully i.e. more mathematical questions.
\end{tcolorbox}
\end{mdframed}
\newpage

\begin{mdframed}[linewidth=1pt, roundcorner=5pt, linecolor=gray!30, backgroundcolor=gray!5]
\textbf{Question 8} 

Which among the following is incorrect for o-machines?

a) Oracle Turing machines

b) Can be used to study decision problems

c) Visualizes Turing machine with a black box which is able to decide cerain decion problems in one operation

d) None of the mentioned
\vspace{5pt} \par
\begin{tcolorbox}[colback=green!5, colframe=green!40!black, title=Solution]
Answer: d

Explanation: In automata theory, an o- machine or oracle machine is a abstract machine used to study decision problems. The problem the oracle solves can be of any complexity class. Even undecidable problems like halting problems can be used.
\end{tcolorbox}
\end{mdframed}
\newpage

\begin{mdframed}[linewidth=1pt, roundcorner=5pt, linecolor=gray!30, backgroundcolor=gray!5]
\textbf{Question 9} 

RASP stands for:

a) Random access storage program

b) Random access stored program

c) Randomly accessed stored program

d) Random access storage programming
\vspace{5pt} \par
\begin{tcolorbox}[colback=green!5, colframe=green!40!black, title=Solution]
Answer: b

Explanation: RASP or Random access stored program is an abstract machine that has instances like modern stored computers.
\end{tcolorbox}
\end{mdframed}
\newpage

\begin{mdframed}[linewidth=1pt, roundcorner=5pt, linecolor=gray!30, backgroundcolor=gray!5]
\textbf{Question 10} 

Which of the following is not true about RASP?

a) Binary search can be performed more quickly using RASP than a turing machine

b) Stores its program in memory external to its state machines instructions

c) Has infinite number of distinguishable, unbounded registers

d) Binary search can be performed less quickly using RASP than a turing machine

e) More than two options are incorrect
\vspace{5pt} \par
\begin{tcolorbox}[colback=green!5, colframe=green!40!black, title=Solution]
Answer: d

Explanation: In theoretical computer science, the random access stored program( RASP ) machine model is an abstract machine used for the purpose of algorithm development and algorithm complexity theory.
\end{tcolorbox}
\end{mdframed}
\newpage

\begin{mdframed}[linewidth=1pt, roundcorner=5pt, linecolor=gray!30, backgroundcolor=gray!5]
\textbf{Question 11} 

State true or false:

Statement: RASP is to RAM like UTM is to turing machine.

a) true

b) false
\vspace{5pt} \par
\begin{tcolorbox}[colback=green!5, colframe=green!40!black, title=Solution]
Answer: a

Explanation: The Rasp is a random access machine model that, unlike the RAM has its program in its registers together with its input. The registers are unbounded(infinite in capacity); whether the number of registers is finite is model-specific.
\end{tcolorbox}
\end{mdframed}
\newpage

\newtopic{The Language of Turing Machine-2}

\begin{mdframed}[linewidth=1pt, roundcorner=5pt, linecolor=gray!30, backgroundcolor=gray!5]
\textbf{Question 1} 

The class of recursively enumerable language is known as:

a) Turing Class

b) Recursive Languages

c) Universal Languages

d) RE
\vspace{5pt} \par
\begin{tcolorbox}[colback=green!5, colframe=green!40!black, title=Solution]
Answer: d

Explanation: RE or recursively enumerable is only called the class of recursively enumerable language.
\end{tcolorbox}
\end{mdframed}
\newpage

\begin{mdframed}[linewidth=1pt, roundcorner=5pt, linecolor=gray!30, backgroundcolor=gray!5]
\textbf{Question 2} 

A language L is said to be Turing decidable if:

a) recursive

b) TM recognizes L

c) TM accepts L

d) recursive \& TM recognizes L
\vspace{5pt} \par
\begin{tcolorbox}[colback=green!5, colframe=green!40!black, title=Solution]
Answer: d

Explanation: A language L is recursively enumerable if there is a turing machine that accepts L, and recursive if there is a TM that recognizes L.(Sometimes these languages are alse called Turing-acceptable and Turing-decidable respectively).
\end{tcolorbox}
\end{mdframed}
\newpage

\begin{mdframed}[linewidth=1pt, roundcorner=5pt, linecolor=gray!30, backgroundcolor=gray!5]
\textbf{Question 3} 

Which of the following statements are false?

a) Every recursive language is recursively enumerable

b) Recursively enumerable language may not be recursive

c) Recursive languages may not be recursively enumerable

d) None of the mentioned
\vspace{5pt} \par
\begin{tcolorbox}[colback=green!5, colframe=green!40!black, title=Solution]
Answer: c

Explanation: Every recursive language is recursively enumerable but there exists recursively enumerable languages that are not recursive. If L is accepted by a Non deterministic TM  T, and every possible sequence of moves of T causes it to halt, then L is recursive.
\end{tcolorbox}
\end{mdframed}
\newpage

\begin{mdframed}[linewidth=1pt, roundcorner=5pt, linecolor=gray!30, backgroundcolor=gray!5]
\textbf{Question 4} 

Choose the correct option:

Statement: If L1 and L2 are recursively enumerable languages over S, then the following is/are recursively enumerable.

a) L1 U L2

b) L2 $\cap$ L2

c) Both L1 U L2 and L2 $\cap$ L2

d) None of the mentioned
\vspace{5pt} \par
\begin{tcolorbox}[colback=green!5, colframe=green!40!black, title=Solution]
Answer: c

Explanation: Both the union and intersection operations preserve the property of recursive enumerablity(Theorem).
\end{tcolorbox}
\end{mdframed}
\newpage

\begin{mdframed}[linewidth=1pt, roundcorner=5pt, linecolor=gray!30, backgroundcolor=gray!5]
\textbf{Question 5} 

If L is a recursive language, L' is:

a) Recursive

b) Recursively Enumerable

c) Recursive and Recursively Enumerable

d) None of the mentioned
\vspace{5pt} \par
\begin{tcolorbox}[colback=green!5, colframe=green!40!black, title=Solution]
Answer: c

Explanation: If T is a turing machine recognizing L, we can make it recognize L' by interchanging the two outputs. And every recursive language is recursively enumerable.
\end{tcolorbox}
\end{mdframed}
\newpage

\begin{mdframed}[linewidth=1pt, roundcorner=5pt, linecolor=gray!30, backgroundcolor=gray!5]
\textbf{Question 6} 

Choose the appropriate option:

Statement: If a language L is recursive, it is closed under the following operations:

a) Union

b) Intersection

c) Complement

d) All of the mentioned
\vspace{5pt} \par
\begin{tcolorbox}[colback=green!5, colframe=green!40!black, title=Solution]
Answer: d

Explanation: The closure property of recursive languages include union, intersection and complement operations.
\end{tcolorbox}
\end{mdframed}
\newpage

\begin{mdframed}[linewidth=1pt, roundcorner=5pt, linecolor=gray!30, backgroundcolor=gray!5]
\textbf{Question 7} 

A recursively enumerable language L can be recursive if:

a) L' is recursively enumerable

b) Every possible sequence of moves of T, the TM which accept L, causes it to halt

c) L' is recursively enumerable and every possible sequence of moves of T, the TM which accept L, causes it to halt

d) None of the mentioned
\vspace{5pt} \par
\begin{tcolorbox}[colback=green!5, colframe=green!40!black, title=Solution]
Answer: c

Explanation: Theorem- If L is a recursively enumerable language whose complement is recursively enumerable, then L is recursive.
\end{tcolorbox}
\end{mdframed}
\newpage

\begin{mdframed}[linewidth=1pt, roundcorner=5pt, linecolor=gray!30, backgroundcolor=gray!5]
\textbf{Question 8} 

A language L is recursively enumerable if L=L(M) for some turing machine M.


 
\begin{center}\includegraphics[width=0.5\textwidth]{automata-theory-questions-answers-freshers-q8.png}\end{center}
 

Which among the following cannot be among A, B and C?

a) yes  w $\in$ L

b) no w $\notin$ L

c) M does not halt w $\notin$ L

d) None of the mentioned
\vspace{5pt} \par
\begin{tcolorbox}[colback=green!5, colframe=green!40!black, title=Solution]
Answer: d

Explanation:


 
\begin{center}\includegraphics[width=0.5\textwidth]{automata-theory-questions-answers-freshers-q8a.png}\end{center}
\end{tcolorbox}
\end{mdframed}
\newpage

\begin{mdframed}[linewidth=1pt, roundcorner=5pt, linecolor=gray!30, backgroundcolor=gray!5]
\textbf{Question 9} 

State true or false:

Statement: An enumerator is a turing machine with extra output tape T, where symbols, once written, are never changed.

a) true

b) false
\vspace{5pt} \par
\begin{tcolorbox}[colback=green!5, colframe=green!40!black, title=Solution]
Answer: a

Explanation: To enumerate a set means to list the elements once at a time, and to say that a set is enumerable should perhaps mean that there exists an algorithm for enumerating it.
\end{tcolorbox}
\end{mdframed}
\newpage

\begin{mdframed}[linewidth=1pt, roundcorner=5pt, linecolor=gray!30, backgroundcolor=gray!5]
\textbf{Question 10} 

A Language L may not be accepted by a Turing Machine if:

a) It is recursively enumerable

b) It is recursive

c) L can be enumerated by some turing machine

d) None of the mentioned
\vspace{5pt} \par
\begin{tcolorbox}[colback=green!5, colframe=green!40!black, title=Solution]
Answer: b

Explanation: A language L is recursively enumerable if and only if it can be enumerated by some turing machine. A recursive enumerable language may or may not be recursive.
\end{tcolorbox}
\end{mdframed}
\newpage

\newtopic{Turing Machine and Halting}

\begin{mdframed}[linewidth=1pt, roundcorner=5pt, linecolor=gray!30, backgroundcolor=gray!5]
\textbf{Question 1} 

Which of the following regular expression resembles the given diagram?


 
\begin{center}\includegraphics[width=0.5\textwidth]{automata-theory-questions-answers-turing-machine-halting-q1.png}\end{center}
 

a) \{a\}*\{b\}*\{a,b\}

b) \{a,b\}*\{aba\}

c) \{a,b\}*\{bab\}

d) \{a,b\}*\{a\}*\{b\}*
\vspace{5pt} \par
\begin{tcolorbox}[colback=green!5, colframe=green!40!black, title=Solution]
Answer: b

Explanation: The given diagram is a transition graph for a turing machine which accepts the language with the regular expression \{a,b\}*\{aba\}.
\end{tcolorbox}
\end{mdframed}
\newpage

\begin{mdframed}[linewidth=1pt, roundcorner=5pt, linecolor=gray!30, backgroundcolor=gray!5]
\textbf{Question 2} 

Construct a turing machine which accepts a string with ‘aba' as its substring.

a)
 
\begin{center}\includegraphics[width=0.5\textwidth]{automata-theory-questions-answers-turing-machine-halting-q2a.png}\end{center}
 

b)
 
\begin{center}\includegraphics[width=0.5\textwidth]{automata-theory-questions-answers-turing-machine-halting-q2b.png}\end{center}
 

c)
 
\begin{center}\includegraphics[width=0.5\textwidth]{automata-theory-questions-answers-turing-machine-halting-q2c.png}\end{center}
 

d)
 
\begin{center}\includegraphics[width=0.5\textwidth]{automata-theory-questions-answers-turing-machine-halting-q2d.png}\end{center}
\vspace{5pt} \par
\begin{tcolorbox}[colback=green!5, colframe=green!40!black, title=Solution]
Answer: c

Explanation: The language consist of strings with a substring ‘aba' as fixed at its end and the left part can be anything including epsilon. Thus the turing machine uses five states to express the language excluding the rejection halting state which if allowed can modify the graph as:


 
\begin{center}\includegraphics[width=0.5\textwidth]{automata-theory-questions-answers-turing-machine-halting-q2e.png}\end{center}
\end{tcolorbox}
\end{mdframed}
\newpage

\begin{mdframed}[linewidth=1pt, roundcorner=5pt, linecolor=gray!30, backgroundcolor=gray!5]
\textbf{Question 3} 

The number of states required to automate the last question i.e. \{a,b\}*\{aba\}\{a,b\}* using finite automata:

a) 4

b) 3

c) 5

d) 6
\vspace{5pt} \par
\begin{tcolorbox}[colback=green!5, colframe=green!40!black, title=Solution]
Answer: a

Explanation: The finite automata can be represented as:


 
\begin{center}\includegraphics[width=0.5\textwidth]{automata-theory-questions-answers-turing-machine-halting-q3.png}\end{center}
\end{tcolorbox}
\end{mdframed}
\newpage

\begin{mdframed}[linewidth=1pt, roundcorner=5pt, linecolor=gray!30, backgroundcolor=gray!5]
\textbf{Question 4} 

The machine accept the string by entering into hA or it can:

a) explicitly reject x by entering into hR

b) enter into an infinte loop

c) explicitly reject x by entering into hR and enter into an infinte loop

d) None of the mentioned
\vspace{5pt} \par
\begin{tcolorbox}[colback=green!5, colframe=green!40!black, title=Solution]
Answer: c

Explanation: Three things can occur when a string is tested over a turing machine:

a) enter into accept halting state

b) enter into reject halting state

c) goes into loop forever
\end{tcolorbox}
\end{mdframed}
\newpage

\begin{mdframed}[linewidth=1pt, roundcorner=5pt, linecolor=gray!30, backgroundcolor=gray!5]
\textbf{Question 5} 

d(q,X)=(r,Y,D) where D cannot be:


 
\begin{center}\includegraphics[width=0.5\textwidth]{automata-theory-questions-answers-turing-machine-halting-q5.png}\end{center}
 

a) L

b) R

c) S

d) None of the mentioned
\vspace{5pt} \par
\begin{tcolorbox}[colback=green!5, colframe=green!40!black, title=Solution]
Answer: c

Explanation: D represents the direction in which automata moves forward as per the queue which surely cannot be a starting variable.
\end{tcolorbox}
\end{mdframed}
\newpage

\begin{mdframed}[linewidth=1pt, roundcorner=5pt, linecolor=gray!30, backgroundcolor=gray!5]
\textbf{Question 6} 

Which of the following can accept even palindrome over \{a,b\}

a) Push down Automata

b) Turing machine

c) NDFA

d) All of the mentioned
\vspace{5pt} \par
\begin{tcolorbox}[colback=green!5, colframe=green!40!black, title=Solution]
Answer: c

Explanation: A language generating strings which are palindrome is not regular, thus cannot b represented using a finite automaton.
\end{tcolorbox}
\end{mdframed}
\newpage

\begin{mdframed}[linewidth=1pt, roundcorner=5pt, linecolor=gray!30, backgroundcolor=gray!5]
\textbf{Question 7} 

Which of the functions can a turing machine not perform?

a) Copying a string

b) Deleting a symbol

c) Accepting a pal

d) Inserting a symbol
\vspace{5pt} \par
\begin{tcolorbox}[colback=green!5, colframe=green!40!black, title=Solution]
Answer: d

Explanation: Different turing machines exist for operations like copying a string, deleting a symbol, inserting a symbol and accepting palindromes.
\end{tcolorbox}
\end{mdframed}
\newpage

\begin{mdframed}[linewidth=1pt, roundcorner=5pt, linecolor=gray!30, backgroundcolor=gray!5]
\textbf{Question 8} 

If T1 and T2 are two turing machines. The composite can be represented using the expression:

a) T1T2

b) T1 U T2

c) T1 X T2

d) None of the mentioned
\vspace{5pt} \par
\begin{tcolorbox}[colback=green!5, colframe=green!40!black, title=Solution]
Answer: a

Explanation: If T1 and T2 are TMs, with disjoint sets of non halting states and transition function d1 and d2, respectively, we write T1T2  to denote this composite TM.
\end{tcolorbox}
\end{mdframed}
\newpage

\begin{mdframed}[linewidth=1pt, roundcorner=5pt, linecolor=gray!30, backgroundcolor=gray!5]
\textbf{Question 9} 

The following turing machine acts like:


 
\begin{center}\includegraphics[width=0.5\textwidth]{automata-theory-questions-answers-turing-machine-halting-q9.png}\end{center}
 

a) Copies a string

b) Delete a symbol

c) Insert a symbol

d) None of the mentioned
\vspace{5pt} \par
\begin{tcolorbox}[colback=green!5, colframe=green!40!black, title=Solution]
Answer: b

Explanation: A turing machine does the deletion by changing the tape contents from yaz to yz, where y belongs to (S U \{\#\})*.
\end{tcolorbox}
\end{mdframed}
\newpage

\begin{mdframed}[linewidth=1pt, roundcorner=5pt, linecolor=gray!30, backgroundcolor=gray!5]
\textbf{Question 10} 

What does the following transition graph shows:


 
\begin{center}\includegraphics[width=0.5\textwidth]{automata-theory-questions-answers-turing-machine-halting-q10.png}\end{center}
 

a) Copies a symbol

b) Reverses a string

c) Accepts a pal

d) None of the mentioned
\vspace{5pt} \par
\begin{tcolorbox}[colback=green!5, colframe=green!40!black, title=Solution]
Answer: c

Explanation: The composite TM accepts the language of palindromes over \{a, b\} by comparing the input string to its reverse and accepting if and only if the two are equal.
\end{tcolorbox}
\end{mdframed}
\newpage

\newtopic{Programming Techniques-Storage and Subroutines}

\begin{mdframed}[linewidth=1pt, roundcorner=5pt, linecolor=gray!30, backgroundcolor=gray!5]
\textbf{Question 1} 

A turing machine has \_\_\_\_\_\_\_\_\_\_\_\_ number of states in a CPU.

a) finite

b) infinte

c) May be finite

d) None of the mentioned
\vspace{5pt} \par
\begin{tcolorbox}[colback=green!5, colframe=green!40!black, title=Solution]
Answer: a

Explanation: A turing machine has finite number of states in its CPU. However, the states are not small in number. Real computer consist of registers which can store values (fixed number of bits).
\end{tcolorbox}
\end{mdframed}
\newpage

\begin{mdframed}[linewidth=1pt, roundcorner=5pt, linecolor=gray!30, backgroundcolor=gray!5]
\textbf{Question 2} 

Suppose we have a simple computer with control unit holding a PC with a 32 bit address + Arithmetic unit holding one double length 64 bit Arithmetic Register. The number of states the finite machine will hold:

a) 2
(32*64)

b) 2
96

c) 96

d) 32
\vspace{5pt} \par
\begin{tcolorbox}[colback=green!5, colframe=green!40!black, title=Solution]
Answer: b

Explanation: According to the statistics of the question, we will have a finite machine with 2\textasciicircum{}96 states.
\end{tcolorbox}
\end{mdframed}
\newpage

\begin{mdframed}[linewidth=1pt, roundcorner=5pt, linecolor=gray!30, backgroundcolor=gray!5]
\textbf{Question 3} 

In one move a turing machine will:


 
\begin{center}\includegraphics[width=0.5\textwidth]{automata-theory-questions-answers-programming-techniques-storage-subroutines-q3.png}\end{center}
 

a) Change a state

b) Write a tape symbol in the cell scanned

c) Move the tape head left or right

d) All of the mentioned
\vspace{5pt} \par
\begin{tcolorbox}[colback=green!5, colframe=green!40!black, title=Solution]
Answer: d

Explanation: A move of a turing machine is the function of the state of finite control and the tape symbol just scanned.
\end{tcolorbox}
\end{mdframed}
\newpage

\begin{mdframed}[linewidth=1pt, roundcorner=5pt, linecolor=gray!30, backgroundcolor=gray!5]
\textbf{Question 4} 

State true or false:

Statement: We can use the finite control of turing machine to hold a finite amount of data.

a) true

b) false
\vspace{5pt} \par
\begin{tcolorbox}[colback=green!5, colframe=green!40!black, title=Solution]
Answer: a

Explanation:


 
\begin{center}\includegraphics[width=0.5\textwidth]{automata-theory-questions-answers-programming-techniques-storage-subroutines-q3.png}\end{center}
 

The finite control not only contains state q but also three data, A, B, C. The following technique requires no extension to the Turing Machine model. Shaping states this way allows to describe transitions in more systematic way and often to simplify the strategy of the program.
\end{tcolorbox}
\end{mdframed}
\newpage

\begin{mdframed}[linewidth=1pt, roundcorner=5pt, linecolor=gray!30, backgroundcolor=gray!5]
\textbf{Question 5} 

Statement 1: Multitrack Turing machine.

Statement 2: Gamma is Cartesian product of a finite number of finite sets.

Which among the following is the correct option?

a) Statement 1 is the assertion and Statement 2 is the reason

b) Statement 1 is the reason and Statement 2 is the assertion

c) Statement 1 and Statement 2 are independent from each other

d) None of the mentioned
\vspace{5pt} \par
\begin{tcolorbox}[colback=green!5, colframe=green!40!black, title=Solution]
Answer: a

Explanation: Cartesian product works like a struct in C/C++. For Example: Computer tape storage is something like 8 or 9 bits in each cell. One can recognize a multi track tape machine by looking at the transitions because each will have tuples as the read and write symbols.
\end{tcolorbox}
\end{mdframed}
\newpage

\begin{mdframed}[linewidth=1pt, roundcorner=5pt, linecolor=gray!30, backgroundcolor=gray!5]
\textbf{Question 6} 

A multi track turing machine can described as a 6-tuple (Q, X, S, d, q0, F) where X represents:

a) input alphabet

b) tape alphabet

c) shift symbols

d) none of the mentioned
\vspace{5pt} \par
\begin{tcolorbox}[colback=green!5, colframe=green!40!black, title=Solution]
Answer: b

Explanation: The 6-tuple (Q, X, S, d, q0, F)  can be explained as:

Q represents finite set of states,

X represents the tape alphabet,

S represents the input alphabet

d represents the relation on states and the symbols

q0 represents the initial state

F represents the set of final states.
\end{tcolorbox}
\end{mdframed}
\newpage

\begin{mdframed}[linewidth=1pt, roundcorner=5pt, linecolor=gray!30, backgroundcolor=gray!5]
\textbf{Question 7} 

Which of the following statements are false?

a) A multi track turing machine is a special kind of multi tape turing machine

b) 4-heads move independently along 4-tracks in standard 4-tape turing machine

c) In a n-track turing machine, n head reads and writes on all the tracks simultaneously.

d) All of the mentioned
\vspace{5pt} \par
\begin{tcolorbox}[colback=green!5, colframe=green!40!black, title=Solution]
Answer: c

Explanation: In a n-track turing machine, one head reads and writes on all the tracks simultaneously.
\end{tcolorbox}
\end{mdframed}
\newpage

\begin{mdframed}[linewidth=1pt, roundcorner=5pt, linecolor=gray!30, backgroundcolor=gray!5]
\textbf{Question 8} 

State true or false:

Statement: Two track turing machine is equivalent to a standard turing machine.

a) true

b) false
\vspace{5pt} \par
\begin{tcolorbox}[colback=green!5, colframe=green!40!black, title=Solution]
Answer: a

Explanation: This can be generalized for n- tracks and can be proved equivalent using ennumerable languages.
\end{tcolorbox}
\end{mdframed}
\newpage

\begin{mdframed}[linewidth=1pt, roundcorner=5pt, linecolor=gray!30, backgroundcolor=gray!5]
\textbf{Question 9} 

Which of the following is/are not true for recursively enumerable language?

a) partially decidable

b) Turing acceptable

c) Turing Recognizable

d) None of the mentioned
\vspace{5pt} \par
\begin{tcolorbox}[colback=green!5, colframe=green!40!black, title=Solution]
Answer: d

Explanation: In automata theory, a formal language is called recursively enumerable language or partially decidable or semi decidable or turing acceptable or turing recognizable if there exists a turing machine which will enumerate all valid strings of the language.
\end{tcolorbox}
\end{mdframed}
\newpage

\begin{mdframed}[linewidth=1pt, roundcorner=5pt, linecolor=gray!30, backgroundcolor=gray!5]
\textbf{Question 10} 

According to Chomsky hierarchy, which of the following is adopted by Recursively Ennumerable language?

a) Type 0

b) Type 1

c) Type 2

d) Type 3
\vspace{5pt} \par
\begin{tcolorbox}[colback=green!5, colframe=green!40!black, title=Solution]
Answer: a

Explanation: Recursively Ennumerable languages are type 0 languages in the Chomsky hierarchy. All regular, context free, context sensitive languages are recursively enumerable language.
\end{tcolorbox}
\end{mdframed}
\newpage

\newtopic{Multitape Turing Machines}

\begin{mdframed}[linewidth=1pt, roundcorner=5pt, linecolor=gray!30, backgroundcolor=gray!5]
\textbf{Question 1} 

A turing machine with several tapes in known as:

a) Multi-tape turing machine

b) Poly-tape turing maching

c) Universal turing machine

d) All of the mentioned
\vspace{5pt} \par
\begin{tcolorbox}[colback=green!5, colframe=green!40!black, title=Solution]
Answer: a

Explanation: A multitape turing machine is an ordinary turing machine with multiple tapes. Each tape has its own head to control the read and write.
\end{tcolorbox}
\end{mdframed}
\newpage

\begin{mdframed}[linewidth=1pt, roundcorner=5pt, linecolor=gray!30, backgroundcolor=gray!5]
\textbf{Question 2} 

A multitape turing machine is \_\_\_\_\_\_\_\_ powerful than a single tape turing machine.

a) more

b) less

c) equal

d) none of the mentioned
\vspace{5pt} \par
\begin{tcolorbox}[colback=green!5, colframe=green!40!black, title=Solution]
Answer: a

Explanation: The multitape turing machine model seems much powerful than the single tape model, but any multi tape machine, no matter how many tapes, can be simulated by single taped TM.
\end{tcolorbox}
\end{mdframed}
\newpage

\begin{mdframed}[linewidth=1pt, roundcorner=5pt, linecolor=gray!30, backgroundcolor=gray!5]
\textbf{Question 3} 

In what ratio, more computation time is needed to simulate multitape turing machines using single tape turing machines?

a) doubly

b) triple

c) quadratically

d) none of the mentioned
\vspace{5pt} \par
\begin{tcolorbox}[colback=green!5, colframe=green!40!black, title=Solution]
Answer: c

Explanation: Thus, multitape turing machines cannot calculate any more functions than single tape machines.
\end{tcolorbox}
\end{mdframed}
\newpage

\begin{mdframed}[linewidth=1pt, roundcorner=5pt, linecolor=gray!30, backgroundcolor=gray!5]
\textbf{Question 4} 

Which of the following is true for two stack turing machines?

a) one read only input

b) two storage tapes

c) one read only input \& two storage tapes

d) None of the mentioned
\vspace{5pt} \par
\begin{tcolorbox}[colback=green!5, colframe=green!40!black, title=Solution]
Answer: c

Explanation: Two-stack Turing machines have a read-only input and two storage tapes. If a head moves left on either tape a blank is printed on that tape, but one symbol from a "library" can be printed.
\end{tcolorbox}
\end{mdframed}
\newpage

\begin{mdframed}[linewidth=1pt, roundcorner=5pt, linecolor=gray!30, backgroundcolor=gray!5]
\textbf{Question 5} 

Which of the following is not a Non deterministic turing machine?

a) Alternating Turing machine

b) Probabalistic Turing machine

c) Read-only turing machine

d) None of the mentioned
\vspace{5pt} \par
\begin{tcolorbox}[colback=green!5, colframe=green!40!black, title=Solution]
Answer: c

Explanation: A read only turing machine or 2 way deterministic finite automaton is a class of model of computability that behaves like a turing machine, and can move in both directions across input, except cannot write to its input tape.
\end{tcolorbox}
\end{mdframed}
\newpage

\begin{mdframed}[linewidth=1pt, roundcorner=5pt, linecolor=gray!30, backgroundcolor=gray!5]
\textbf{Question 6} 

Which of the turing machines have existential and universal states?

a) Alternating Turing machine

b) Probalistic Turing machine

c) Read-only turing machine

d) None of the mentioned
\vspace{5pt} \par
\begin{tcolorbox}[colback=green!5, colframe=green!40!black, title=Solution]
Answer: a

Explanation: ATM is divide into two sets: an existential state is accepting if some transitions leads to an accepting state; an universal state is accepting if every transition leads to an accepting state.
\end{tcolorbox}
\end{mdframed}
\newpage

\begin{mdframed}[linewidth=1pt, roundcorner=5pt, linecolor=gray!30, backgroundcolor=gray!5]
\textbf{Question 7} 

Which of the following is false for Quantum Turing machine?

a) Abstract machine

b) Any quantum algorithm can be expressed formally as a particular quantum turing machine

c) Gives a solution to ‘Is a universal quantum computer sufficient'

d) None of the mentioned
\vspace{5pt} \par
\begin{tcolorbox}[colback=green!5, colframe=green!40!black, title=Solution]
Answer: c

Explanation: ‘Is a universal quantum computer sufficient' is one of the unsolved problem from physics.
\end{tcolorbox}
\end{mdframed}
\newpage

\begin{mdframed}[linewidth=1pt, roundcorner=5pt, linecolor=gray!30, backgroundcolor=gray!5]
\textbf{Question 8} 

A deterministic turing machine is:

a) ambiguous turing machine

b) unambiguous turing machine

c) non-deterministic

d) none of the mentioned
\vspace{5pt} \par
\begin{tcolorbox}[colback=green!5, colframe=green!40!black, title=Solution]
Answer: b

Explanation: A deterministic turing machine is unambiguous and for every input, there is exactly one operation possible. It is a subset of non-deterministic Turing machines.
\end{tcolorbox}
\end{mdframed}
\newpage

\begin{mdframed}[linewidth=1pt, roundcorner=5pt, linecolor=gray!30, backgroundcolor=gray!5]
\textbf{Question 9} 

Which of the following is true about Turing's a-machine?

a) a stands for automatic

b) left ended, right end-infinite

c) finite number of tape symbols were allowed

d) all of the mentioned
\vspace{5pt} \par
\begin{tcolorbox}[colback=green!5, colframe=green!40!black, title=Solution]
Answer: d

Explanation: Turings a- machine or automatic machine was left ended,right end infinite.Any of finite number of tape symbols were allowed and the 5 tuples were not in order.
\end{tcolorbox}
\end{mdframed}
\newpage

\begin{mdframed}[linewidth=1pt, roundcorner=5pt, linecolor=gray!30, backgroundcolor=gray!5]
\textbf{Question 10} 

Which of the following is a multi tape turing machine?

a) Post turing Machine

b) Wang-B Machine

c) Oblivious turing Machine

d) All of the mentioned
\vspace{5pt} \par
\begin{tcolorbox}[colback=green!5, colframe=green!40!black, title=Solution]
Answer: c

Explanation: An oblivious turing machine where movements of various heads are fixed functions of time, independent of the input. Pippenger and Fischer showed that any computation that can be performed by a multi-tape Turing machine in n steps can be performed by an oblivious two-tape Turing machine in O(n log n) steps.
\end{tcolorbox}
\end{mdframed}
\newpage

\end{document}